% =====================================================================
%  R. Nielsen, THE SPECULATIVE LOGIC IN ITS MAIN OUTLINES.  Fascicles 1--3.  Copenhagen, 1841--42.  192 pp.
%  (Den speculative Logik i dens Grundtræk; Fascicle 4 is not in the scanned volume, and the text breaks off on p. 192.)
%  ENGLISH TRANSLATION, from transcription.tex (Danish).
%  GENERATED by ebook-method/speculative-logik/tr_build.py --write from this template and the five translators' parts in
%  texts/nielsen/speculative-logik/.parts/tr/ (TRANSLATOR-BRIEF.md; terms: GLOSSARY.md, approved 2026-10-07).
%  Title: "Grundtræk" = main outlines, fundamental features; the author's note calls the book "these main outlines".
%  Conventions: TRANSLATION-PLAYBOOK.md.  Page numbers of the 1841--42 print in the margin.  Emphasis as in the Danish:
%  \textbf = the print's bold Fraktur, \emph = its letter-spacing.  Footnote marks *) **) restart on every printed page.
% =====================================================================
\documentclass[12pt,a4paper,openany]{book}
\usepackage[utf8]{inputenc}
\usepackage[T1]{fontenc}
\usepackage{libertinus}
\usepackage{libertinust1math}
\usepackage{textalpha}
\usepackage{graphicx}
\DeclareUnicodeCharacter{0254}{\reflectbox{c}}  % ɔ (U+0254) = reversed c
\DeclareUnicodeCharacter{221E}{\ensuremath{\infty}}
\DeclareUnicodeCharacter{2032}{\ensuremath{'}}
\DeclareUnicodeCharacter{2033}{\ensuremath{''}}
\DeclareUnicodeCharacter{2034}{\ensuremath{'''}}
\DeclareUnicodeCharacter{2057}{\ensuremath{''''}}
\DeclareUnicodeCharacter{2212}{\ensuremath{-}}
\DeclareUnicodeCharacter{2299}{\ensuremath{\odot}}
\DeclareUnicodeCharacter{A75B}{r}   % r rotunda
\usepackage[english]{babel}
\usepackage{enumitem}
\usepackage[protrusion=true,expansion=true]{microtype}
\usepackage{setspace}
\usepackage[margin=1.3in]{geometry}
\usepackage{fancyhdr}
\usepackage{hyperref}
\hypersetup{pdftitle={The Speculative Logic in Its Main Outlines}, pdfauthor={Rasmus Nielsen}, hidelinks}
\setlength{\headheight}{15pt}
\pagestyle{fancy}
\fancyhf{}
\fancyhead[LE,RO]{\thepage}
\fancyhead[C]{\textit{R. Nielsen: The Speculative Logic in Its Main Outlines (1841--42)}}
\renewcommand{\thefootnote}{\ifcase\value{footnote}\or *)\or **)\or ***)\else
  \arabic{footnote})\fi}
\setstretch{1.3}
\usepackage{marginnote}
\newcommand{\opage}[1]{\marginnote{\footnotesize\textit{[#1]}}}
% headings as in the Danish edition: Chapter (bold, with its title), Section A./B./C. (bold), a)/b)/c) (spaced),
% the § number and its argument (bold), and the centred "Transition." (spaced)
\newcommand{\capitel}[2]{\chapter*{\centering\textbf{#1}\\[0.3em]\textbf{#2}}\addcontentsline{toc}{chapter}{#1 #2}}
\newcommand{\afsnit}[1]{\section*{\centering\textbf{#1}}\addcontentsline{toc}{section}{#1}}
\newcommand{\underafsnit}[1]{\subsection*{\centering\emph{#1}}\addcontentsline{toc}{subsection}{#1}}
\newcommand{\paragraf}[2]{\begin{center}§ #1.\\[0.2em]\textbf{#2}\end{center}}
\newcommand{\overgang}{\begin{center}\emph{Transition.}\end{center}}

\begin{document}

% PDF 6, the title of Fascicle 1 (in a Greek-key border in the print, not reproduced)
\begin{titlepage}
\begin{center}
\vspace*{4em}
{\Huge\textbf{The Speculative Logic}}\\[1em]
{\Large in Its Main Outlines}\\[1em]
by\\[1em]
{\Large\textbf{Lic. Rasmus Nielsen,}}\\[0.5em]
Professor Extraordinary of Philosophy at the University of Copenhagen.\\[2em]
\rule{6em}{0.4pt}\\[2em]
{\large\textbf{First Fascicle.}}\\[1em]
\textit{1841.}\\[4em]
{\small Translated from the Danish}
\end{center}
\end{titlepage}
% PDF 7, the verso: the author's note
\thispagestyle{empty}
\vspace*{8em}
These main outlines are to be regarded as a fragment of a philosophical methodology, the first part of which was to
contain the Logic with a prefatory Introduction. The necessity of having a printed guide for the oral lectures has
hastened the publication. The remaining fascicles will follow as the course of lectures announced in the Lecture
Catalogue advances.

\medskip
Copenhagen, 10 November 1841.\hfill\textbf{R. Nielsen.}
\clearpage

% ===== TEXT (pp. 1--192) =====
% --- p. 1 ---
\setcounter{footnote}{0}%
\opage{1}\capitel{Chapter One.}{Being.}
\afsnit{A. Quality.}
\underafsnit{a) Pure Being.}
\paragraf{11}{Being, Nothing, Becoming.}

\emph{Being} is the most comprehensive and precisely for that reason the
most contentless of all positive determinations; as such it has
nothing for its limit. In virtue of this delimitation, being and nothing must exclude each other,
but precisely for that reason also determine each other,
i.e. be posited as each other's predicates. This double
contradiction, with which each of the terms is burdened,
is sublated in becoming, which contains the infinite
alternation of both.
The copula ``is'', which, hidden or overt, expresses the connection between every subject and its predicates, indicates, when we consider it more closely, that all determinately existent things, all concrete determinations, the essential as well as the unessential, have their common medium in being. Being, then, taken immediately, relates to everything as space relates to sensible objects. If we think all objects away, then there remains only
% --- p. 2 ---
\setcounter{footnote}{0}%
\opage{2}the empty place, pure space; so too,
when we abstract from all the special determinations of determinate existence,
we get pure being back.
Insofar, then, as things coincide with
being, they thereby coincide with one another; insofar as they, on the other hand, differ from one another,
they fall outside pure being. Pure being,
therefore, one must not think of as an image of fantasy;
to seek it one need not stare into
the blue air or down into the abyss of thoughts; it is
not to be regarded as a conjured-up shape that only
appears at a mystical midnight, but as a wholly simple
moment of determinate existence (momentum simplicissimum),
which is everywhere equally near to us. Pure being runs
through tree and stone as well as through consciousness
and its thoughts; but the main thing is that the eye not be dazzled by the countless distinction-determinations of heterogeneous objects.

What we must hold fast here, then, is only this, that,
however different things may otherwise be, it nevertheless
holds of them all that they \textbf{are}.

Since there are now no distinction-determinations in
pure being, and since the diversities, if as such they
fall outside being, must be posited as non-being:
all distinction must vanish. This vanishing
has its positive and its negative form. From the positive
side the relation presents itself thus, that being itself
spreads out over all distinction-moments, assimilates
them all to itself, since it is itself the primal element
of which they are all formed; from the negative side, on the other hand,
thus, that every distinction as such falls
outside being and therefore ceases to be. If we
thus speak of the being of A and of B, then they both seem
each for itself to have one moment that is absorbed into
% --- p. 3 ---
\setcounter{footnote}{0}%
\opage{3}being, and another, on the other hand, that falls outside it. A we can
then think of as divided into a and α, B into b and β; if
a and b now fall within being, then α and β fall outside. Since now a in its distinction from b actually exists,
one must again be able to speak of the being of a and of b, and halve a and b anew; the remaining difference
thereby shrinks to half, this half can again
be halved, and so on to infinity, until being
has extended its power over everything and conquered all distinction.
In the negative form we see at once that, since a and b fall
outside being, they are not at all; they are then
annihilated; one posits and recognizes them in their nothing.

In its first form, then, the extension of being must
be presented as: being + being + \dots{}; in the latter, on the other hand, as being + nothing.

In both forms being is posited as unlimited. But the latter mode of expression is far more determinate than
the first; for while the former merely indicates a thought-experiment,
the latter presents the final result. \textbf{Being thus has nothing} \textbf{outside itself}.

But even the result thereby yielded is burdened
with a great ambiguity, and one such as lies
in language itself. The proposition: nothing limits being,
can express: that being exceeds all limits,
that there is not anything that can set it barriers, that it is positively unlimited; that nothing really
can set being a barrier, that where nothingness
comes to dominion, being must cease. Here, then, a definition of the unlimitedness of being is really given. Indeterminate being is thought in
its whole indeterminateness as such, i.e. its indeterminateness
is posited as its determination. But in order to get some
determination into being, which absorbs all determinations, one must connect a determinate thought with the
% --- p. 4 ---
\setcounter{footnote}{0}%
\opage{4}otherwise indeterminate nothing. Being, which sublates all distinction,
thus cannot after all get the better of the most general distinction-moment,
of negation as such,
of non-being, or of nothing. \textbf{Being is then
not} nothing, \textbf{and} nothing \textbf{not being}. But in its pure immediacy this distinction
cannot hold its ground either. If one tries, e.g., to think
being as the sum of all realities whose oppositions
have been levelled out, thus as pure, solid,
universal reality itself, then nothing is to be represented
as the pure abyss of emptiness. Emptiness thus has
its being alongside reality; but it is characteristic
of empty being that it contains only
itself, and the being of reality likewise contains
only itself; reality and emptiness then cease here to
be oppositions: we find nothing in being and being
in nothing. The question: what is pure being,
demands that being affirm itself through its other.
Being thus here seeks its different determinations:
x + y + z. If we could find these, then we could
also define being. We have already found
them; for it is evident that what is sought in x + y + z must be distinction-moments, and since we now know that all
distinction has vanished, we must put in their place: nothing + nothing + nothing. Since now nothing here quite
generally denotes a missing distinction, nothing in this general form = nothing. Therefore: being
comes through nothing to itself: \textbf{being is
nothing}. The proposition: being is nothing, is the purest
expression of the fundamental law that all unity as such
is its own opposition, i.e. that opposition springs
from unity itself. From the positive standpoint we can now say that pure being here is on the point of
% --- p. 5 ---
\setcounter{footnote}{0}%
\opage{5}restoring all the oppositions it formerly
absorbed, although these cannot yet be characterized
except by the common determination that they are everywhere
to constitute the nothing of pure being. In positing its own nothing, being thereby gives all oppositions
their specific being, and thereby posits itself
as the \textbf{absolute} being. From the negative standpoint, on the other hand, the proposition
expresses this, that pure being as such is a
phantom, a contentless abstraction, and that true
being must be sought in the oppositions from which we have
abstracted, i.e. in the things that are. This thought
consciousness holds fast with ease. When everything that is
is taken away, nothing remains. Being stripped of everything
that is, is then pure nothing. We can now also reverse the proposition and say:
\textbf{nothing is pure} \textbf{being}. However negative pure
nothing may be in all its shapes, it nevertheless shows itself
everywhere as a being. The represented nothing is
the empty place from which everything has been taken away; but if one
now represented the place as taken away, one would have to
represent at the same time the place where the place
had been. Nothing thus continues to present itself
as the place that is, though empty, i.e. as the
being whose concrete determinations have been removed,
i.e. as pure being. This representation is also justified
by strict thinking. Pure nothing, namely, is at home not in the world of sensibility but in
that of thought. But the negative operation of thinking, which sublates all immediacy, cannot come
to pure nothing except by negating immediacy
in its most universal shape, i.e. pure being. Pure nothing is thus the reflex of pure
% --- p. 6 ---
\setcounter{footnote}{0}%
\opage{6}being itself, transferred from the form of immediacy into that of reflection.

Herewith we have arrived at two mutually contradictory
propositions: a) being is not nothing (nothing is
not being), and b) being is nothing (nothing is being). Since we can deny neither of these propositions,
while the one nevertheless sublates the other, there is
no other way out than to press the transition of each into
the other.

When we began to think pure being,
thought was driven over to the opposite extreme, pure
nothing, and in the endeavor to hold this fast,
it was led back again to being. Thought was thus
really in movement, because its object was constantly
moving; but it did not notice the movement as
such. It therefore tried first, quite immediately,
to grasp being as resting in itself, and when it came to
nothing, it likewise believed it could hold this fast in its
purity; now, on the other hand, it knows that it is not to settle
down in either of the categories named. The
new category, which elucidates the transition of being and nothing
into each other, is \textbf{becoming}.

When thought is to hold fast the transition as such,
it must move without interruption between the opposed
terms; the thought of being cannot be absolutely simultaneous
with the thought of nothing. In this oscillation
the distinction of the two moments is preserved. Thus we
now see that being must continually alternate with nothing; thus
we have: being, nothing, being. But the movement still seems merely to belong to thought, while both the
opposed terms rest quietly side by side.

The true transition is first recognized when we bring out the unity of the moments. Being is nothing; in
positing itself, it posits its nothing; its
% --- p. 7 ---
\setcounter{footnote}{0}%
\opage{7}transition into nothing thus rests on its developing its
hidden nothing; just as the transition of nothing into being
follows solely from its having being in itself.

In the series of becoming, then, each of the opposed
terms must be \textbf{doubled}: being (a), nothing (a) \textbar{} nothing (b), being (b). For each of the terms contains
both a posited and a not yet posited moment, when we
present it objectively, or, taken from the subjective standpoint, an overt and a hidden side. If we thus begin with immediate being (a), then
it seems indeed at first glance only to want to posit
itself being (a) = being (a), but if we consider it
more closely, it at once turns its other side forward,
namely its hidden nothing (a). In that nothing (a) comes forth, it thereby becomes itself posited, and can therefore, as
nothing (b), become a new and immediate initial term,
which, itself sprung from being (a), must now also
posit its hidden being-moment, being (b). Although
now being (a) is itself nothing (a), and nothing (a) again
nothing (b) etc., there is nevertheless just as much a distinction
between being (a) and nothing (a) as between nothing (a)
and nothing (b); unity and distinction come and
vanish to infinity: it is a restless unity in
a restless distinction. Because of the opposition that constantly obtains
between the overt and the hidden moment, it is
by no means indifferent whether one begins with being
as posited, so that it can posit its nothing, or with a
posited nothing that is to posit its being.

The emergence of being from nothing is an \textbf{arising}; nothing,
on the other hand, results from being through a \textbf{passing away}. But
in consequence of the identity of the moments every arising
is again a passing away, and conversely.
% --- p. 8 ---
\setcounter{footnote}{0}%
\opage{8}Becoming is thus a concept that really results
from being and nothing. But in that all the moments of the concept are posited and developed, they thereby all come to be present. Being then no longer seeks its nothing, it already has
it; nor does nothing pass over into being, for
since it has posited its being, there is no transition to
make. The restless unity of the moments thereby comes to
rest: becoming is sublated.

\underafsnit{b) Determinate Being.}
\paragraf{12}{Something, Other, Change.}

In \emph{determinate being}, i.e. the double being, nothing is with being and being with nothing. Being, which
has sublated its nothing and integrated itself, thereby
becomes something, while nothing, determined by its
being, now also closes together with itself, and shows itself as \emph{other}. Since it is the same being and the same nothing that have here doubled their
unity and distinction, something is itself other, and other something: the truth lies in the \emph{change} of both. Being, nothing and becoming we can now regard as
the fundamental elements from which all subsequent logical
categories develop; these now follow one another
as ever more perfect expressions for specifications
of positive being and of negation in their
reciprocal transition, delimitation and reflection.

Being, which through its nothing delimits itself,
is determinate being; the moments of this self-delimitation
develop as follows: A) In becoming it was shown that
% --- p. 9 ---
\setcounter{footnote}{0}%
\opage{9}nothing came forth from being, just as being from nothing; we
thereby obtain as the first term: being (a), nothing, being
(b). Through the intervening nothing there arises a
nuance of opposition between being (a) and being (b);
for there is a distinction between the being that immediately
precedes the negation and the being that results from
the posited negation. We thus have here not
simply being + being in the unlimited, as
in § 11, where being in its untrue abstract
immediacy was = being, and therefore also nothing;
but being has here doubled itself: it is in determinate being
that being first comes \textbf{to} being. B) Through
the doubled being the moment of nothingness glimmers
forth and vanishes. Nothing vanishes in determinate being:
for the negation that separates being (a) from being (b)
is here to denote only the distinction that underlies
the mediated unity of the two moments. Nothing
\textbf{thereby itself comes to nothing}. If we analyze the content of this
pregnant proposition, it contains: α) that nothingness wholly ceases and is overcome by determinate being,
β) that nothingness (negation) itself acquires determinate being,
i.e. becomes a determinately existent nothing. This determinate being follows
directly from § 11, where we saw that nothing, just because
it was not being, had being in itself. The only
way in which determinate being can get free of the hollow,
empty nothing is that nothing itself is brought into determinate being.
Just as being went together with itself through nothing,
so now nothing (a) goes together through
being (the negation of nothing) with nothing (b) and
becomes one doubled in itself, i.e. a \textbf{determinately existent} nothing. C) Herewith determinate being itself has been reduplicated; it
must therefore be more closely determined as: 1) determinate being
[being (a) nothing, being (b)] delimited by 2) determinate being
[nothing (a) being, nothing (b);] the first term is here
% --- p. 10 ---
\setcounter{footnote}{0}%
\opage{10}to be regarded as the thoroughgoing position; the second term,
on the other hand, as the thought-through negation; the identity of both first constitutes the full, integrated \textbf{determinate being},
which can more determinately be posited as \textbf{determinate existence}.

In determinate existence each of the developed terms must now
be thought for itself in one thought, seen in one view; thereby
the determinately existent being is to be determined as \textbf{something},
the determinately existent nothing as \textbf{other}. For the development
of something and other, being is our firm basis: abstract nothing has vanished; we have left only being's own
positive oppositions. In something we meet an
immediate \textbf{being}, but in its pure immediacy this is merely the first point of determinate existence, which at once
refers to an \textbf{other} being. If this reference
were now continued to infinity, determinate existence would
never come to be something; when, on the other hand, otherness is retorted,
when it becomes a moment through which
what is doubles and integrates itself, it acquires
its proper limit, and is posited as determinately existent\footnote{In sensation we take up a multiplicity of the moments of determinate existence; but owing to the constant vanishing of the one into the other, mere sensation never arrives at a determinate apprehension of something.}.
Just as the being reduplicated in determinate being
has nothingness as a sublated moment, so
the determinate something too now has the moment of otherness
in it; but so long as other itself is not yet developed, something lacks full determinateness. Other is
indeed precisely the true limit of something; if this is not fixed
as such, then something itself flows away in the stream of determinate existence: \textbf{in order that something} \textbf{be able to maintain} itself,
\textbf{other must become} \textbf{something}.
% --- p. 11 ---
\setcounter{footnote}{0}%
\opage{11}This thought contains a significant pregnancy; for it expresses: that other itself is to support
the positive something, i.e. that something must posit its
hidden other in order through this otherness to affirm
itself. Thus the moments: something (a) (the beginning
something), other (the determinate limit), something
(b) (the something resulting from otherness), belong together
in one single immediacy, in order to produce the determinate existence
of something, or the full-born something; next (b) that
other itself becomes something, that it holds itself in its otherness
as such, and sustains its own determinate existence.

The determinate existence of other is determined by a further development
of the determinately existent nothing. That other occupies
the place of negation is evident at once from its belonging to
thought and reflection. All the moments of mediation
that are required in order to develop something must again
be sublated; in something immediacy predominates; consciousness
of any other whatever is pushed back,
so that something may be seen in its own light. Otherness, on the other hand, presents itself at once in a double light;
it does not testify straightforwardly about itself, but always takes
account of the preceding something. Something must be
the \textbf{first}, then follows the \textbf{other} (secundum, from
sequi). But when other enters, it is really
all over with something; other does indeed, in immediately
positing itself, also become something; but this is
surely \textbf{something} quite \textbf{other}.

The two opposed terms thus seem to keep
quiet, each on its own side. Their mutual indifference
stems from their closed-off being. Something is
only itself, and therefore does not care about other. Other
does indicate an exclusion of something; but since it
does not itself lay claim to being everything, but only an
other, this exclusion is no absolute
% --- p. 12 ---
\setcounter{footnote}{0}%
\opage{12}displacement. Both are, so to speak, decent neighbors who
live in peace because each keeps to its own. The relation,
however, cannot be sublated. Since that whereby something is integrated is precisely an other,
and that whereby other sustains itself is a sublated something, they continually result from each other. Something
is the anticipated other, and other the explicated something.
Something thereby continually becomes other, and other
something. That something is itself its other, and other itself
its something, is comprehensible only in that the other
that lies in something is at the same time the something that develops
out of other. Thus something comes to delimit
itself; it finds again only itself in every other, and
posits its infinity.

But other is itself its something in that it has
something for its other. In that something now comes forth
at every point, other too recognizes everywhere its
own otherness, and likewise delimits itself as
the infinite. Thus determinate existence seems to get
the better of nothing; but it will soon be shown that it is
just as much nothing that has itself come to be, and thereby
destroys positive determinate existence. As the more determinate
expression of the relation, otherness everywhere gains the upper hand.
The opposition cannot vanish without our
sinking back down to abstract being. Something is indeed
itself its other, but not tautologically. The truth is that
something \textbf{changes}. Other is the determinate nothing
of something, just as something is the positive nothing for other.
The more sharply something delimits and excludes other in order
to maintain itself, the more it brings about its
own dissolution, since it is itself its other. Self-dissolution thus belongs to the concept of the determinately existent
something. Likewise, when other has something for its
other, the meaning is not that other immediately is
% --- p. 13 ---
\setcounter{footnote}{0}%
\opage{13}= other; but conversely, that other is distinct from other.
In change the positive identity
of something and other is reduced to being a moment of an
infinite delimitation and breaking down of all limits:
the abiding characteristicon is thus finitude
and transience.

Change is concrete becoming, though with the
modification that, just as becoming most nearly denotes a
coming-to-be whereby determinate existence is \textbf{introduced}, so
change most nearly expresses a destruction in which
determinate existence is \textbf{sublated}.

The movement of change is expressed in the following propositions: a)
from something comes other; b) from nothing (i.e.
from the determinate, determinately existent nothing, namely other) comes
something, and finally c) from nothing comes nothing; for
in change something and other are each other's reciprocal
nothing, and precisely for that reason each a nothing for itself.
In change lies the infinite annihilation of determinate existence. It is no foreign power to which determinate existence here succumbs. Annihilation is its own
innermost essence. It is therefore so far from being the case that
true being should perish in the nothing of determinate existence, that it rather prepares its victory out of
alteration and corruptibility themselves. What
holds of every single moment in determinate existence, namely
that it passes over into its other, must necessarily
also hold of the whole of determinate existence itself; in its infinite
alteration this constantly refers to its true nothing. But this nothing cannot in its positive determinateness
be posited as the fragmentary being-for-other;
for in this form it is precisely determinate existence itself.
If determinate existence ended in such a nothing, it would thereby
precisely go together with itself, conquer its whole
% --- p. 14 ---
\setcounter{footnote}{0}%
\opage{14}otherness, and become the opposite of what it had formerly
been. The nothing in which determinate existence vanishes is
precisely the most solid something, which has gone through
and overcome every alteration (being-for-other),
thus genuine being-for-itself.

\underafsnit{c) Being-for-itself.}
\paragraf{13}{One, Many, Exclusion.}

When determinate being has gone through and overcome
every possible change, it determines itself
as positive being-for-itself, i.e. as mediated
quality. Since otherness has vanished,
something posits itself as the indissoluble One, as the
logical atom. Everything other is thus only a repetition
of the One; otherness and opposition must
therefore be sought in plurality. The One, infinitely
identical with itself, splits itself into the \emph{Many}. But
since each of the Many is itself a One that displaces
all others, while they are all absorbed by the One
of unity, the self-contradiction of quality is herewith
raised to the highest. The for-itself negation
therefore expresses itself as an all-sided \emph{exclusion},
which destroys every qualitative limit.
 The result of the self-development of determinate existence is the perfect identity of something and other. But such a being is no longer determinate being; for its criterion was precisely a duality, a finite opposition between the two terms that are. From this it will
% --- p. 15 ---
\setcounter{footnote}{0}%
\opage{15}be seen that the logical sublation of given oppositions
is both their annihilation and their preservation.
In determinate being the point was to find and posit
otherness at all points; in being-for-itself, on the other hand,
the point is to overcome it in all shapes,
and to show it as a moment in the closed self-affirmation.
In determinate being all determinations were external
and relative, because something and other constantly pointed
to each other; in being-for-itself determinateness is absolute; here negation and delimitation
no longer lie outside, but in that very thing which is
for itself. But determinate existence with its moments is nonetheless
preserved; for a) it is from these moments that all determinations in being-for-itself are formed.
Although the whole infinity of all diversities and changes
here sinks together into immediate unity,
our being-for-itself is not on that account pure being.
In pure being the distinctions were not yet
logically posited; in being-for-itself they are posited as
overcome. In order to think the identity of being-for-itself, one must first think the distinction of determinate being.
b) it is so far from being the case that being-for-itself could immediately displace otherness, that it must rather,
in order to identify it with itself, raise
it to a for-itself other. Thereby
being-for-itself itself gets determinate being. c) As being-for-itself
goes together with itself through determinate being, and
thus raises itself above determinate existence, this happens
solely in that all the moments of determinate existence themselves, at
every point, in every form, are raised to being-for-itself.
All this will become more
evident from what follows.

In being-for-itself we must now, in the first place, not
seek ``\textbf{something}''; for since the for-itself something
% --- p. 16 ---
\setcounter{footnote}{0}%
\opage{16}has gone together with its other, it no longer deserves
this name; just as other too, which has
identified something with itself, loses the character of otherness.
--- Still less can we begin here with
the indefinite article's \textbf{a}; for while it is true that
``a'' is a common determination for something and other,
in this way being is common to all logical
categories. ``A'' is wholly indifferent to the distinction
that lies in something and other; if we do not have \textbf{a}
something, then we have \textbf{an} other; the sublated distinction of both
is on the other hand expressed determinately in the category ``\textbf{One}'',
as the real identity-category. Finitude,
duality, has vanished, in that something and other merge
into \textbf{One}. The unity in which something and other are here
sublated must be regarded from a double side. Partly
it is something that sublates other and makes itself ``one''
with it; partly it is conversely other that makes itself ``one''
with something. We therefore have at once a double One, namely
One (a), which began from something, and One (b), which
began from other. But since all distinction between something
and other is overcome, it is from this standpoint wholly indifferent whether we begin with something
or with other; both the terms named are then in truth One. Thus One (c) must express the unity of both.
Had there now been the least opposition between
One (a) and One (b), then One (c) would thereby have
vindicated for itself a meaning of its own as the higher unity
that sublated both the preceding terms; but since
these are \textbf{for themselves} identical, the last term stands wholly
\textbf{for} itself, as a repetition of what is already given.
The three ones are thus in and for themselves identical, which
can be expressed by a new term, One (d); thus
the One produces itself to infinity.
% --- p. 17 ---
\setcounter{footnote}{0}%
\opage{17}\textbf{One} is here set against \textbf{One}; but the opposition
vanishes like a shadow; for One is - One. It
is by its inexhaustible self-production that the One overcomes
everything \textbf{other}.

The only form under which the opposition
is still preserved must then be sought in the production itself.
The producing One in its being-for-itself can
be opposed only to the whole sum of the produced ones,
as this is for itself: \textbf{the One} can be opposed only to
the Many. Unity and plurality are then the duality
that the One here encloses; we shall now see how
this is developed and overcome.

\textbf{One} is the immediate and determinate expression of unity;
but just because the One is itself the full unity,
plurality, although a product of the One, comes to lie outside
unity. But outside unity the One is at the same time
outside itself. One can indeed say that every produced
One is its own unity; but in this form
it then no longer comes forth as produced in opposition
to the producing One. As a term in plurality
every One is delimited by other ones; in
plurality the One breaks apart its own being-for-itself.
But this going-out of the One from itself into plurality
is precisely its going-back to itself and to unity. That the One has gone out beyond itself means, in other
words, that it has become its own limit.
Outside unity, where we should have expected something
other, we meet the One, only in inverted shape. The One
of unity, which is the principle of plurality, itself constitutes
the first term in the series of plurality. As principle
it has the whole of plurality in itself; as the first term of the series
it is itself one of the Many, and coordinated with any other
whatever. Every term in the series of plurality
is consequently itself pregnant with a plurality. The Many
% --- p. 18 ---
\setcounter{footnote}{0}%
\opage{18}2 are the very explication of the One, and the One the concentration
of the Many. Thus the One of unity presents itself
in a double form, namely as the principle of plurality,
and as that which itself results from plurality. Thereby
unity everywhere takes itself back out of its own
plurality and integrates its being-for-itself. The same now holds of plurality as well. Since
the true unity results from plurality, we must
begin here with this. The many ones in their simultaneous existence thus precede unity: their
sporadic character bears witness to the being-for-itself of plurality. But this being-for-itself is as yet only
abstract. Plurality rests on individual ones,
each of which is a unity \textbf{for} itself. In every isolated
One plurality perishes: the individual One is, as
such, the limit of plurality. But this limit
is sublated by plurality itself in a double way: a) In that
plurality forbids the individual One to remain within
its own unity, it directs its determination
toward another One, and while all the ones are externally
referred to one another, their unity thereby comes outside
them all and outside itself; but this halving
of unity, which takes place in such a way that its one moment lies
scattered as center for the individual ones,
while the other moment hovers over them all, is
precisely plurality. Further, b) the moment of unity
that everywhere constitutes the center of the being-in-itself of the individual ones is, according to what precedes, the principle of
plurality. That plurality thus from all sides forms
unity and results from it constitutes its perfect
being-for-itself. But \textbf{the One} itself seems hereby given up as prey to
the most decisive contradiction. It was to be the
true solid identity, and yet two opposed
% --- p. 19 ---
\setcounter{footnote}{0}%
\opage{19}elements break forth from its midst in closed circles. Its unity has its being-for-itself directly over against
the for-itself plurality. It is, however, so
far from the case that this duality could burst
the adamantine \textbf{One}, that the latter, precisely through the determinate existence of the
opposition, raises itself to the highest being-for-itself. The plurality through which unity
in the One has mediated its being-for-itself is the
very same that is in turn mediated through unity itself.
The self-mediation of unity thus falls within the limits
of plurality, and that of plurality within the limits of unity.
Unity first attains its proper being-for-itself through the being-for-itself of plurality. The result is
thus that being-for-itself develops in a series
of for-itself moments: the one being-for-itself
is only for the other; this infinite self-doubling is the means whereby otherness is here conquered.

In this result the development of quality culminates.
Quality is being-for-itself posited in the form of
immediacy. The moments of opposition that pure
being absorbed are hereby given back. The unity of being
is the principle of infinite specification. The carrying-through
of this specification forms an atomistic
element in logic, which dialectic can never dissolve
without itself vanishing. Every category is a logical
atom whose being-for-the-others is uninterruptedly mediated
through its being-for-itself.\footnote{The category treated here plays a great role in theology, whose task it is precisely to mediate the being-for-itself of God with the being-for-itself of the world. The Hegelian disciples, however, seem not to grant it its due, however masterfully the master himself has developed it in his way. The reason for this is explicable. People have amused themselves so royally at seeing how, by the magic power of dialectic, the one could constantly turn over into the other, so that in the much-praised philosophy of spirit everything has in the end become a being-for-other. --- The principle of the thoroughgoing mythical treatment of the life of Jesus is the idea of self-consciousness's own production of Christ. Christ, the Christ-consciousness, is then to be for human consciousness (i.e. for other) without really being for itself. --- Thus, to be sure, everything becomes quite fluid; the qualities in this speculative fluid have no aftertaste; for there are none. It is pure philosophy, purified of all elements of representation; its pure being is purer than pure water.}
% --- p. 20 ---
\setcounter{footnote}{0}%
\opage{20}This positive form of being-for-itself expresses perfect immediacy. Every moment must be taken \textbf{for} itself; for it is as it is. From
so complete a positivity, therefore, nothingness,
which in determinate being hid itself under the mask of otherness,
seems to be wholly displaced. Where could it now
force its way in, since all places are surely occupied by solid
ones and affirmative qualities? But if negation is not in positive being-for-itself
itself, it must have a being-for-itself. This
it has had, during our development so far, in the observing
consciousness.

When the question was of the \textbf{limit of} the One,
the question was really of its \textbf{negation}; but instead of
the negative operation of self-negation, there presented
itself at once the positive expression of negation, i.e. the Many. Since
the positive moments are thus everywhere posited in
the place of negation, they are mutually each other's
% --- p. 21 ---
\setcounter{footnote}{0}%
\opage{21}negation; being-for-itself must therefore be developed through
all its negative relations and movements.

The One is, as we have seen, the identity of something and
other. This identity has a wholly negative character,
for it comes forth through the change of change.
In change the negation that is was at work
in that something and other were for each other other
and nothing; the one nothing thus alternated with the
other. This alternation bore an external stamp; for
although something became nothing each time it passed over
into other, it was not on that account any the less still a positive
something in itself. Negation was thus merely relative: nothing had determinate being, but as yet no being-for-itself.
In the thoroughgoing change, on the other hand, it is shown that something could not become other, i.e. its
nothing, if it were not already in itself an other. But
in that other is made into something, it becomes other to itself
and thereby nothing in itself. The unity of something and other
is thus to be posited as the doubled nothing in
its most perfect likeness to itself: the positive
One is of negative descent. We can now concentrate
all these moments of negation in the following way:

Something is its other in that it is itself its nothing;
for 1) something is = something; therein lies its being (a):
it is itself and not other. 2) Something is = other:
the being that something has lies not in itself, but in
other. In its other it obtains its positive being (b).
According to b, a must be the first negation of something, nothing (a): something \textbf{is not} what it is; it is in itself nothing,
therefore it reaches out beyond itself in order in b to \textbf{be} what it is not. Since b expresses what something is not,
it is in turn negated by a and thereby reduced to
a nothing b. In order that the double being in something no
longer be burdened with the relative
% --- p. 22 ---
\setcounter{footnote}{0}%
\opage{22}opposition of determinate existence, but merge more perfectly into \textbf{One}, each
of the terms must itself be its nothing; without this reciprocal
self-negation they cannot become positively identical. Thus the otherness taken up into something transforms itself into
a pervasive negation. In that the One is posited,
it acquires at once a positive and a negative side. The positive side comes forth in that being (a)
is identified with being (b); but the identification goes
through negation; nothing (a) is identified with nothing (b).

From this it now becomes clear that the One is its own negation. The negative side represents to us here
retorted otherness. Outside the One, then, there lies no
nothing; for the for-itself negation moves
in and through the One; still less anything whatever
other; for the whole of otherness is enclosed in the negative
movement. The One has gone into itself and has taken everything
with it: an \textbf{outside} no longer holds. But this identity of the positive and the negative
side calls forth the unrest of contradiction in \textbf{the One} itself. The moments that constitute the positive side are the
very same as those that determine the negative. Here
there is not a position + a negation, but a position
that is itself negation, and conversely. If we call,
e.g., the positive side A, the negative B, then the One
can hold itself in the moment A only by
taking up the negative movement of B, and making itself
into what B is; but B is precisely thoroughgoing
negation. Thus A becomes the negative of itself;
positivity falls over onto B, which in turn absorbs the
negative moments in A. During the continued rotation of A and B
the One vibrates uninterruptedly between itself and its
nothing.
% --- p. 23 ---
\setcounter{footnote}{0}%
\opage{23}But this nothing is itself a One, which makes its preceding
One into nothing. Therefore the for-itself negation is to be characterized as \textbf{exclusion}. In exclusion, the sharpest limit of quality, lies
precisely the whole hardness with which closed
positivity separates itself from closed negativity.
This shows itself when we follow it through all
its moments of movement, as a) beginning,
b) reciprocal, c) absolute exclusion. Beginning exclusion is the One's own coming-to-be.
The One posits itself in the moment A, thereby it excludes
the positivity in B; B now vindicates its positivity
and excludes the positive A. Under this dialectical
pulse-beat the One must move, although it stands
still. It moves; for scarcely has it fixed its
positivity before this is at once broken down by negation.
But it nevertheless stands still; for there is and remains only
one single One; when the one One comes, the other is
excluded. Thereby the One enters into a relation to
itself that is no relation, and yet a relation;
for the One that arises always presupposes the one that
passes away: exclusion becomes reciprocal.

Every One thus has the property of positing itself
and excluding another One; but this in turn posits
itself, and excludes another, etc. Therefore the One exists
only in plurality. In plurality all the
individual ones divide between them self-maintenance and exclusion
of the rest; these two activities condition
each other. But the equilibrium of plurality is continually sublated;
for what the one does, they all do. If the one One deduces its existence from itself, then
the rest do the same: all the ones in plurality are
at one stroke. If the one cannot exist except by
excluding all the rest, then these too cannot
% --- p. 24 ---
\setcounter{footnote}{0}%
\opage{24}come into existence before they have accomplished the exclusion of
the one and of one another. This is precisely
what is desperate in the for-itself negation, that it
knows no restriction. In logical despair
the One splits itself apart in order to save its identity
in plurality; but plurality rests on the maintenance
of the One and perishes in its dissolution. Hereby
all qualitative barriers are burst.

What is surprising in this contradiction stems from
a residual immediacy. One concludes according to the principium exclusi medii: either the individual ones are excluded,
or they are not. But such a dilemma cannot hold its ground when exclusion becomes absolute. For according to what precedes, the ones are posited
in their individuality by exclusion itself; when this
is now accomplished on the individual One, then justice is thereby done
to it. The exclusion of all individuals is the position of all;
every One then first posits itself in and \textbf{for} itself
when it posits itself in and \textbf{for} all the rest: absolute
exclusion is an assimilation.

\overgang

When it was said above, in the positive development, that
unity through the being-for-itself of plurality obtained
its perfect being-for-itself, the main thought was
that the one positive opposite was to integrate
the other. Therefore rounded immediacy came forth
everywhere, and everything presented itself in the form of
quality.

Now, on the other hand, since thoroughgoing negation has sublated
the whole qualitative immediacy, the otherness that was pushed back
and hidden has been duly drawn forth,
and being-for-itself has \textbf{out of itself} become a being-for-other. Plurality is again recognized as unity's
% --- p. 25 ---
\setcounter{footnote}{0}%
\opage{25}other; but the One is no longer split apart into the Many,
for its being-in-itself is at the same time its being outside
itself. In that this \textbf{outside} is given back, unity
continues itself by a pliant transition into
other, produces its otherness everywhere and assimilates
it to itself. All specifications are now to be regarded
as extensions and restrictions of one and the
selfsame One. This One is no longer qualitative, for
it has no being-for-itself outside the specifying
movement. Self-extension and self-restriction are
its own immanent determinations. Thus the
stiff, fixed limit of qualitative being-for-itself has now
become elastic and mobile; it therefore determines itself
as

\afsnit{B. Quantity.}
\underafsnit{a) Pure Quantity.}
\paragraf{14}{Continuity, Discreteness, Limitation.}

\emph{Quantity}, the unity of being-for-itself
and being-for-other, is the real limit.
In that the limit develops its identity with the adjacent,
it becomes delimited by itself, and is
in this its self-extension determined as \emph{continuity}. But the middle of the limit is the exclusive negation
of everything adjacent, and thereby its own negative
unity; as such the limit falls apart from itself into
sheer exclusive limit-moments, and is discreteness. \emph{Limitation}, the posited unity of continuity and
% --- p. 26 ---
\setcounter{footnote}{0}%
\opage{26}discreteness, constitutes the real
determinate being of the limit, or its full self-delimitation.

Determinate being was only a being for other; therefore
it passed through change over into being-for-itself. By absorbing otherness, being-for-itself took up
its exclusive negation into itself, and therefore had to go out
beyond itself; quantity, which is the higher unity of both,
is now to develop the higher other of determinate being out of the One of being-for-itself,
and must therefore be determined as the
transition of both into each other, i.e. as the \textbf{limit}. The
immediate or qualitative limit is a negative
middle between two positive opposites: every
quality has its limit. In quality the limit,
as the real middle between something and other, is the pure negation that is in itself. The negative
side of the limit, on the other hand, which turns toward something,
is at once assimilated with the positive something itself,
just as the side turned toward other is absorbed into
the latter's positivity. The limit of quality is therefore positive
determinateness. Determinateness is negation, but a
negation that is absorbed as a vanishing moment in position:
quality is a being-for-itself that must be taken positively, as it is. But if the limit that separates
two qualities is as such neither the one
nor the other of these, then both must, in the limit itself,
be regarded as absent, and that not merely from each other,
but also from the middle of the real limit-line. They then do not border immediately on
each other, but on the common limit. We thus have:
adjacent, limit, adjacent\dots{} to
infinity. --- But just as being (§ 11) passed over into nothing because, as such, it had nothing in itself, and conversely,
so the limit and the adjacent
% --- p. 27 ---
\setcounter{footnote}{0}%
\opage{27}here continually pass over into each other. If we call the limit's being-for-itself a, its being for the adjacent b, the adjacent's being for the limit α,
its being for itself β, then it is shown that it is
the limit that in the adjacent continually borders
on itself. For since b lies between the limit
and the adjacent, it is thereby itself to be regarded as the real negative and exclusive limit,
for which a and α are adjacent moments. But
just because b is now the for-itself limit,
α becomes the limit-negation between b and β. In this
way the limit continues to shift itself in infinite
progress. This, then, is characteristic of the limit,
that in exclusive fashion it is always both in and outside
itself. Limit-mediation makes the adjacent
\textbf{into} limit in the same instant as it makes the limit
into an adjacent. In this self-extension of the limit quantity is posited as \textbf{continuity}.

In continuity it is shown that the limit cannot
be anything stationary. Each time the limit is posited
with a certain immediate determinateness, it is surrounded
by an adjacent; but this is, as we have seen,
itself limit, and consequently the limit itself. Therefore: only
by the limit's going through all its delimitation-moments, positing and sublating all its determinations,
can it be grasped as true limit. But in
the self-shifting of the limit there lies at the same time an infinite
repetition of its exclusive negation. What
is first determined as the pure for-itself limit
becomes in the next moment an adjacent; but every
adjacent is likewise a for-itself limit:
the one limit thereby excludes the other; but since
each limit is determined in the same way, then
limit is always = limit; the limit thus excludes
% --- p. 28 ---
\setcounter{footnote}{0}%
\opage{28}itself, posits itself thereby as separated from itself:
in this completed self-separation the limit is discreteness. The pure limit then dissolves into a plurality
of limits, such that each in turn contains for itself
a limit-plurality. In the limit there is at
once both continuity and discreteness, or, more precisely expressed, the limit is itself the true middle of continuity and
discreteness. It is this sharp separation
of both moments that underlies
the Eleatic presentation of the logical impossibility of movement. One can, namely, in the line of continuity posit
any two discrete points, e.g.
A and B in the line \textbf{A---c---B}; the line itself is wholly indifferent
to these; for it can be extended beyond
both: the discreteness of the points is posited in external fashion in the flowing continuity of the line.

The question of the possibility of movement is now reduced
to the question of what conditions its result. Before he who moves can come from
A to B, he must first have come half the way,
namely from A to c. But the distance from A to c can
again be halved, and the halving of the halved distances
continues without cease. The absurdity of movement
then lies in this, that even the shortest way cannot be covered
before a still shorter one has been run through: movement
presupposes itself to infinity. The result
of the very first movement would thus have to be: 1) a
\textbf{determinate} continuity (i.e. continuity with a finitely
small extension); for without a certain distance between
starting point and end point movement
cannot be thought; 2) an \textbf{indivisible} continuity (with an infinitely
small extension); for otherwise the movement
would not be the first. The contradiction thus lies
in this, that one demands a continuity without discreteness,
% --- p. 29 ---
\setcounter{footnote}{0}%
\opage{29}while one nevertheless proceeds from the consideration that
every continuum can be separated into discrete moments.
Thereby one comes to think movement outside
movement itself. For prior to the very first starting point
no movement can be given, since
it is the first; but when movement stands still at
the very first point, then it can surely never reach the
second. Thus one comes also to think
rest outside rest itself. Every rest is the result
of a preceding movement. Movement cannot
now come to rest at the second point, because it
constantly rests at the first; but neither can it
rest at the first, because no movement has gone before it.

The truth is that completed movement contains
continuity and discreteness in indissoluble unity.
The Eleatics have merely confined themselves to demonstrating both
opposites, and instead of recognizing the transition of the one into the other, so as thus to comprehend movement
from movement itself, they have pressed continuity and
discreteness each for itself, and reduced movement to
a discrete and consequently absurd standing-still. --- Only movement can sublate the sharp opposition between
continuity and discreteness. Just as becoming was an
uninterrupted transition from being to nothing, and thereby
at once affirmed and negated both moments, so movement is now a constant hovering
between standing-still (the moment of discreteness) and self-shifting (the continuity-moment), and consequently the medium
in which both are posited.

The result of movement is thus a continuity determined by
discreteness, delimited by a discreteness determined by
continuity, or, in other words, the determinate existence of quantity as quantum in continuous
% --- p. 30 ---
\setcounter{footnote}{0}%
\opage{30}and discrete magnitude. Discrete magnitude, since
it has continuity in it, is a plurality of separate continua, just as continuous magnitude is a unity of
several discrete moments. If one now tries to transform
discrete magnitude as such into continuous,
or conversely the latter into the former, then one comes to
the antinomy of infinite divisibility.

a) Continuous magnitude can be divided; for it
has the moment of discreteness in it, and each of the discrete
parts are themselves continua. By division, then,
continuous magnitude is dissolved into a plurality of continuous magnitudes, which again each for itself can be divided:
continuous magnitude is \textbf{infinitely divisible}.

b) Since the result of every division is a continuous magnitude, it is so far from the case that the continuation of division
could sublate the continuity of the magnitude,
that it rather becomes a repeated confirmation
of it: continuous magnitude is consequently absolutely
\textbf{indivisible}.

This contradiction one is now accustomed to veil
by presenting every magnitude as consisting of
infinitely many, infinitely small parts. Such an explanation
displays an infinitely beautiful mathematical balance,
in which the one thoughtlessness outweighs the other. One is in the first place to think that an infinite division is completed, that is, comes to an end; next, that the
parts that are yielded by the infinite division are
infinitely small (- 0), i.e. \textbf{no} parts.

It becomes evident from this that divisibility and
indivisibility mutually condition each other. The parts are
themselves magnitudes and the magnitudes parts. In that the
one magnitude is delimited by the other, it is delimited in itself by
its parts, and thus limitation,
or infinite delimitation, is realized.
% --- p. 31 ---
\setcounter{footnote}{0}%
\opage{31}The elements of limitation are sheer magnitudes and
sheer parts. Each part is as such a singularity;
each magnitude, on the other hand, a composite. But since the part
is itself magnitude, and the magnitude part, then
the composite passes over into the single, and the latter into the former. % print: enkelte. (stray stop) (misprint)
Each single thing must as such be posited in its positivity,
like each of the ones in being-for-itself. But
here the single One is posited with the express determination that it is to be delimited by another single
One. The point here is then not, as before, the return of unity
to itself through plurality,
but a \textbf{determinate} plurality and a unity delimited by plurality. The limiting unity and plurality
are determined by a threefold act: 1) each of
the terms of plurality is separated and posited for itself; 2) the posited
terms are taken together, whereby magnitude, as encompassing
unity, becomes composite; 3) the comprehending
unity is again posited for itself as a new
singularity. In this separating and taking-together
limitation is numeration.

\underafsnit{b) Number.}
\paragraf{15}{Number-Unit, Amount, Innumerability.}

In limitation the unity of continuity is delimited on all sides by the plurality of discreteness; the mutual
determination of both is to be posited as number. In the continuity of number the discrete plurality is at
once absorbed and excluded; with this determination
continuity is \emph{number-unit}, and the discreteness
of plurality is \emph{amount}. Since each of the ones of the amount
% --- p. 32 ---
\setcounter{footnote}{0}%
\opage{32}is a number-unit that contains its amount, then
number moves infinitely within itself, and yet goes
everywhere outside itself: \emph{innumerability} is the self-dissolution of number.

From the beginning of the logical development a
determination of numeration has hovered before thought; for without
such, no repetition, no opposition can
be held fast. The moment of numeration expressed itself clearly
in the \textbf{other} (secundum) of determinate being, still more clearly
in the \textbf{One} and \textbf{Many} of being-for-itself. But with categorical accent numeration could not come forth in quality.
In the determination of other the point was really
to show what it was in itself, as a being
(alterum, aliud), in opposition to something; its determination as a secundum was, on the other hand, wholly external
and fleeting. Even in being-for-itself the concern was
really only the qualitative genesis of the ones. The One of being-for-itself
was therefore most nearly to be regarded as a
One of identity; even plurality was qualitative, for
it was determined merely as the sublated other of the for-itself unity; therefore the question was not yet asked:
\textbf{how many}? Now, on the other hand, since unity has become a
continuity that has taken up plurality into itself as its own
discreteness, quantity posits itself as a unity perfectly
determined through its discrete ones.
Hereby the logical movement of thought comes as near as possible to illogical
immediacy; for in numeration it confines itself to a pointing-out of the individual
ones, such that each is held fast for itself and
only in external relation to the rest. The result of numeration is the concept of number. Admittedly no numeration can in reality
take place without the determination of number hovering before consciousness, but in the
% --- p. 33 ---
\setcounter{footnote}{0}%
\opage{33}immediate counting: 1+1+1 etc., each of the ones \textbf{is}
still most nearly apprehended as qualitative being-for-itself; only through the comprehending operation of counting-up does
the determinate concept of a unit, a two etc.,
i.e. of \textbf{number} as such, come forth. The most perfect
being-determination at which we have now arrived must then
be posited as a number-determination. All determinate existence is \textbf{a}
number\footnote{From this standpoint the Pythagorean philosophy must be judged.}, for it is a unit that repeats itself
to infinity, but with each repetition obtains a
new and wholly specific meaning. Every unit is
immediately one with itself, but in this being-for-itself
it has its peculiar place in the series of ones. What is characteristic of the number-unit is precisely this, that
with all its variety it is nevertheless in every moment so absolutely
delimited. Insofar as the number-unit is regarded as
such, e.g. a five as opposed to a six, it is absolutely
identical with itself, and has excluded all plurality.
But this exclusion is just as much a sublation
and taking-up of the ones of plurality; for only in
relation to a delimited plurality can unity be determined.
Only by positing 1 + 1 + 1 + 1 + 1 do we obtain
the unity 5. But in the moment of plurality each
of the ones is to be posited for itself and to follow a necessary determinate
order; in that plurality is thus encompassed
by unity, each of its ones is counted.
These counted ones of plurality now constitute the second
main side of number, namely \textbf{amount}. The amount is precisely
the limit between unity and plurality; for while
on the one side it points to the discrete ones,
on the other side it expresses their
% --- p. 34 ---
\setcounter{footnote}{0}%
\opage{34}taking-together in the number-unit. The amount always states the reciprocal
determination of unity and plurality: so great
a unity, so great a plurality\footnote{The peculiar development of unity and amount is given in the deduction of the four arithmetical operations, but here the number-determinations pass over into mathematics' own specifications in such a way that they can no longer become moments in speculative logic. A concrete development of numerical relations in their various transitions would require a thoroughgoing philosophy of mathematics. For while mathematics immediately presupposes its peculiar axiomata and by stringent definitions of the understanding characterizes its problems and justifies its procedure by the operations themselves and their results, the philosophy of mathematics would, on the other hand, have to deduce the peculiarity of the operations and of the results from the very concept of quantity itself. Mathematics obtains its exact character by the sharp delimitation with which it keeps all heterogeneous elements away. But precisely the circumstance that mathematical results can be applied to real determinate existence shows that the number-determinations are abstracted from a more concrete world. A philosophy of mathematics ought therefore to demonstrate the one-sidedness and contradiction with which all mathematical determinations are burdened, and thereby make it evident that this science, however independent it may seem to the immediate understanding, nonetheless dialectically sublates itself in order to become a moment in an all-embracing cognition of the world.}.
In the amount, therefore, the inner contradiction of number comes
clearly to appearance. For unity, plurality and the counted ones
each have in the amount their abstract being-for-itself,
while they are nevertheless only for one another.
The number-unit is given in its simplest and most general
form in the number 1. But 1 is always = 1.
This tautological likeness of unity with itself
% --- p. 35 ---
\setcounter{footnote}{0}%
\opage{35}33
makes it necessary that one, in order to \textbf{come} to a development
of the ones of plurality, must take to one's aid another determination which, moreover, by its nature lies wholly
outside the concept of number. For the thought of 1+1 demands that one separate 1 from 1. In the sphere of
representation this separation takes place very
easily, since one here generally counts the determinately existent
objects themselves. But if one now abstracts from
all specific qualities, then one comes to posit
1, 1, 1\dots{} This positing looks indeed at first glance % print: Diekast (for Øiekast) (misprint)
like a repetition, but according to the pure
concept of unity it really is not one. Thought cannot separate
the one 1 from the other if one does not look about
for heterogeneous criteria, e.g. time and succession. But succession lies wholly outside the concept of
number. That 1 is a sign that can be written in several different
places, that it can be uttered by repeating
the same word several times, does not concern abstract
unity as such. When, therefore, a real
numerical separation is to take place, one must have a certain
quality, no matter which, in mente. This quality
is now to be regarded as the other of quantity, as the
negation through which it takes itself back.
Herewith, then, number goes outside itself and thus realizes
the first main requirement of quantity.

But quality itself obtains no qualitative
determination whatever through numeration itself\footnote{Hence it also comes about that one cannot add up different qualities, e.g. 4 tables and 3 chairs. By such an adding-up one would confuse quantity-determinations with quality-determinations, which, according to the nature of quantity, is absurd.}. As it is
% --- p. 36 ---
\setcounter{footnote}{0}%
\opage{36}to number wholly indifferent whether we count the stars in
the sky or the grains of sand on the seashore, so is
it also indifferent to each single star whether
it is posited in the amount as the 1st or as the 10th. On this indifference of number toward everything other is grounded the innumerability of the qualities of determinate existence. The things in the world can indeed be counted, but only on
the presupposition that they present themselves with external,
contingent restriction and delimitation. One can
count the buttons on a coat; but there are in the world never
so many coats and never so many buttons; there
could, after all, as far as these immediate qualities themselves are concerned, be many more. That God, on the other hand, has numbered
the hairs of men's heads is comprehensible only in that
he who posits the measure of things at the same time posits the deeper
unity of quality and quantity. Were he,
on the other hand, first to take the hairs on our heads as given, and begin to count 1, 2, 3, then numeration
would never be brought to an end. From this it is evident that
number negates its otherness in order to determine itself;
quality and quantity cannot delimit each other:
\textbf{number counts only} \textbf{itself}. Thus we find % print: sinde (for finde) (misprint)
the second main requirement of quantity, namely that its
being \textbf{outside} itself is to be precisely its own being in
itself, determinately expressed.

But in that number excludes its other, it gives itself
its otherness. It is this otherness that we shall now
consider more closely.

Since each number is so immediately itself that one
can come from one number to another only by the admixture
of heterogeneous determinations, the
individual numbers on this account hold themselves in abstract external
opposition to one another. The abstract being in
itself of number is its equally abstract being \textbf{outside} itself. This
% --- p. 37 ---
\setcounter{footnote}{0}%
\opage{37}externality of number knows no limit. Every
One in the plurality of the amount can be posited as number-unit for
a new amount to infinity: \textbf{the numbers are innumerable}.

The \textbf{innumerability} of number contains: a) the infinite reproduction of number-determination; b) the
vanishing of number-determination; c) the logical sublation of number.

a) In the specific numerical expressions quantity is made finite
and defined by quantum + quantum +\dots{}
The definitive quanta can be counted; but however long
one continues addition and multiplication,
the finite quanta never come to an end. Quantity,
which sublates every limit, constantly meets new limits,
and thus never comes to accomplish the sublation of the
limit; it oversteps every amount and yet
never reaches out beyond the finitude of numbers. Innumerability,
as the unceasing reproduction of definitive number,
is hereby posited as the omnipresence of the limit in quantity, as \textbf{infinite finitude}.

b) When quantity is expressed in a sum of
finite quanta, then one would expect that in this
perfect delimitation there was posited a determinate limit
for a maximum and minimum; but the greatest in finitude has a finite greatness; there is therefore
always a still greater; the smallest finite magnitude
has a limit to its smallness; beyond this limit
is found a still smaller magnitude. Further, there lies
between two given finite magnitudes an infinite
amount of other finite magnitudes. Between 1
and 3 lies not merely 2×1, but 4×½, 8 ×¼ etc.

Greatness and smallness are thus in infinite finitude wholly relative magnitudes. A minimum and
maximum can by external comparison here be posited
and sublated everywhere. Thus the discreteness of numbers flows away in the stream of continuity. Innumerability, as
% --- p. 38 ---
\setcounter{footnote}{0}%
\opage{38}the dissolution of number, is a sublation of all posited
limits, \textbf{an infinity resulting}
\textbf{from finitude itself}\footnote{As an example taken from mathematics of quantitative infinity in its double form, one may cite the quotient of $\frac{1}{1-x}$, which (on the presupposition: $x > 0$ and $< 1$) gives $1 + x + x^2 + \dots$. The continued development of the quotient makes evident at once the relation between amount and innumerability and also the contradiction that lies in the approximation to the infinitely great and the infinitely small. At whatever term one breaks off the division, this breaking-off shows itself to be arbitrary. With each new term we obtain an increased definitive numerical value: but it is nevertheless always too small; with every continuation of the division we approach a more adequate expression: but the adequate limit itself flees unceasingly further away. The absolute limit can therefore be said to be just as far away from $1 + x + x^2 + \dots x^{1000}$ as from $1 + x$, whereby the whole approximation, regarded from the absolute standpoint, loses its meaning. This is, as Hegel too has developed it, the contradiction in the infinite progression: ``Das Ovantum hat es in seinem Begriffe ein Jenseits zu haben. Dieß Jenseits ist erstlich das abstracte Moment des Nichtseyns des Qvantums: dieses lost sich an sich selbst auf; so bezieht es sich auf sein Jenseits als auf seine Unendlichkeit, nach dem qvalitativen Momente des Gegensatzes. Aber zweitens steht das Qvantum in Continuitet mit diesem Jenseits; das Qvantum besteht eben darin, das Andere seiner selbst, sich selbst äusserlich zu seyn, also ist dies Aeusserliche ebenso sehr nicht ein Anderes als das Qvantum: das Jenseits oder das Unendliche ist also selbst ein Qvantum. Das Jenseits ist auf diese Weise aus seiner Flucht zurückgerufen, und das Unendliche erreicht. Aber weil dieß zum Diesseits geworden wieder ein Qvantum ist, ist nun wieder eine neue Gränze gesetzt worden; diese, als Qvantum, ist auch wieder von sich selbst geflohen, ist als solches über sich hinaus, und hat sich in sein Nichtseyn, in sein Jenseits von sich selbst repellirt, das ebenso perennirend zum Qvantum wird, als dieses sich von sich selbst zum Jenseits abstößt''. (Wissensch. d. Logik. Werke 3 Band p. 265). If the development of the quotient is to be adequately concluded (as by $x^n + \frac{x^{n+1}}{1-x}$), then the progression is hereby reduced to the infinity lying in the relation between dividendus and divisor, and the amount is posited as one with innumerability. In this way true quantitative infinity is given in the expression $\frac{1}{1-x}$ itself. For since dividendus and divisor, each for itself, are determinate finite numbers, their mutual relation is here a determinate, constant relation. But this relation, delimited in itself, is precisely the source of the unlimited numerical magnitude that is evolved in the quotient. The expression $\frac{1}{1-x}$ thereby takes on the character of an intensive magnitude, whose extensity lies in the very evolution of the quotient.}.
% --- p. 39 ---
\setcounter{footnote}{0}%
\opage{39}c) In the mediation of amount and innumerability
quantity becomes master of its limit. With the innumerability
of the amount it sinks down into sheer delimited quanta
and thereby \textbf{posits} \textbf{all limits}. The delimited quantum is finite, since it has a limit and constantly
points out beyond it in order to obtain its determinateness in another
quantum. But since this transition into other
lies in the quantum's own determination, since this other
is itself quantum, the delimitation is at the same time the sublation of the limit. Finitude mediated with itself
is then infinity itself. Through this mediation
of finite quanta quantity takes itself
% --- p. 40 ---
\setcounter{footnote}{0}%
\opage{40}back, and \textbf{sublates all its limits}. Insofar
now as all delimitations are posited and taken up in
quantity, the significance of numbers is preserved; but in that
the limit is everywhere sublated, number is
reduced to a vanishing moment. Instead of the immediately delimited number-unit,
we then get a quantity-unity unlimited in its self-delimitation, i.e. \textbf{an intensive magnitude}; and instead of
the amount a quantity-evolution in which innumerable
quanta uninterruptedly fix and sublate all possible number-determinations,
i.e. \textbf{an extensive magnitude}.

\underafsnit{c) Absolute Quantity.}
\paragraf{16}{Intensity, Degree, Approximation.}

By the infinite position and sublation of the limit, quantity moves within itself, and is absolute quantity. When the discreteness of amount is posited as a sublated moment in the continuity of innumerability, the being-in-itself of quantity is determined as intensity. Each of the moments of intensity has all the others sublated into unity with itself, and is consequently itself an intensity. The relation of the intensities is expressed in \emph{degree}. Degree is the qualitative determinateness of quantity, but still abstractly quantitative. The specific delimitation is reached only by an \emph{approximation}, whose insufficiency compels quantity, closed in itself, to go together with its qualitative other in modality
% --- p. 41 ---
\setcounter{footnote}{0}%
\opage{41}. In the dialectical development of the categories we can distinguish: a) the successive alternation of the opposed terms, which forms a continuous series, while it, so to speak, advances in a straight line; and b) the necessary emergence of the one opposition from the other, whereby the speculative unity of both is cognized, and the movement takes on the character of a circulation. Thus the series: being, nothing, being\dots{} can be continued without end; but in becoming it is cognized that nothing emerges from being itself and being from nothing. If one has thus seen through the constant relation between the first two terms of the series, then the meaning of all the subsequent ones is thereby grasped; for however often they are repeated, they are still only a tautological repetition of the first two. The problem that the rectilinear movement calls forth is always solved in the circular movement. As long as quantity is merely apprehended under the external opposition of limit and limitlessness, of amount and innumerability, this category contains an insoluble contradiction. One then holds fast to the demand that quantity is to be determined by discrete quanta, and that these must be stated in numbers. But precisely because one is caught in the finitude of amount, one is just as immediately overcome by innumerability and infinity. Stated determinately, the contradiction runs thus:

: Number has meaning; things must be counted if they are to be grasped; but now they cannot be counted: as surely as space is infinite, so surely are the \emph{quanta} of quantity innumerable\footnote{Cf. the Indian myth of Brima.}.

When, on the other hand, one sublates the series and makes clear to oneself the constant relation between the determinate
% --- p. 42 ---
\setcounter{footnote}{0}%
\opage{42}quantum and the indeterminate quantity, between amount and innumerability, then the abstract opposites will sublate one another and the contradiction will vanish. By the finite quantum's delimiting itself, the movement of quantity is reduced to a circular movement. This now completes the delimitation; for the one finite quantum leads continually over into the other quantum (a) over into quantum (b), just as being into nothing but every barrier is at the same time lifted; for the finite quantum finds itself again in its negation [quant. (b) leads again over into quant. (a)]. Quantity hereby closes itself off as absolute \textbf{quantity}. Therefore it is now no longer necessary either to step from amount to amount in order to count the innumerable. The one amount is delimited by the other; but amount is = amount; \textbf{amount is its own} \textbf{other}, it goes together with itself: therein lies the mediated innumerability, which consequently cannot be confused with blind indeterminability. By this circular mediation the intolerable length and breadth of quantity is sublated. The eternal fundamental relation between amount and innumerability, in which its disjecta membra, the sporadic quanta, are concentrated, can now be seen in one view and grasped in one speculative thought: \emph{intensity} is the truth of quantity.

In absolute intensity, now, all immediate discretenesses must in the first place vanish: its intension is that they should all be in one another and go together into one; but by their complete vanishing intensity itself would vanish; therefore it has at the same time a decided tendency also to posit the being of the finite quanta outside one another and to give them an extension. It is this ceaseless tension between intension and extension that maintains the limit and the infinite alternation of amount and innumerability.
% --- p. 43 ---
\setcounter{footnote}{0}%
\opage{43}Under extension the one finite quantum is delimited by the other; but this mutual delimitation, which forms extensity, is only relative; in intensity, on the other hand, all quanta have their common absolute limit: as many sublated limits in intensity, so many posited limits in extensity. In this way amount is fixed. But since intensity has sublated all possible limits, and thus has them in itself, every possible delimitation must also be able to be posited in extensity; thereby innumerability is preserved.

Although intensity is now the true limit for all extensive quanta, it nevertheless cannot have its being outside these; for in that case it would itself be reduced to an extensive moment; everything would then thereby be absorbed into extensity, and extensity itself would have to vanish, since it cannot subsist without intensity. As a consequence of this, intensity must sustain discreteness and continuity by always positing a \textbf{discrete} moment as point of departure for the \textbf{sublation} of all the others. If, then, regard is had to a single determinate moment by itself, it lies in the nature of the matter that intensity is opposed to extensity. For the intensive strength of the single moment consists in its having all the other moments in itself. But according to extensity the moments are outside one another. It is consequently a plain absurdity that a single moment should, in one and the same sense, both have and not have the others in itself. From this standpoint one concludes rightly: \textbf{the greater the extensity}, \textbf{the less} intensity. The moments that in extensity have their being outside one another are, however, the very same as those that in intensity have their being in one another;
% --- p. 44 ---
\setcounter{footnote}{0}%
\opage{44}extensity is consequently intensity's own discrete and distinct expression: \textbf{the greater the intensity}, \textbf{the greater the extensity}. Since intensity is now to acquiesce in a single discrete moment while it sublates the discreteness of all the other moments, and each single moment in intensity and extensity must here have the same right, it holds of them all that each of them by itself is to be posited in its discreteness by the sublation of the others. On account of this, absolute intensity must extend itself in a \textbf{series of} intensive \textbf{moments}. Every moment is an intensity, but only in relation to the other sublated moments. According to the mutual distinction of the moments, the sublation must take place in a different way in each single one of them. Thus the discreteness of the intensities develops. But this difference is still only quantitative. This quantitative distinction cannot be sought in the greater and lesser extent of the moments; for their immediate extension has vanished. The extensity of the single discrete moment lies solely in the sublation of the others, and is consequently its very intensity. Just as little can the amount of the sublated moments form any difference; for when any single one whatever is posited in its discreteness, all the others are thereby posited as sublated. The amount of the sublated moments is always equally large. The distinction can then arise only from the \textbf{developed relativity} of \textbf{the discrete} intensities.

Every intensive moment sublates an innumerability of moments, in which it itself is to be sublated innumerable times. Insofar as it has all the others in itself, it does contain the whole intensity; but since it in turn falls back, as sublated, to the intensive discreteness of the others, it is only a relative expression for absolute intensity.
% --- p. 45 ---
\setcounter{footnote}{0}%
\opage{45}Absolute intensity then posits every relative intensity as a sublated moment for another relative intensity: the more sublated intensities the relative intensity contains, the more it approaches the absolute. Hereby absolute intensity unfolds itself in a series of \textbf{lesser} and \textbf{greater} intensities.

On account of this one must distinguish between a twofold sublation of the discrete moments. a) The general sublation reduces discreteness to its origin. It posits the continuity and unity of intensity as the primary, the discreteness and plurality of extensity, on the other hand, as the derived and secondary. In the universal sublation the discrete quanta as such are not yet posited. But only when something is posited can it actually be sublated. Therefore there now follows b) the specific sublation, which presupposes certain discrete moments as already posited, in order, by their absorption into a subsequent moment, to form an actual relative intensity. The peculiar strength of the relative intensity thus depends not on how many moments it contains according to the general sublation (for it contains all), but on how many it has in itself according to the specific sublation, that is to say, how many relative intensities it has under it. The first and weakest intensity would now have to arise by the first discrete moment's being posited through a general sublation of the others; but since the relative intensity is never formed by any merely general, but only by a specific sublation, the weakest intensity would be no intensity at all. The intensities are on account of this innumerable, the transition from the lowest to the highest is infinite; the truth of the movement lies in gradation: \textbf{the series of relative intensities} \textbf{rounds itself off} in \textbf{the circle of degrees}.
% --- p. 46 ---
\setcounter{footnote}{0}%
\opage{46}Degree posits intensity with its determinate relativity. Every degree is as such a quantitative something by itself. But the single degree, which has sublated other degrees in itself, is according to its determination to be delimited and sublated in turn by a higher degree: its being-for-itself refers continually to another being-for-itself. The scale of degrees thus presents us with a series of for-itself moments, which are nevertheless all for one another mutually. This being-for-itself is now quantity's own quality, and this quality the mediated relativity. The transitions and nuances of gradation are in and for themselves innumerable; but relativity as such is the abiding determinateness that constitutes what is characteristic in all possible gradations, and therefore sublates the infinite series of degrees. Gradation expresses precisely the inflexible demand that there be oppositions, fixed and unshakeable difference-points in the true, absolute quantity: \textbf{degree separates itself from} \textbf{itself in order to} \textbf{refer itself to} \textbf{itself}. Since the opposition belongs to degree's own self-determination, there must be contained in the lower degree a moment that does not let itself be sublated by the higher: there must be a something (a quale) that stabilizes the distinction between the single degrees. Under this positive specification degree has its other in itself; it lies in its determinate existence that it must change. The single degree, which is itself something by itself, is thereby at the same time a degree of something that it cannot completely express. This something then reaches out beyond the lower degree, effects its sublation, and specifies itself in a higher degree, which thereby obtains its being-for-itself. But however often the sublation is repeated, the positive something will still always place itself outside the single degrees. Therefore it also shows itself with a twofold
% --- p. 47 ---
\setcounter{footnote}{0}%
\opage{47}character. It fixes in the first place the limit between degree and degree, and next forms the common determinate limit for all degrees. It is this something whereby degree becomes its own other, and thereby at the same time something other than itself. Since each degree is something by itself, change is a transition from the \textbf{one} degree to the \textbf{other}; but as every degree develops itself as a degree of the other degrees, degree is always reduced to the tautological determination of being a degree of itself. The change of degree thereby sublates itself, and gradation loses its meaning. But when degree thus determines itself as a something without gradation, it is no longer degree; as the degree-change stops, degree has absolutely changed and has become something quite other. Yet as degree loses itself in other, it gains itself; for this otherness becomes the proper basis whereby discreteness is restored, and gradation comes into its own. Each single degree now sustains its specific meaning, and the degrees become one another's mutual other, and set one another limits in relation to the qualitative other, which posits the absolute limit of all. This other, however, cannot so appear in the single degrees that it is itself specified; for thereby the degrees would be separated into nothing but qualitative differences, and discreteness would become so predominant that continuity and with it discreteness itself vanished. From this it follows that \textbf{the quality} \textbf{that underlies} \textbf{the actual gradation} \textbf{is not itself} \textbf{subject} to gradation\footnote{From this standpoint the Stoics denied degrees in virtue and crime.}.
% --- p. 48 ---
\setcounter{footnote}{0}%
\opage{48}But when the specification is owed to the quantitative something, and this something, as quality, falls outside gradation, then gradation falls outside itself; for it cannot do without the moments of specification. Degree can therefore be regarded as the general negation and limit by which quality completes its own specifications, since every term in the series is determined by a degree. In this way quality comes to develop itself in a gradation; but as quality turns over into its other and becomes degree, that which was formerly posited as degree turns over into its other and becomes quality. The result of this is then this, that the change of degrees or the true gradation is at once \textbf{an alternation of degree} \textbf{with degree and of} \textbf{degree with degree-negation}.

Herein lies a contradiction that comes to clear development in \textbf{approximation}.

Approximation posits the degree of something (quantity with quality) as point of departure. But, not content with the way in which the qualitative something is expressed in a certain degree, one strives forward toward a higher degree, which is to give the quality a truer determination. One does not, however, want another something, does not want to pass over to another quale; the otherness one seeks solely in the changed degree. Approximation \textbf{is} a continued \textbf{gradation}. What occasions the gradation is solely the demand that the quality in question is to have its abiding, adequate expression. Since degree is here the only determination of quality, the given quale must receive different determinations in the different degrees. By gradation, then, quality itself undergoes a specification. The qualitative something would obtain its absolute expression only when the highest degree was attained; but degree \textbf{is}
% --- p. 49 ---
\setcounter{footnote}{0}%
\opage{49}quantity, the amount of degrees innumerable: the adequate determination of quality is consequently never attained. If there is now a relative distinction between the qualitative something posited in a lower and in a higher degree, then there must be an infinite distinction between its determination as it can be given in relative degrees and as it was to be given in the absolute degree. In approximation there thus lies the insufficiency of gradation; the qualitative something that one strives to express in the relative degrees thereby shows itself as a qualitative other that lies outside all degrees, i.e. as an unattainable goal\footnote{Subjective idealism, which postulates immortality on account of the endless strife between reason and drive, and panlogism, which urges the infinite process of personification of spirit but denies God's absolute personality that is for itself, both find themselves in the unresolved contradiction of approximation.}. Precisely because, on the standpoint of approximation, one so immediately trusts oneself to gradation and places the qualifying otherness in degree itself, one finds oneself just as immediately compelled to regard the delimiting other as lying outside every degree. When, on the other hand, it is brought to consciousness that the other of degree is everywhere both quantitative and qualitative, and consequently lies both in and outside degree, then the goal shows itself at once as both absent and present (quod petis hic est): the true approximation is an infinite evolution.

\overgang

Quality develops itself as a relation between for-itself moments. The qualitative limit is fixed and immovable; it cannot be changed without
% --- p. 50 ---
\setcounter{footnote}{0}%
\opage{50}quality itself being changed: : quality is as such the fixated determinateness.

Quantity, on the other hand, is the self-delimitation of the limit, which has its purest being-for-itself in number. But the self-delimitation of the limit is the sublation of every determinate limit, and number reproduces itself to infinity. Quantity is the movable limit; indeterminate increase and diminution is its determinateness. Quality and quantity, then, regarded in their immediacy, are indifferent to one another and therefore cannot delimit one another.

But this indifference, which stems from their immediate, categorical distinction, sublates itself.

For as quality becomes its own other, it finds itself again on the far side of the fixated limit; its self-delimitation is therefore its indeterminate self-expansion and self-restriction; it thereby becomes something other than itself and is determined as quantity. The fundamental movement of quantity is limit-sublation. But the limit cannot be lifted without its having first been posited: continuity requires discreteness. In the posited limit and in discreteness as such we thus find again the qualitative determinations in quantity. As quantity qualifies itself and posits its other, it becomes, as approximation showed, the negative of itself, i.e. quality. The result, then, is that quality mediates its self-delimitation through quantity, and that quantity delimits itself by quality.
% --- p. 51 ---
\setcounter{footnote}{0}%
\opage{51}\afsnit{C. Modality.}
\underafsnit{a) Pure Modality.}
\paragraf{17}{Quantitative Quality, Qualitative Quantity, Modification.}

The mutual relation of the qualitative and quantitative moments is developed in \emph{modality}. Every quale must as such qualify itself. Its qualification is a self-movement through all relative delimitations and determinations, and as a consequence of this a quantification: the quale delimited in itself is a \emph{quantitative quality}. The determinate quantum results from a quantifying. Quantification is a development of a determinate relation between discrete, i.e. qualitatively distinct, quanta: every actually determinate quantum must therefore be posited as \emph{qualitative quantity}. The transition of qualification and quantification into one another is a \emph{modification}, which introduces the determinate existence of quality and quantity.

The categories are positive expressions for hypostatized thoughts. Because thought receives its fixed shape in the category, it appears with a certain restriction and, as a consequence of this, with a determinate one-sidedness. --- The consciousness of this one-sidedness calls forth contradiction, and this in turn the dialectical movement; therefore dialectic always goes against the categories.
% --- p. 52 ---
\setcounter{footnote}{0}%
\opage{52}In the quality-categories thought everywhere seeks its foothold in the fixed, immovable limit: being is \textbf{not} nothing, something not other, the One not the Many. But dialectic always breaks down the posited barrier and expresses the dissolution of the limit in a special transition-category (becoming, change etc.). --- The transition-categories themselves are, with respect to the class of categories in question, partly \textbf{relative}, partly absolute. The relative transition-category merely shows the insufficiency of single subordinate categories and thereby introduces the thought-determination in higher categories of the same class; the absolute transition-category, on the other hand, presents the developed class of categories in its whole one-sidedness, and thereby leads over into a system of new categories\footnote{Thus becoming is a relative transition-category; for although it sublates the qualitative limit between being and nothing, it nevertheless introduces a new qualitative limit between the something and other of determinate existence. Only when the qualitative limit has reached its highest determinateness as exclusive limit, and exclusion in its absolute carrying-through has shown itself as assimilation, is every immediate quality-limit to be regarded as sublated.}.

As the last and highest quantity-limit is sublated, it is hereby brought to clear cognition that every posited limit is variable. Since now the variable limit as such is the criterion of quantity, it follows from this that what comes to determinate utterance in the quantity-categories is nothing other than what has already implicitly moved through the whole dialectic of quality. But as the qualitative determinations are volatilized in the quantity-categories, dialectic must again lift this one-sidedness and continually
% --- p. 53 ---
\setcounter{footnote}{0}%
\opage{53}restore the vanishing limit. It was these attempts at restoration that in the preceding section led to ever higher and fuller categories, and thereby at last to the cognition that the quantitative determinations had to be integrated by repeating the vanished quality. Thus: one begins with the \textbf{being}, immediately given qualities (qualia). Their immediacy is sublated as they go through their \textbf{becoming}. This becoming, in which the qualities vanish, at first takes on the character of a quantification. But since quantification is the qualities' own genesis, it must in its development show itself as the proper qualification and thus lead to qualitative results. It is this fundamental movement that determines the whole of \textbf{modality}.

Every true quality must qualify itself. As long as quality has not yet gone through its qualification, it represents itself as a being, a sensuously given, that belongs merely to representation. The whole sensuous world is in its immediacy a sum of such unqualified qualities. By reducing the world of sense to its simplest elements, one arrives at crass \textbf{atomism}\footnote{Even the Herbartian Realia are unqualified qualities.}. The atoms are sensuous immediacies; since they lack qualification, they necessarily also lack quantification, and are consequently \textbf{without extension}. Thus no atom has the moment of otherness in itself either; the atom presents only the abstract something whose other is external and indifferent: distinction follows upon distinction without any transition. The medium in which the atoms are found is therefore also an external medium (space, the void), and the power that connects and separates them an external, incomprehensible power. This, then, is the logical contradiction of the sensuous atoms,
% --- p. 54 ---
\setcounter{footnote}{0}%
\opage{54}that they are to have a being without becoming, a determinate being without inner change, a being-for-itself without mutual exclusion and assimilation.

Pure being represents itself first as given quality, as \textbf{being} being; nothing can likewise present itself as \textbf{being} nothing. But as being results from nothing, it thereby posits nothing as a limit that it itself sublates, while nothing, in the reverse direction, takes itself back through the sublated limit of being. Both \textbf{qualify} and extend themselves as they \textbf{become}, each in its own \textbf{manner}: the quality of modality is a \textbf{quantitative quality}.

Modal being is posited as the primal quality that by its qualification produces all being determinations. But in the mere process of qualification everything fixed becomes fluid and vanishing, and being becomes the intensive magnitude that extends itself \textbf{without measure or manner}. Abstract qualification is a mere quantification; it is seen here as a coming-to-be that can never come to be anything, a restless change that never leads to being-for-itself, a unity without plurality, and is consequently a one-sided abstraction. Since abstract quantification never brings it to fixed quantitative determinateness, it can never constitute determinate quanta either. The quality-moments, which come and go to infinity, are innumerable, and to the primal quality itself, which is absorbed into the pure process of quantification, number cannot be applied at all. Quantification must therefore be posited as a moment for qualification, and this in turn lead to determinate qualities. As beginning-moment the primal quality or general quality must be the most abstract; the further, on the other hand, qualification is continued, the more concrete % print: nden (for uden) (misprint)
% --- p. 55 ---
\setcounter{footnote}{0}%
\opage{55}become the qualities resulting from it. It is by these groups of more universal and more special qualities\footnote{We use ``quality'' here in the widest sense. The distinction between the qualities as independent existences and as non-independent moments can only be brought out later.} (genera and species) that quantity properly comes into its own and becomes a moment in determinate existence. For just as the determination of quantity perishes when all qualities vanish in the general quality, so would it also lose its meaning if the determination of quality were driven to its extreme point and speciality held fast by itself. Only by the more special qualities being reduced to a more universal quality does a \textbf{qualitative} number-unit come about, and, as the universal quality again points to its specifications, a \textbf{qualified} amount. By, e.g., reducing 4 tables and 3 chairs to the common quality ``furniture'', a qualitative unity is obtained with which numeration can begin. But if one now stays with the more universal determination, then furniture = furniture; only by bringing out the specification is the \textbf{one} piece of furniture seen in its distinction from the \textbf{other}, and one can now count (determine the \textbf{amount} of) the separate pieces of furniture. Modal \textbf{quantity} is always \textbf{qualitative}. The transition of qualification and quantification into one another (quod modum efficit) can now be determined more closely as \textbf{modification}.

In modality something and other are each to be posited as \textbf{qualified} qualities. These qualities now subsist with relative independence in a system of qualitative moments. The determination for the qualitative something is always this, that it must go
% --- p. 56 ---
\setcounter{footnote}{0}%
\opage{56}under in other; other is its \textbf{absolute} limit; where other begins, something must cease. But in order that something can go under, it must first arise. It does not arise at one stroke; quality must qualify itself; it has its being in its becoming. Therefore the two moments that in abstract becoming immediately went together into one are now also each separated out and specified in modality: modal becoming is modified by a determinate opposition between \textbf{arising} and \textbf{passing away}.

As something begins its qualification, it posits its first qualitative moment, a moment in which it can be said both to be and not to be. Qualification therefore advances successively, as it continually posits new quality-moments and preserves those already posited. But the quality that results from this is as such restricted; it is a something that is delimited on all sides by other. In order to form the full delimitation of something, other must therefore both go before and follow after. The preceding other is something by itself, the subsequent likewise; the movement is this, that just as something arises by a negation of the antecedent other, so the following other emerges in turn from the negated something. But this opposition, which thus enters into otherness itself, is here to be seen merely as a moment for the determinate existence of the intermediate something. During its arising something must continually change. The change consists in its excluding the presupposed other by assimilating its otherness and making itself its own other. In this modification of change something thus has its subsistence, its \textbf{abiding}\footnote{``Blive'' expresses both \textit{manere} and \textit{fieri}.}. Something abides by its
% --- p. 57 ---
\setcounter{footnote}{0}%
\opage{57}becoming (qualification) being posited as a moment for its being (quality). But in the assimilation of other, something has taken its own nothing up into itself. The moments of otherness came from other when this was to pass over to something, but now it is again the determination of something to pass over into other; the otherness-moments by whose assimilation it has reached its subsistence therefore develop into what they actually are, namely the negation of something itself. Something subsists in a nothing that in turn develops into positive being in the qualification of the subsequent other\footnote{The fundamental thought that something becomes its own other through something other than itself is of pervasive influence on the philosophical cognition of the world. Atomism, which acquiesces in the view that something has other outside itself, leads to materialism; for it belongs to spirit to become its own other to itself: the Orientals, who let the primal quality be its own other to itself without properly positing anything other than itself, are pantheists. Fichtean idealism does cognize that the I is to assimilate its otherness to itself through its other; but since this other is not qualified in itself, it remains standing as an immediate barrier. Christianity cognizes God in his personality; for here he is his own other (the Son) through an other qualified in itself (the world).}. The subsistence of something thus introduces its passing away; with this modification the change of something is an annihilation. But if something now passed away in the same instant as it arose, there would be neither coming-to-be nor annihilation: each of the modifications of change must have its extension. In order to sustain its modal oppositions, qualification modifies itself into quantification.
% --- p. 58 ---
\setcounter{footnote}{0}%
\opage{58}As something, during qualification, posits and sublates its relative limits in order to reach the absolute limit, it posits itself as a determinate \textbf{continuum}, and when the absolute limit is overstepped, other begins its continuation. While discreteness, then, during the whole coming-to-be is absorbed into the qualified continuity, discreteness on the other hand predominates in the limit itself between something and other, and thereby qualifies itself in discrete continua. The quality-unity of something is thus to be determined as an \textbf{intensive} magnitude that \textbf{extends} itself in the system of the qualitative moments. In this way quality and quantity both come each into its own. Quality requires its fixed, determinate limit; this it now obtains by transferring the variable onto quantity. But quantity varies only as long as the qualification in question lasts. By the restriction of quality, quantity obtains its fixed, immovable limit, indeed even becomes, as that which comprises all the qualitative moments of the particular quality, the only determinate expression for quality in its wholeness: where the absolute limit of quality lies can be decided only by quantity. What is constant in quality designates what is changeable in quantity; but the constant limit of quantity expresses the whole quality's own change. Thus modification leads to

\underafsnit{b) Measure.}
\paragraf{18}{The Immediate Measure, the Qualified Measure, the Modification of Measure.}

The qualified quality expresses its determinate existence by the determinately existent quantity, and is
% --- p. 59 ---
\setcounter{footnote}{0}%
\opage{59}thus determined in \emph{measure}. For quality in its mere determinate being, the quantitative determination is superficial, i.e. an \emph{immediate} measure. Where, on the other hand, the being of quality is posited in unity with its coming-to-be, quantification develops itself out of qualification itself, and measure is qualified. The \emph{qualified}, immanent measure of quality is determined by a relation to the immanent measures of other qualities; by this relation measure itself is infinitely \emph{modified}. Qualification must lead to quality, as becoming to determinate being, and immediacy with the determinate limit is restored. Thereby the determinate quantum now also results from quantification. Quality has thus become quantum; they have their determinate existence in and with one another. But in the same way as something has quantified itself and in its determinate being shows itself as magnitude, other has also obtained its quantification and \textbf{is} likewise a magnitude. Both are then quantitative qualities and in this respect completely equal to one another. But when all opposition between something and other is sublated, determinate being itself vanishes in pure being. It is now a matter, then, of finding a new and higher distinction between something and other, that is to say, of determining the limit between the qualitative magnitudes mutually.

The \textbf{qualitative} limit showed itself (§ 14, p. 26) as the exclusive negation; but since the bordering qualities were immediate and as a consequence of this without extension, they themselves were reduced to mere limit-moments. Quality was therefore absorbed into the self-movement of the limit, developed itself as \textbf{quantitative} limit
% --- p. 60 ---
\setcounter{footnote}{0}%
\opage{60}(p. 27), and thereby produced an extension that sublated all delimitation.

The \textbf{modal} limit, on the other hand, is fixed and determinate like the qualitative, for it separates \textbf{something} from \textbf{other}; but something and other are \textbf{magnitudes}; the modal limit is therefore movable like the quantitative, for it gathers in its determinateness the relative limit-moments of the magnitudes. The modal limit thus has two sides, whose mutual relation is developed in \textbf{measure}.

Since both main sides of measure are each by itself modal, they both contain the immediate unity of the constant limit of quality and the variable limit of quantity, but in opposite ways. For while the one side contains the movable limit as a moment sublated in the constant one, and is thereby posited as a quality with \textbf{determinate} extension, the other, on the contrary, lets every constant limit-moment vanish in the varying limit, and shows itself as a continuum with \textbf{indeterminate} extension. Extension is then the characteristic common to both sides; this is precisely what is superficial in the immediate \textbf{measure}, that all quality-determinations are here absorbed into mere extension.

But quality continually takes itself back through the quantitative moments. The specific development of determination arises from the constant and the variable limit being united in each of the sides: thus they can be \textbf{measured} against one another. In measuring, the determinate continuum is posited as \textbf{measuring-rod}. The measuring-rod is a quantity that represents itself in representation; \textbf{how} large it is must be immediately intuited: its extension, i.e. its specific quantity, is the proper determinateness of the quality. The constant
% --- p. 61 ---
\setcounter{footnote}{0}%
\opage{61}limit that is posited in and with the continuum of the measuring-rod now becomes the means whereby the indeterminate continuum, i.e. that which is measured, is reduced to determinate limits. The measured magnitude thus develops both its moments: in and for itself it has a variable limit that not only makes possible but even requires constant limit-determinations; this last demand is satisfied by the application of the measuring-rod. But the determinate continuum of the measuring-rod, which has the variable limit in itself, now also requires that the hidden moment of changeableness be developed and determined. With its immediate determinateness the measuring-rod can be extended and restricted at will; but in relation to the measured magnitude what is arbitrary in its extension becomes delimited. The measuring-rod is a given unity that finds its amount in the measured magnitude.

Since the immediate measure merely expresses the superficial relation of the quantities, it follows from this that every deeper distinction must be excluded, and the quality-peculiarity that qualifies the magnitudes to be measured against one another must be immediately given\footnote{Therefore line is immediately measured with line, surface with surface.}.

But modal determinate being develops itself as an abiding (§ 17), in which the becoming of qualification continually alternates with the being of quality. Since quality \textbf{is} as it \textbf{becomes}, it must be quantitatively determined as an \textbf{intensive} magnitude that extends itself in a series of intensive moments (§ 16. p. 44). But \textbf{something}, the determinate quality, arises and abides only by assimilation and exclusion of \textbf{other} (§ 17. p. 56). The quality coming to be can therefore only be determined
% --- p. 62 ---
\setcounter{footnote}{0}%
\opage{62}i.e. \textbf{measured} by another intensive quality, which is likewise qualified by its extension.

As now the one of these qualities is posited as indeterminate magnitude, whose involuntary increase and diminution is merely indicated by an immediate measuring-rod, the other, on the contrary, will maintain its qualitative peculiarity by assimilating to itself the variable magnitude in a peculiar way in different quantities. (Thus, e.g., water takes up a different quantity of heat).

The assimilating quality has a delimited being-for-itself, and, as a consequence of this, likewise only a delimited assimilation of its variable quantum, which here plays the role of otherness. (Water as water takes up only a \textbf{certain} quantity of heat). It is now this assimilation that in its whole relativity and delimitation receives its development in the \textbf{qualified measure}.

Quality as such must be given with a determinateness that is, consequently also with a certain assimilation of its other, i.e. of the variable magnitude. (Water at the very thawing-point has a \textbf{determinate} heat). Since now a qualifying becoming must go before every qualitative being, one can only seize and hold fast the constant quality in positive unity with the variable magnitude by a leap that belongs to immediacy as such. (At the zero point of the thermometer one can, e.g., have both water and ice). But since the quality that is must \textbf{become} just as much as \textbf{be}, it reduces its relative being to a vanishing \textbf{state}, and documents its qualitative movement by letting the variable states succeed one another.

Everything thus depends on the dialectical relation between the constant quality and its changeable
% --- p. 63 ---
\setcounter{footnote}{0}%
\opage{63}state. In pure immediacy quality is held fast by itself and the varying states by themselves. Therefore one cannot comprehend either how quality preserves unchanged likeness with itself although it continually takes up different quanta of the variable magnitude; still less does one see why quality, after it has shown itself indifferent toward a series of varying states, is at last altered at one stroke by a mere increase of quantity.

These contradictions are lifted by consideration of the qualification of measure itself: a) Since the constant quality as such has a certain quantum of the variable magnitude in itself, it lies herein that its quantitative variety expresses precisely the quality's own determinateness, and that this, as a consequence, cannot be indifferent toward any state whatsoever. It is only through the \textbf{changes} that quality can develop its \textbf{constant} character; these therefore belong to its own qualification. When it is therefore said that the same quality sustains itself in all alternating states, the truth is this, that quality first expresses itself completely in all the states that are posited and sublated.

b) Since now, accordingly, no single state satisfies quality, it would never obtain its adequate expression if the moments of a lower state could not be preserved in a higher. But it is here that gradation develops in its full meaning. The external, immediately determined quantity of the changeable magnitude is, in every given state, assimilated into peculiar unity with quality. Quality therefore has magnitude in itself as determinate intensive magnitude, that is, by an even distribution of the intensity-moments that are to characterize the single states, as degree. % print: ædæqvate (for adæqvate) (misprint)
% --- p. 64 ---
\setcounter{footnote}{0}%
\opage{64}Degree is here posited in and by a quality; therefore there is here a \textbf{lowest} and a \textbf{highest} degree. In the lowest degree the quantity of the variable magnitude is sublated as a moment for the state in which quality begins its determinate being; in the highest, on the other hand, quality appears absolutely identified with the state itself. (At the 100th degree (C) water has assimilated to itself all the heat that, according to its peculiar quality, it is able to take up). c) As long, now, as the relative intensity does not express quality's completed act of assimilation, so long can quality overstep every quantity-limit and change its state without itself being changed; when, on the other hand, assimilation has reached its limit, and the adequate state sets in in which the quality developed through the quantitative changes acquiesces, then it lies in the nature of the matter that this state cannot be changed by a raised degree without quality being altered with it. The surprising change of quality is thus introduced by the change of quantity. Hereby the constant quality is grasped as the immanent measure of the changeable magnitude.

In the measure thus qualified we then have: α) a quantity-unity positively determined by quality (from 0° to 100° (C.) an equal quantum of heat is always contained, if other circumstances do not act on the water); β) an amount determined by this unity, which is posited and immediately observed as the variable magnitude in the quality goes through its gradations; and finally γ) a true identity of unity and amount (the heat of water at the 100th degree contains 100 degrees of heat).
% --- p. 65 ---
\setcounter{footnote}{0}%
\opage{65}The immanent measure develops itself through innumerable gradations (between freezing point and thawing point there lie not merely 80 or 100 degrees, but, according to infinite divisibility, cf. p. 30, an arbitrary quantity); as a consequence of this every number-determination is sublated. But where number has vanished, the specific limit of measure is erased; if the gradations of the qualified measure are to be determined, then the marks of the degrees must be set off arbitrarily, and gradation be immediately observed by them.

But when the qualified measure thus depends on the immediate measuring-rod, it is thereby absorbed into unqualified determinations and itself becomes immediate. This immediacy has a twofold side: partly, namely, it lies in the arbitrary division of the measuring-rod, partly in the constant quality's own limits, as this by its alteration immediately indicates when the variable magnitude in it has overstepped its measure.

The variable magnitude itself, on the other hand, is not altered (heat continues to be heat after the water has become steam), but seeks within the limits of another constant quality another immanent measure etc. Hereby the qualified measure develops itself in a series of immanent measures, and thus becomes external to itself. From this externality there now arise the specifications of measure itself. Since the variable magnitude is determined by the totality of the constant qualities in which it finds its immanent measures, its quantum is hereby specified; but the constant qualities too are quantitatively specified, for their qualitative peculiarity is determined solely by the particular quantity that each of them takes up of the variable magnitude. If these quantities are now compared with one another mutually, then measure is reduced to constant
% --- p. 66 ---
\setcounter{footnote}{0}%
\opage{66}ratios. The varying immediacy of quantity is hereby sublated, and the immediate measuring-rod, whose arbitrary division now no longer has any influence on the determination of the constant ratios themselves, is reduced to an external means whereby the immanent measure measures itself. (For measuring the temperature it is indifferent whether the thermometer is divided into 80° or into 100°).

Reduced to the constant ratio, the immanent measure has its specific expression in the exponent of the ratio. Every immediate quality is as such one-sided; it seeks its negation and determination in an opposed quality, as this one does in that; both are then to be regarded as each other's sides, and the complete quality as the unity of both. The developed quality is always a double-quality, which in unity with itself exposes the ratio between two qualitative sides. But as the qualitative sides determine one another mutually and negate one another, qualitative immediacy vanishes in their mutual relation; the one is thus taken up into the other, the one is contained in the other, the increase of the one entails, if the ratio is to subsist, an increase of the other: their mutual determination is altogether quantitative. Since the ratio, as the exponent shows, can remain the same provided only that the increase or diminution is undertaken in a corresponding manner with both sides (1:2 = 2:4), quantity has obtained in the exponent a fixed, specific character, mediated through the quantitatively variable sides. This specifically determined quantity expresses precisely the explicated quality, as the equilibrium of the quantitative sides expresses their qualitative determinateness. Thus the complete development of the immanent measure is to be sought in the exponent,
% --- p. 67 ---
\setcounter{footnote}{0}%
\opage{67}But in the system of immanent measures developed above this exposition is not yet carried through. It is indeed evident that the specific quality everywhere conditions the specification of quantity and conversely; but these specifications are not yet deduced from one another. The variable magnitude, which can be regarded as the common side of the constant qualities, as the general quality running through them all, is taken as it is immediately found. Since no increase or diminution ever brings it to give up the peculiarity of its quality, the whole alteration that it produces in the constant sides in which it finds its measure is also a mere modification. The constant qualities therefore retain the whole stamp of immediacy, from which it then also follows that the exponents likewise show themselves as immediate number-determinations.

If specific qualification is to be absorbed into specific quantification, the qualities must forfeit their immediacy in quantity, in order to realize the true quality in the quantitative ratio. Every immediate quale must then show itself as it is, namely as the neutrum that is neither itself nor its other, but waits for another neutrum that will make it its other, while it itself makes itself an other for this one. This qualification is represented in chemical neutralization, where the exponent can be said to be the developed quality itself.

Let A and B be the qualitative sides (substances) from whose union the qualitative exponent (the body) E results; their immediate qualities are then no longer to be sought in E, for in that case E would merely be a composition of A and B, but no quality by itself; but as the immediate qualities are negated and sublated in quantity, their
% --- p. 68 ---
\setcounter{footnote}{0}%
\opage{68}quantity itself, on the other hand, is preserved: in quantitative respect is
E = A + B.

Here the mutual specification of quality and quantity now shows itself most clearly of all. Which quality-determinations lie in A and B can be known only from the quantitative ratio under which they combine with one another; but on the other hand one cannot determine this ratio either, except precisely by testing the qualitative peculiarity of both and asking them themselves under which quantitative conditions they will unite. From this it is explained: α) that A and B must always unite in a certain ratio in order to produce a certain qualitative exponent; β) that from the nature of such a ratio one cannot immediately conclude that it is the only one; for since in its exponent it expresses only one determinate quality, both sides may, under a new ratio-exponent, present new quality-moments; finally A and B may each by itself be sides for several other immediate qualities: A can, e.g., unite with C, D etc.

The relation between the sides themselves and their exponents is now this, that the former are absorbed as moments into quantitative ratios, while the latter, on the contrary, represent the qualitative being-for-itself. In the series: E (A + B), E′ (A + C), E″ (B + C), E‴ (C + D) A can be determined altogether quantitatively; it thereby shows itself as the variable magnitude that finds its immanent measures in the constant qualities B and C. If, on the other hand, A and B are regarded as constant qualities, then C represents the variable magnitude; the same holds of B, which as common quality for A and C likewise takes on a quantitative character and is taken up by these in different quantities. The exponents E, E′, E″, on the other hand,
% --- p. 69 ---
\setcounter{footnote}{0}%
\opage{69}exclude one another as for-itself qualities, and thereby express the specific ratio of the quantitative sides.

Yet even this opposition between the exponents and their sides must, when it is further developed, sublate itself. The exponents, which by negation of the quality of the sides have gone through their qualification, themselves take on the stamp of immediacy, and can thus again become sides for exponents of higher order. (E + E‴ = E⁗), and since every immediate quality as such presupposes a qualification (p. 53), the subordinate sides are themselves to be regarded as exponents Insofar now as all exponents are reduced to sides, all quality-determinations are sublated in quantity; insofar, on the other hand, as all sides are posited as exponents, qualitative being-for-itself takes itself back at all points.

The tension that hereby sets in between quality and quantity is lifted in \textbf{affinity}.

According to affinity the qualities cannot unite indiscriminately, but must be disposed toward one another; their mutual assimilation entails an exclusion, whereby the quality-moment as such asserts itself. The greater or lesser disposition of the qualities toward one another does not here stand primarily in relation to the quantum of the sides; rather, the union itself is a neutralization that grants quality its right. Since all qualities are to be regarded as results of the qualifications of the primal quality (p. 54), they are all mutually related; but since specification necessarily belongs to the forming of the exclusive limit-determinations of quality, their kinship is not equally great. Affinity \textbf{has its degrees}.

As a moment in the system of qualities no qualitative side can be considered in isolation. Every new union
% --- p. 70 ---
\setcounter{footnote}{0}%
\opage{70}is a dissolution of previously contracted unions. By the neutralization of A and B an immanent measure is constituted. But union with B presupposes its separation from C; with C, A also forms an immanent measure; this must then be dissolved in order that the former can be realized. The \textbf{elective combinations}\footnote{That the categories of affinity and elective combination hold not only in chemistry but also in psychology, Goethe has shown in ``die Wahlverwandtschaften''.} thus present the most complete evolution of the qualified measure. The immanent measures, whose totality expresses the whole self-development of measure, are not here compared with one another in a merely external way; rather, the one presupposes the other and proceeds from the other; the coming-to-be of the one is the annihilation of the other: the immanent measures thereby modify and negate one another mutually, and measure is sublated.

\underafsnit{c) Absolute Modality.}
\paragraf{19}{Indifference, Difference, Pregnant Being.}

The highest distinction and unity of unity and distinction to which being as such can bring it is posited with the immanent measure, whose self-development closes itself off in absolute modality. As the total unity of the distinct moments, the immanent measure preserves all distinctions in its unity, so that the separation of unity and distinction is quite indifferent: modality is determined as the indifference of being. But the moments of measure are themselves immanent
% --- p. 71 ---
\setcounter{footnote}{0}%
\opage{71}measures, each of which concentrates all being-determinations in itself; in this way the union of unity and distinction is indifferent, and being is absorbed into infinite \emph{difference}. Since the whole being-unity is in every distinction, distinction and unity are everywhere determined in the selfsame way; herewith all being-determinations vanish in \emph{pregnant} being.
It is the task of philosophy to cognize the true unity in the true distinction. The unity that has distinction outside itself is itself distinct from its distinction and thereby reduced to a distinction-moment. In the same way the distinction that has unity outside itself also goes into one with itself, and thereby constitutes a side of unity. The true unity must unite both sides of unity in itself, i.e. be posited as a unity of distinction and unity. But the unity that negates and excludes distinction is, as it appears as a distinction-moment, precisely thereby distinct from itself: it is, in its separation from distinction, in itself distinction. Likewise the abstract distinction that negates the immediate unity, in order to be absorbed as a unity-moment into the developed unity, is distinct from itself.
The separation of unity and distinction is thus no longer immediate, but results from the self-development of both. While the higher unity, through the immediate and thereby subordinate distinction-moments, goes together with itself and predominates, distinction takes itself back from the developed unity-moments as a higher distinction co-ordinated with the concrete unity.
But when the developed distinction is co-ordinated with
% --- p. 72 ---
\setcounter{footnote}{0}%
\opage{72}the developed unity, the latter can no longer predominate; it is then reduced to a distinction-moment for a higher distinction between unity and distinction.

This ascending development forms a dialectical scale, which marks off its determinate steps in categorical immediacy. As long, then, as a logical sphere has specific categories for a higher unity, it is not yet dialectically exhausted; when, on the other hand, the category is developed that expresses all the specific determinations of the sphere in highest unity, the determinations themselves also appear in their highest distinction; if it then turns out that unity, when it is held fast, loses itself in distinction, and that distinction, when it is sustained, vanishes in unity, then there is neither distinction nor unity; the categories are then insufficient, and the sphere in question is sublated.

The unity and distinction of being, with the specific transitions of both into one another, are considered from the standpoint of quality, of quantity and of modality.

The immediate qualitative unity in its most abstract shape is \textbf{pure being} itself; but as it reduces all distinctions to non-being and makes them \textbf{nothing}, it is itself reduced to an immediate distinction-moment, namely a being distinct from nothing. It is this immediate distinction that in \textbf{becoming} yearns for a higher unity, since each of its sides shows itself distinct from itself. \textbf{Determinate being} is now the more concrete unity that has overcome the empty nothing, and cannot be reduced to a distinction-moment by pure negation; it is no immediate combination of being and nothing, but a union of being and being mediated through negation. But as being in determinate being is to be determined by being, it must separate itself from itself, and distinction enters anew in
% --- p. 73 ---
\setcounter{footnote}{0}%
\opage{73}\textbf{something} and \textbf{other}. This distinction is higher than the distinction between being and non-being; for both distinction-moments have, each by itself, a determinate being. Hereby determinate being itself has become a distinction-moment, and unity is broken. For the distinction lies not merely in a separation of something and other from one another, but is at the same time a distinction between determinate being as such, regarded from the standpoint of unity, and determinate being in something and other, determined by its own distinction-moments, thus: a distinction between unity and distinction. Insofar as determinate being is given up on all sides to distinction, it exhausts itself in \textbf{change}.

As the change of change there results the predominating unity of quality, namely being-for-itself, which lets something and other go into one. But as the One negates all other outside itself, it thereby comes outside itself; the One is broken asunder into the Many, and distinction again displaces unity in the reciprocal exclusion of the ones. Since, however, reciprocal exclusion develops through absolute exclusion into an assimilation, in which the separated Many flow together again into one, distinction cannot sustain its preponderance either. A restful equilibrium that is, between unity and distinction, quality can never bring about; for if unity is held fast, everything goes into one, and distinction vanishes; if, on the other hand, one vindicates distinction, then in the Many it is all over with unity. Here, then, there is only an alternation of unity and distinction, a restless yearning for equilibrium; but equilibrium itself must be sought in \textbf{quantity}.

The category that stands on the limit of quality and quantity, of distinction and unity, is the limit itself. In the qualitative limit unity is negative, distinction positive; here, then, there is specific distinction between unity and distinction: the quantitative limit,
% --- p. 74 ---
\setcounter{footnote}{0}%
\opage{74}on the other hand, which, as it makes the bordering into limit and the limit into the bordering, and in the movement of continuity has the resting points of discreteness, posits unity (continuity) just as positively as distinction (discreteness); in the quantitative limit there is therefore unity of unity and distinction. Limitation, which shows the omnipresent limit, now contains the task of everywhere finding and positing the distinction of unity (every continuum is as such discrete), the distinction of distinction (every discretum is as such continuum and thereby discreteness in itself) and the unity of distinction. This distinction and unity of unity and distinction is given in \textbf{number}. The one of the number-unit differs from the qualitative One precisely in this, that the latter negates and excludes plurality, while the former, on the contrary, vindicates its unity by the discreteness of plurality. The plurality discrete in itself preserves its distinction from unity by the units of \textbf{amount} being specifically ordered; since every numbered unit is assigned its place, quantitative plurality cannot, like the qualitative, be dissolved again into a single standing One. But the distinction of discreteness does not interrupt the unity of continuity either, for each of the ones of amount is itself a unity for a new amount. From \textbf{innumerability}, the infinity of the number-unit dissolved into an infinite amount, there develops the predominating unity for which the distinction of number-discreteness is posited as a sublated moment. This elastic unity, pregnant with all specific distinction-moments, is \textbf{intensity}, which in \textbf{degree} posits the whole developed distinction by positing itself. Unity is here its own distinction, and thereby comes to predominate, but each distinction-moment has the whole unity in itself, and is thus itself its unity; this then points to the preponderance of distinction. The predominating
% --- p. 75 ---
\setcounter{footnote}{0}%
\opage{75}distinction is now no longer a distinction between quantity and quantity, for every such distinction is already exhausted, but rather an exclusive distinction between quantity and non-quantity, in which the qualitative is taken up again. That quantity is not by this exclusion reduced to an immediate qualitative moment, whereby dialectic would sink back onto the overcome standpoint of quality, is prevented in \textbf{approximation}, whose true development is an \textbf{evolution}. According to approximation the exclusive quality-moment is introduced by the pliant, assimilating quantity itself, and according to evolution the whole quantitative gradation is itself an expression for the specific positions of quality.

The distinction between the qualitative and the quantitative moment that hereby sets in lies solely in the \textbf{manner} in which being is determined in each of them. \textbf{Modality} yields the higher unity. For quantity does not come in from outside as an arbitrary auxiliary determination for quality, but is quality itself, determined in another manner. Likewise it follows from the dialectic of quantity itself that it must determine itself in a manner that is precisely opposed to the quantity-determinations as such. Since now the qualitative manner of determination is thus opposed to the quantitative, and ``everything depends on the manner'', the all-sided distinction is realized as a distinction between two classes of being-determinations. Hereby the unity of modality is dissolved into two modal sides, \textbf{quantitative quality} and \textbf{qualitative quantity}, each of which by itself is the whole modality. But as each of the sides separates itself from itself, the modality of quality is itself that of quantity, and unity is striven for in \textbf{modification}. Modification expresses precisely the fleeting distinction in the preponderant unity. Quality
% --- p. 76 ---
\setcounter{footnote}{0}%
\opage{76}and quantity are both sides of the \textbf{same} being: the specific quantity-categories must in their full extent be integrated by corresponding categories of quantity. The indissoluble and specific unity of both classes of categories is posited at one stroke in \textbf{measure}. But the unity of measure too determines itself to distinction: both sides in the immanent measure are, each by itself, themselves immanent measures. In these immediate and infinite series of measures the one measure is separated from the other by each of the sides, when it is regarded as measure, being posited as unity for sublated distinction-moments, while conversely every measure, insofar as it is itself a side, yields a distinction-moment for a still higher unity. In this way the immanent measure everywhere comes outside itself. a) The highest immanent measure would be that in which all others were sublated as negated sides; but it lies in the dialectic of every measure itself to be posited as side. b) The highest measure should, as that in which all other measures are posited, be the most concrete of all; but every measure is, as the negation of the sides, a pure being, which in its sheer immediacy presents only itself: every immanent measure is therefore, as such, equally abstract and equally concrete. The distinction of the absolute measure from the relative ones therefore seems to have to be reduced to the modal genesis whereby the one measure becomes from the other. c) The highest measure would then be that from which all the others would have to be deduced. But the highest measure is, as negation and sublation of all lower ones, precisely itself deduced from these.

The absolute measure is hereby determined as the modal intensity that, under a gradual approximation, seeks its adequate expression in the relative measures without, however,
% --- p. 77 ---
\setcounter{footnote}{0}%
\opage{77}ever finding itself in them: \textbf{the highest measure is} \textbf{unattainable}.

As that whose determinations can never be exhausted in any finite, attainable measure whatsoever, the unattainable measure separates itself from all relative ones, and is thereby itself posited as a distinction-moment. But separated from relativity, determined solely by its negation, it is in its being-for-itself itself relative. Its inflexible determination is to specify its content in the relative measures, just as these too, each by itself, always express the absolute in a certain degree.

The becoming in which the absolute measure continually passes over into the relative, and this into that, is \textbf{measure}'s own infinite \textbf{modification}; by this it strives through all its distinction-moments toward its highest unity. Since every special and relative manifestation of the absolute measure is only a modification of the same, and the distinction of modification only a fleeting distinction, the infinite series of modification must end in \textbf{absolute modality}.

Measure in its highest unity gives itself its own sides; indeed, measure is itself its own side, and the sides, each by itself, measure itself. In absolute modality every distinction has thus vanished. All the distinction-moments that were held back by pure being, repressed in non-being, have now gradually come forth from their nothing as passing being-determinations, but have thereby at the same time shown themselves in their true nothingness, forfeited their meaning in the dialectical process, and thus betrayed their indifference. The unity that now presents all distinctions in this their mutual indifference, and by its indifference toward them all shows itself indifferent toward itself, can only be determined as \textbf{absolute indifference}.
% --- p. 78 ---
\setcounter{footnote}{0}%
\opage{78}Every being-unity is always posited as the indifference of the sublated distinction-moments; but such an indifference is still relative as long as it can itself become a distinction-moment. Thus in the moment of pure being, being is posited as the indifference of the distinctions, but distinction was latent in nothing; determinate being therefore became the indifference of being and nothing, and being-for-itself in turn sublated something and other in the indifference of the One. The whole dialectic of being can be regarded as a successive reducing to higher and higher indifferences. Absolute indifference is the \textbf{all} of being; for all being-determinations culminate in it; but the being of the determinations in indifference is precisely their non-being; indifference is the negative unity in which all distinctions sublate one another; as such, indifference is the \textbf{zero} of all-being. The zero is the standing-still nothing\footnote{See Kant's Versuch den Begriff der negativen Grössen in die Weltweisheit einzuführen. Werke (G. Hartenstein) Erster Band p. 21.}, which by the indifferentiating of distinction is torn out of the stream of dialectic. Since absolute indifference is indifferent toward all distinction, it is consequently also indifferent toward the distinction that there is between the highest unity and its sublated distinction-moments, and thus indifferent toward itself: \textbf{indifference is the logical} \textbf{readiness of being for all}.

The immanent measure is the neutrality of the sides; none of the sides (momentorum neutrum) is preserved as such in the neutralization; the being of measure is the non-being of the sides. As indifference, as the immanent measure, now itself constitutes its own sides, it thereby sets itself aside: \textbf{the being of indifference} is eo ipso \textbf{its non-being}. In order rightly to show its indifference toward itself, the indifferent unity determines itself to distinction, and being opens itself up in \textbf{infinite difference}.
% --- p. 79 ---
\setcounter{footnote}{0}%
\opage{79}Herewith it is then given that being concentrates all its determinations in any single determination whatsoever. The whole being is itself pure being, the whole being is a nothing, a something, an other, quantity and measure; \textbf{the whole being is} \textbf{all and zero}.

Every difference of being is consequently itself the whole being-unity. With absolute difference the unity of distinction and unity is predominant insofar as every difference becomes its own unity. But as each single difference in its unity with itself is an indifference of all the other differences, and reduces their distinction-moments to the nullity of indifference, it must itself vanish as an indifferent moment in every other difference, and thus differs from itself. The indifference of difference consists precisely in this, that all differences are equally valid. While being, then, in absolute differentiating, places the unity of its totality in every difference, it at the same time reduces every difference to indifference, posits all distinction-moments as vanishing, the unity of indifference, on the other hand, as the abiding; there then sets in once more a distinction between (the indifferent) distinction and (the abiding) unity. But since absolute indifference is indifferent toward itself, and is thereby absorbed into infinite difference, the distinction between unity and distinction is just as indifferent as the unity of distinction and unity.

The result of this is that being cannot be without giving itself a determination, and that the being-determinations continually negate one another. Unity displaces distinction, and distinction absorbs unity: the one determination cannot be at rest for the other. Even the unrest in which all-being hereby finds itself cannot even be expressed in any becoming whatsoever; for becoming
% --- p. 80 ---
\setcounter{footnote}{0}%
\opage{80}presupposes a long-vanished distinction between being and non-being. Being is as it is not, and is not as it is: \textbf{all-being is pregnant}.

According to this pregnancy every being-determination always means at the same time the opposite of what it states. Herein lies the highest criterion for dialectic's having exhausted and used up all the determinations of being. Absolute being will, in the successive negation of all distinction-moments, go together into unity with itself; but by the very negation of all being-moments it is itself infinitely negated and reduced to non-being. But non-being \textbf{is}, and is even all-being, since it is all-being itself that has posited itself in the moment of non-being. In its affirmation, then, being is negated, and in its negation at the same time affirmed. Every specific determination that is now to express this being affirmed in its negation can then no longer be held fast from the side of being alone, for thereby it would again be negated and pass over into non-being; just as little can it be posited as the pure negation of being, for every non-being is sublated again in being. Every true determination must therefore be posited at once in the moment of being and of non-being, and categorically express the pregnancy of negation and affirmation. But in the sphere of being this pregnancy is an indifference different in itself and a difference indifferent in itself; here, then, are not the higher determinations themselves, but only the dialectical birth-pangs under which the tried being is reborn in new categories.
% --- p. 81 ---
\setcounter{footnote}{0}%
\opage{81}
\capitel{Chapter Two.}{Essence.}
\underafsnit{Introduction.}
\paragraf{20}{The Essential, the Unessential, Relativity.}

Pregnant being is being and non-being in the unity and distinction of both. In positing its non-being as a moment for an affirmative, determinate being, and showing itself as \emph{the essential}, it must at the same time posit being as a moment vanished in non-being, and is thereby determined as \emph{the unessential}. Thus a distinction that simply is gets maintained in pregnant being (the standpoint of criticism). But all being is absorbed in becoming. Insofar as non-being is itself being, the unessential can continually be raised from its immediacy to the essential (the standpoint of ingenuity); insofar as all being, on the other hand, is itself non-being, the essential constantly perishes in the unessential (the standpoint of skepticism). The truth of the essential and the unessential lies in the infinite \emph{relativity} of both, which sublates the fluctuation of the
% --- p. 82 ---
\setcounter{footnote}{0}%
\opage{82}standpoints, deduces the one central movement from the other, and thus reduces the dialectical development to essentially distinct spheres.

Pregnant being is a double being in a unity of being. This being, doubled within itself, is however, by virtue of the characteristic pregnancy of the unity, no longer to be regarded as a being limited by non-being, for such a qualitative limit would sublate the being of both in each other. We therefore do not have here merely a being which, with non-being as its limit, becomes ``something'', while non-being, limited by being, becomes an ``other'', but a being in which being is, in \textbf{unity} with a being in which non-being is, and in which being therefore is not. It is precisely this being-in-itself of the reduplicated being itself that distinguishes the self-doubling of pregnant being from determinate being.

Every category of being expresses as such only pure abstract immediacy, an immediacy that is in turn sublated by the individual negation. In the categories ``something'', ``the One'', ``the immanent measure'', being is indeed presented with richer and richer determinations, but in each of these it is nevertheless always found over against another being, in which it must in turn vanish. What is characteristic in the determination can therefore everywhere be reduced to: \textbf{an immediate immediacy}, \textbf{limited by an immediate} \textbf{negation}. % print: overoveralt (syllable doubled) (misprint)

The dialectical progress in being takes place in that the corresponding nothing (non-being) becomes more concrete as the content of being is more and more developed, but the mutual relation of both is constantly dissolved in an alternation that continually prevents being from coming to itself. If being in a category is held fast with positive immediacy (as ``something'', ``the One'', ``measure''), then
% --- p. 83 ---
\setcounter{footnote}{0}%
\opage{83}precisely the being by which it was to integrate its determination is posited as immediate non-being; there is thus in being itself no non-being, and in non-being no being at all. Since there is no categorical preservation of immediacy in negation or of negation in immediacy, immediacy can as little as negation obtain any true being-in-itself. Pregnant being, by contrast, which posits a categorical unity of being and non-being, seeks its higher specific determinations in a \textbf{pregnant} immediacy \textbf{with a pregnant negation}.

If the pregnancy of immediacy and of negation is developed, then the opposition between the essential and the \textbf{unessential} sets in.

Being in its abstract immediacy is determined as pure, distinctionless being, i.e. as being without all negation. But this immediate immediacy has nonetheless been won through a negation, for only by negating the distinctions does one arrive at pure being. If this is now held fast in immediate equality with itself, it gets nothing as its limit, is thereby negated and thus comes \textbf{outside} itself. If it is to remain in itself and affirm itself as constant being, then it must indeed open itself to negation and concede to the distinction-moments their being, but instead of letting this being of the distinctions determine itself as its own non-being, it must conversely posit their non-being as its being; only thus does it negate the negation, reduplicate its immediacy and come forth as a being in being, i.e. as an \textbf{essential} being.

But since non-being has had to give up its positive being-moment to essential being, it itself retains only a negative being. It must indeed be acknowledged that that which is non-being also \textbf{is}, but seen in its
% --- p. 84 ---
\setcounter{footnote}{0}%
\opage{84}own negation it can no longer assume the character of pure immediacy and appear as the positive other over against the essential. The being that non-being has is precisely a being negated in itself. That which is non-being must thus pass for what it is, namely for the side of being that is abandoned by affirmative being itself. Since its non-being is its being, this being is itself a non-being. We thus get non-being in its own non-being, i.e. a pregnant negation. That which is non-being, with non-being as its being, is thus determined as the \textbf{unessential}.

The distinction between the essential and the unessential is hereby developed as a distinction of dignity, but the dignity itself depends on the relation of both. The essential is essential only with \textbf{regard} to the unessential, just as conversely the unessential is unessential only with \textbf{regard} to the essential. Each of the opposed terms doubles itself in the further development of the relation.

The immediate being of the essential is, as such, unessential, and assumes the character of essentialness only through a privation of the being that would otherwise belong to the unessential. The essential is thus essential only in the unessential, and so has the distinction of dignity in itself.

The unessential does not, as such, belong to itself; isolated from the essential it is only an immediate non-being, and consequently neither unessential nor essential. Herein lies: a) that the negated being of the unessential rests in a positive other, with regard to which it is determined, that it consequently cannot itself make itself what it is, but must fetch its unessentialness from outside; b) that the essential, which gives the unessential its unessentialness, thereby characterizes its
% --- p. 85 ---
\setcounter{footnote}{0}%
\opage{85}own essentialness solely in this, that this essentialness, unessential in itself, is in consequence nothing other than the very unessentialness in the unessential; thus: unessentialness is the unessential's own essentialness.

The relation must then be more closely determined thus: that the selfsame thing which in the one term is essential is in the other unessential, and since each of the terms refers itself to itself as to its negation, the selfsame term is in one regard essential, in another regard unessential. The distinction of dignity is entirely dependent on the different regard (diversus respectus). But no regard can take place without there being something to which one looks. This something can be sought neither in the essential nor in the unessential, since it is precisely this distinction of dignity that is to be determined by the regard itself. That to which one looks must therefore present itself as something immediately given, as an abstract essentialness valid for both terms. If, e.g., A is such a given essentialness, with regard to which a is posited as the essential, b on the other hand as the unessential, then A, as that on which the distinction of dignity between the subordinate terms rests, is the only thing that in this relation can be called the truly essential. But if A is the essential, then, considered in and for itself, it is at the same time both essential and unessential; its hypothetical dignity then depends on an opposed term B, which likewise can in itself be both essential and unessential. The distinction of dignity between A and B must therefore be determined by a still higher essentialness, etc.

\textbf{Insofar} as B, then, is the unessential, A is the essential, and the distinction of dignity deduced from it between a and b stands firm; \textbf{insofar} as B, on the other hand, is posited as the essential, A becomes the unessential,
% --- p. 86 ---
\setcounter{footnote}{0}%
\opage{86}and the subordinate dignity-moments are altered; insofar, finally, as both A and B, according to higher regards, are each in its own way recognized as essential, their subordinate moments become now more, now less essential, and the modification goes on to infinity.

The characteristic one-sidedness of \textbf{criticism} becomes evident from this. The critical method presupposes an essential distinction between the essential and the unessential; it is so far from confounding these dignities of being that its whole endeavor is rather directed at discerning and separating them from one another. But since criticism grasps this \textbf{essential} distinction merely as something that immediately \textbf{is}, and everywhere seeks to mark the qualitative limit by the demarcation lines of the understanding, it must continually change standpoint, and so loses its way in a labyrinth of contradictory hypotheses and different regards.\footnote{These ``different regards'' are, so to speak, the squalls of external reflection, which, if they gain the upper hand, must cause every philosophical idea to capsize. It is with the critical philosophy as with critical science in general: it always says: Insofar — insofar! But the insofar-philosophy is an old man's philosophy, which despite all \textit{reservationes mentales} must \textit{crepere} from scrupulosity. The various divisions into \textit{Articuli fidei fundamentales} et \textit{non fundamentales} can, as far as the Church's doctrinal system is concerned, furnish an example.}
 The next advance is that the different regards, with all the dignities of being that follow from them, are grasped in their whole mutual relativity, whereby every immediate determination, whether in its pure immediacy it presents itself as essential or as unessential, is sublated. This sublation of all superficial
% --- p. 87 ---
\setcounter{footnote}{0}%
\opage{87}immediacy characterizes the standpoint of \textbf{ingenuity} and of wit, where one seeks in everything a deeper meaning (ὑπόνοια).

On the one side, then, one has here the whole sum of determinate-existence-moments, which are all unessential because they are all immediate; on the other side, by contrast, pure essentialness, which, in giving all determinations their essential meaning, thereby at the same time determines itself. The task is then to demonstrate the relation of the immediate moments to essentialness; for in this relation they are all essential. On the standpoint of ingenuity, where even the most insignificant thing receives its significance, the distinction that \textbf{is} between the essential and the unessential is thus sublated. The method aims at negating abstract immediacy, which as such is one with unessentialness itself; but this negation is a reinterpretation, whereby every moment entirely loses its unessential side, and essentialness itself is identified with the sum of the relative essence-moments. Ingenuity is still caught in the critical dilemma: either A must be essential or unessential. The being of essence in the totality of the relative moments is thus after all still an immediate being; one procures only a superficial immanence, without bringing it to consciousness that essentialness must necessarily have unessentialness, as such, in itself. The one-sidedness therefore also betrays itself in a double way: partly, namely, one can develop essentialness only by a successive reinterpretation of the unessential moments, whereby the development loses itself in an infinite becoming and approximation that never reaches absolute essentialness itself; partly the essentialness presented by the reinterpretation is drawn down into the relativity of the moments to such a degree that in them it
% --- p. 88 ---
\setcounter{footnote}{0}%
\opage{88}loses all essential meaning, and wit ends in witlessness.\footnote{Examples of this we have in superficial teleology, tasteless typology and arbitrary allegorizing.}

For if the unessential has itself become essential and been posited as the immediate expression of essentialness, then the essential has hereby also become unessential. With the vanishing of unessentialness the fleeting negation vanishes, as that which concerns only external consideration, and only pure abstract immediacy remains; but this is precisely the true characteristicon of unessentialness.

\textbf{Skepticism} shows us how the immediately essential constantly passes over into the unessential; every affirmative being dissolves in the skeptical ``tropes''. α) Pure immediacy, the thing as such, \textbf{is} only \textbf{for} a consciousness; the being of the thing therefore does not affirm itself, but the affirmation belongs to self-consciousness, which is the thing's non-being. Since it is thus not the thing itself, but only the subjective assertion about the thing, that has significance, every thing, as such, is \textbf{unessential}: ``Man is the measure of things''.

But the affirmative being of consciousness characterizes itself in that it can negate any determinate being whatever: its self-affirmation lies precisely in the negative operation whereby it negates everything, itself too. The essentialness of self-consciousness thus consists in its being able to posit itself as \textbf{unessential}.
β) If the negation by which consciousness negates its own being and that of all things is in turn negated, then the negative operation does thus indeed assume a
% --- p. 89 ---
\setcounter{footnote}{0}%
\opage{89}positive tinge; but the essential immediacy in which the being restored by the negation coincides with itself is nevertheless only the indeterminate and all-determining opining, pure postulating. But every postulating must postulate something. Insofar as something is postulated, other is negated. Now since something is postulated only because it is postulated, and other is negated merely because it is negated, one can conversely just as well negate something and postulate other. If the ground why something is affirmed is not to lie in the affirmation itself, then it must be sought in an other that is already affirmed; but if this affirmation in turn is not to be groundless either, then it again presupposes an other, and so on to infinity. The axioms of dogmatism are all groundless and in consequence \textbf{unessential}.
If one finally reduces the affirmation of something and other to their reciprocal relation, so that consciousness is affirmed because it affirms the thing, and the thing in turn is affirmed by the affirmation of consciousness, then a ``Diallelus'' sets in, which makes the confirmation dependent on that which itself depends on
the confirmation.

The strife between these different standpoints arises from a one-sided conception of the infinite \textbf{relativity} of pregnant being.
a) Relativity in all determinate being presupposes a being-for-itself. That something is relative means that it is just as much for itself as for other. This side is held fast in criticism.
True relativity expresses the being of the relative moments for one another. Herein lies partly a mutual affirmation, partly a mutual negation.
% --- p. 90 ---
\setcounter{footnote}{0}%
\opage{90}Ingenuity grasps the relation from the affirmative standpoint; what is negated here is only the isolated being-for-itself of the moments. The one immediate determination coincides with the other, and as all being-moments are joined to and supported by one another, all-being consolidates itself as the positive essentialness. Ingenuity recognizes the pregnant immediacy of being, but on the other hand does not notice that the one immediate being-moment is precisely, as such, the negation of the other, and that true being continually vanishes as one tries to grasp it and hold it fast.
Skepticism takes relativity from its negative side; for it the sum of all the moments of being is only a totality of mutual negations. Immediacy vanishes each time it surfaces, and with it every affirmation vanishes too. Skepticism has an eye only for the pregnant negation, under which all determinations mutually consume and dissolve one another, but it does not feel the infinite presence of essentialness in all negative movements, and can therefore never reach the categorical expression of true being. % print: Die (for Øie) (misprint)
Absolute relativity, on the other hand, presents to us the pregnant immediacy in the pregnant negation. While skeptical relativity expresses itself only in an external relation between negative moments, so that its negation is successively infinite without being absolute, absolute relativity on the contrary concentrates all negative moments in a negative unity, in which every negated determination refers itself to all the others as to itself, and thereby raises itself to affirmative immediacy. According to relativity, the one immediate
% --- p. 91 ---
\setcounter{footnote}{0}%
\opage{91}being-determination is the negation of the other; herein lies: α) that immediate being \textbf{is} negation, and thus the negative of itself; β) that its negativity \textbf{is} the very other in which it is negated; γ) that this other also \textbf{is} the negative of itself.

All negative moments thus \textbf{are}, namely as negated, and the one is what the other is, for it is a common determination for them all that each of them for itself is the negative of itself. Every negative determination is thus in its negativity itself and only itself alone; this its relation to itself constitutes the first main side of immediacy. Next, the negative moment in its negativity is precisely the negative of itself; it is itself all other negative moments. It therefore does not put itself in the place of other and exclude otherness, as the One does in being-for-itself; otherness \textbf{is} in the negative moment as negated; but negated otherness is otherness in its own otherness, thus otherness in its truth: the negated moment is itself its other, since it is an other of itself. This relativity expresses the second main side of immediacy.

But this pregnant immediacy, which is negativity's \textbf{own immediacy}, is consequently at the same time \textbf{immediacy's own negativity}.

Negative immediacy is the unity of being in which all negated moments sink together as in their universal negation. The universal negation is thus a negation of all negative moments and, as such, their affirmative unity. As the unity that affirms the negative moments, it does indeed affirm its own other; but since its relation to its other is its absolute relation to itself, it affirms itself in affirming its own negativity. Every
% --- p. 92 ---
\setcounter{footnote}{0}%
\opage{92}immediate being-determination that can vanish in it has thus already vanished, and every one that can come \textbf{has} already come; as the negation of all coming and vanishing, the unity is the positive center of all-being, its absolute remaining-in-itself. This central being, purified in negation, is the \textbf{essence} of all-being. As the immediacy of negativity, essence is the affirmation of negativity in negativity itself; every other being that shows itself with pure immediacy must then just as immediately show its immediacy negated; every immediate being will then henceforth have to be regarded as immediacy's own negativity, as a being that immediately expresses its non-being, i.e. as \textbf{semblance}\footnote{Think here, e.g., of the thinker's absolute relation to himself through his thoughts.}.
 It will be evident from this that the pregnant immediacy and self-affirmation which essence vindicates through the pregnant negation can only be posited and developed as all the moments of determinate existence, as semblance and being-reflexes, are sublated in the dialectic of essence.
a) The self-determination of essence thus accomplishes itself in a
circular movement, during which it takes up all the moments of being: \textbf{the dialectic of essentiality must have} \textbf{its own peculiar sphere}. b) The same holds of semblance, the essential negation of essence. Although the being of semblance must be deduced from essence, it nevertheless has a being; just because the being of essence is the original and that of semblance the deduced, the proper dialectic of semblance can only develop by positing essentiality as a negated moment; semblance thus everywhere sublates essence, and
% --- p. 93 ---
\setcounter{footnote}{0}%
\opage{93}moves within itself as phenomenon. \textbf{Thereby} \textbf{the phenomenon obtains its logical sphere}. c) But as essence and phenomenon, each in its sphere, round themselves off into dialectical totalities, it turns out that these are only relative. Their true distinction of dignity must be sought in their absolute relativity. According to absolute relativity the totality is referred to itself as to its other, and to its other as to itself. \textbf{The pervasive unity of the totalities}, \textbf{which will overcome the} \textbf{essential} dualism, \textbf{unfolds itself} in \textbf{the sphere of actuality}.

\afsnit{A. The Sphere of Essentiality.}
\underafsnit{a) Pure Essence.}
\paragraf{21}{Essence, Semblance, Reflection.}

\emph{Pure essence}, i.e. the absolute unity of being, which sublates all distinction in negating all immediate and relative being-moments, is, by virtue of this universal negation, first to be determined as the absolute nothing. But the immediate determinations of being are annihilated in the absolute nothing only in that they annihilate one another; since they are thus posited in their nothing, they are posited in their truth. Finite being is hereby recognized for what it really is, namely negativity's own fleeting immediacy, pure \emph{semblance}. Identity, the essential unity of essence and semblance, is the unity of distinction, and opposition, the essential
% --- p. 94 ---
\setcounter{footnote}{0}%
\opage{94}distinction of both, is the distinction of unity: \emph{reflection} expresses the self-movement of essence and its mirroring in all the moments of semblance. % print: udtryk- (for udtrykker) (misprint)

Although being in the preceding chapter has separated its categorical determinations, and has thereby gradually been dissolved into a multiplicity of limited and modified moments, yet with all this it is the \textbf{one and the same} being that, by negating its immediate equality with itself, has posited its distinction-moments, and, by a successive negation of every specific determination, has sublated all its immediate categories, in order to coincide with itself in absolute unity.

In \textbf{pure being} all specific being-determinations are suppressed, and their place is taken by pure nothing; in \textbf{absolute indifference} the distinct determinations develop at once into indifference and into equal validity. Insofar now as their indifference is their invalidity, indifference presents itself as the blind unity of being; but insofar as the moments, on the other hand, in their indifference each for itself have the same validity, indifference is abandoned to infinite difference, and being never attains any true determination. In pure \textbf{essence}, by contrast, all being-moments are not merely developed and posited, but also negated and sublated; their immediate validity is hereby determined as their true invalidity, and essence posits itself as the absolutely valid, i.e. as the true unity of being. Since essence, as absolute being itself, cannot change, it follows that it cannot immediately be ``something'' either; for every ``something'' stands over against an ``other'', and is as such subject to ``change''. Just as little can essence be held fast as
% --- p. 95 ---
\setcounter{footnote}{0}%
\opage{95}an immediate being-for-itself whose qualitative unity would exclude all relative being-elements, for thereby these, as for-itself, would in turn have to exclude it. Through such a reciprocal exclusion essence would then present itself only as ``One'' among ``Many'', and thus be drawn into the circle of finite relations.

The only determination of essence then becomes that of negating every determination, whereby it then again seems to present itself as pure being. But the essential distinction between being and essence is nonetheless, when more closely considered, conspicuous. For while pure being, in its logical innocence, as yet knows of no negation, and is therefore, if we may say so, surprised by the delimiting nothing, so as to be determined by it in its whole indeterminateness, essence on the other hand has taken up negation, and with it the principle of all determinations, into itself. Essence is not negated by anything, since it rather itself negates everything; it does not \textbf{let} itself be determined, but is itself that from which all determination proceeds. The determinateness of essence is thus to be posited as its absolute self-determining.

But although essence, as being concentrated in itself, is now regarded as the logical wellspring from which all the moments of being emanate, and consequently must itself yield determinations to infinity, this its determinability is nevertheless by no means to be confused with the readiness of indifference to assume any determination. Indifference falls prey to difference, and makes do, so to speak, with the being-moments as chance may have it, just because it must \textbf{let} itself be determined. It does indeed have negation under all particular shapes in itself, but it does not itself master this negation; its negation therefore constantly alters its
% --- p. 96 ---
\setcounter{footnote}{0}%
\opage{96}being, and thus prevents it from all self-determination. Essence, on the other hand, has negation itself in its power, and therefore also masters all relative being-determinations. Every immediate being that attempts to escape the central unity and affirm its own immediacy is then at once overtaken and overpowered by the irresistible negation of essence, which compels it to confess the Whence? and Whither? of its movement. The immediacy that is becomes from its nothing; it thus comes from the negative of itself, i.e. from essence; in its becoming it therefore also goes back into its nothing, i.e. into the negation of essence.

Since essence negates every immediate being, it must consequently also negate the immediacy of its own being. By its two-edged negation it does indeed win for itself a pregnant immediacy, but its two sides always constitute one single, purely immediate unity. According to the pure immediacy of this unity, essence is, in its self-determination, continually on the point of becoming ``something'' and of presenting itself as something that is. Now if every immediate something has fallen prey to the negation of essence, then this negation must also continually sublate the being of essence itself. Absolute being is thus posited as the in-itself absolute nothing, and the positive side determines itself as the negative. If essence is in itself negativity's own immediacy, then it is eo ipso also immediacy's own negativity: \textbf{pure essence is}, \textbf{as such}, \textbf{pure semblance}.

The relative, finite nothing, which negates only a certain something, is in itself just as much a being as a nothing; it therefore also presents itself as a determinately existent nothing, and assumes the positive stamp of otherness; but the in-itself absolute nothing, which tolerates no immediate being whatsoever, can on the other hand never itself enter into determinate existence.
% --- p. 97 ---
\setcounter{footnote}{0}%
\opage{97}The result then seems to be this: that essence, according to its self-determination, must deny itself every being-determination, that it thus does not exist at all, and that only finite being, seen in its own infinity, i.e. in the wealth of all its moments, is true being.\footnote{In this result atheistic logic acquiesces; for atheism, logically taken, is the denial of essence.}
 But if essence is without determinate existence, then determinate existence must also be essenceless, and, as such, present itself in its whole unessentialness. The dialectical development of determinate being does indeed show that something must itself be its other, and that it thereby becomes something other than it was before; but this sublation of the immediacy of something is, on the standpoint of being, only an alternation with a following immediacy. The negation of change is thus merely a progression from the one immediacy to the other. What is negated here is only the given, determinate immediacy-moment; the negation is therefore partial and successive. The something that perishes is annihilated only as \textbf{this} or that determinate something. But every something that follows must always be determined as \textbf{this} or that: the ``\textbf{determinate} something'' thus remains, throughout the whole change, really unchanged. Herewith, then, the being of change is sublated and posited as a semblance. How far change is to hold depends solely on how far something can immediately separate itself from other. As long as there is a distinction that is between something and other, so long is there also a becoming and a transition from the latter to the former. But according to the pregnant negation the immediate being of something is really its non-being, its being as not-something. Something is in itself a semblance; from this it
% --- p. 98 ---
\setcounter{footnote}{0}%
\opage{98}comes that its transition into other is no transition. The other of something then, in consequence, \textbf{is} not a real other either, but only the reflected semblance of other, by which something points beyond itself. As other now casts its semblance on something without itself being absorbed in it, it reserves its pure immediate being for itself. Other is thus now the real something, the real, which seems to have saved itself out of the nothingness of semblance. But if other is now held fast as something, then it is, according to what precedes, again a semblance in itself and a reflected semblance of other. Thus determinate existence with all its multiplicity becomes transparent, and its immediacies clear into innumerable reflected semblances. Each semblance in this realm of finitude is a fleeting immediacy that in its negation points beyond itself to something that is; but in the place of what is a new semblance continually enters, and affirmative being flees away into infinity. In semblance itself there is not anything, and every something outside semblance transforms itself into semblance; thus determinate existence stands as absolute emptiness: ``Alles Seyn ist Schein''.\footnote{Herewith the standpoint of acosmism is indicated. The acosmist and the atheist can unite in recognizing the transience of the world. But they grasp it in opposite ways. The atheist regards everything from the standpoint of being and change, and can therefore indeed concede the perishability of things in the particular; but since the partial immediacies are continually restored in change, he insists, as far as the whole is concerned, on the eternal apriority and imperishability of the world. The acosmist, on the other hand, sees through the nothingness of the whole of determinate existence, in that he grasps, from the standpoint of essence, the content of all-being as semblance of semblance and the world as the realm of illusions.}
 The nothingness in which all-being now seems to evaporate must, however, itself vanish as semblance when the
% --- p. 99 ---
\setcounter{footnote}{0}%
\opage{99}essence-negation is carried through. What is terrifying in the thought-movement of nihilism lies in this, that it presupposes a distinction of dignity between immediate being on the one side and essenceless, illusory semblance on the other. It ceaselessly seeks a being which in its truth and purity should be freed from every semblance, and since it then everywhere finds that such a being can only be determined as the pure other of semblance, and that every such other, as the negation of semblance, is, when one wants to hold it fast, a reflected semblance of semblance and in itself a semblance, it leads consciousness to the disconsolate result that everything is \textbf{only} semblance.

One rejects semblance because one trusts in the pure immediacy of being. If, on the other hand, it becomes evident that it is precisely immediate being which, when more closely tested, is itself semblance, then being has its truth in semblance; semblance then stands, as such, above mere being, the relation of dignity turns around, and one must admit that every immediacy that is supposed to lie outside semblance is \textbf{only} a being. Now if true being is to be sought in semblance, then it is hereby given that every other being that is not sublated in semblance must for just that reason be regarded as a real semblance-being. The essential separation between true and untrue being will therefore now have to be posited as a \textbf{separation between true} \textbf{and false semblance}.

Pure semblance is in the first place an immediacy that is, which represents only itself; but as fleeting immediacy, i.e. as immediacy's own negativity, it must next just as immediately turn over into its negation. This sudden reversal, which repeats itself everywhere in the infinite series of semblance, always brings with it, for naive, credulous consciousness, a
% --- p. 100 ---
\setcounter{footnote}{0}%
\opage{100}surprise. He who does not suspect that immediacy is itself negativity will, when he sees the pure immediacy of semblance, want to capture and hold fast positive being in it; but as the semblance-immediacy now at once betrays the secret of its negation, he sees himself deceived, and will then leap over semblance entirely in order to grasp reality. It is this infinite self-separation, in which semblance continually repeats itself anew, that produces the illusion: ``\textbf{Semblance deceives}''.

The single semblance in its self-negation refers to another semblance; the flight of immediacy and positivity shows itself as an uninterrupted transition from semblance to semblance; thus the affirmation of semblance lies in its negation. But as semblance through its negation mirrors and confirms itself, it is completely sublated; for it conflicts with the determination of semblance to affirm and mirror only itself alone. The absolute semblance, which shines through all reciprocal semblances, is then at the same time the true other of semblance, the real in semblance, whose immediacy is not fleeting but positive and abiding; thus \textbf{essence}, the \textbf{absolute negation} of semblance in \textbf{semblance itself}, is \textbf{clarified as the pure} \textbf{light-} (τὸ

The light is the unity of semblance; this unity is won through the infinite self-separation of semblance; it is thus the unity of distinction, and as such \textbf{identity}. As the affirmation of semblance in its negation, identity is semblance-unity and itself semblance; but as the negation of semblance in its affirmation, as its sublation in other, it is essence-unity, indeed essence itself, and distinction is semblance. Since identity is thus both semblance and essence, it is at once the negation of distinction and at the same time itself distinction. Both these
% --- p. 101 ---
\setcounter{footnote}{0}%
\opage{101}moments must then be posited and negated in the infinite reciprocal semblance between essence and semblance, i.e. in \textbf{reflection}.

The all-round development of reflection characterizes its three main shapes: 1) As reflection of \textbf{movement} it expresses the infinite transition of identity and distinction into one another; 2) as \textbf{reflection of determination} it fixes the opposition of identity and distinction; and finally 3) as \textbf{absolute} reflection it shows the true reciprocal semblance of movement and determination.

1. The reflection of movement is a) \textbf{positing}, b) presupposing and c) \textbf{opposing}.

a) Immediacy, pure abstract identity, negates itself and vanishes as semblance in presenting itself as its own distinction; the immediate thus passes over into the negative. This transition, however, is no becoming; what immediacy hereby \textbf{becomes}, namely the negative of itself, it already \textbf{was} before the transition took place. Here, then, one does not begin with the immediate; rather the reflection-immediacy \textbf{results} as an identity-medium between two negative extremes. The transition thus takes place from the negative to the negative, whereby negativity refers itself to itself, so that immediacy is only the unity in which the relation is held fast. Distinction, the negative that proceeds from immediate identity, is precisely that from which identity, or the immediacy that is, has itself proceeded. It is the negation of immediacy that sublates identity; but since immediacy, as semblance, is itself this negation, identity in turn sublates distinction. The return of the immediate to itself through negation is
% --- p. 102 ---
\setcounter{footnote}{0}%
\opage{102}the very return of negation to itself through immediacy. Immediacy is posited in resulting from the sublated negation. But it is not posited by the abstractly negative, for this too is resulting, and is itself \textbf{posited} as the sublation of the immediate.

What is \textbf{posited} and sublated in this transition is thus a fleeting moment, a semblance that vanishes as it comes, and comes again in the same instant as it vanishes; essence, on the other hand, is the infinite \textbf{positing} and \textbf{sublating}, and as such movement and reflection itself.

As pure positing, essence is a movement from semblance to semblance, and consequently a movement that does not move itself, i.e. a semblance-movement that is itself posited by the unmoved. As the essence-identity, through the sublation of the posited being of semblance, would affirm itself, it is absorbed in an essential distinction; for it becomes at once both the unmoved, which posits but is not posited, and the negative of this, namely the positing movement that is itself posited. As the not-posited but positing, essence is the \textbf{principle} of movement, movement on the other hand, with all that it deposits, semblance and \textbf{result}. But as the principle of movement, the principle must express its essentiality in movement itself; it thus proceeds from its own posited and sublated movement, and is accordingly itself result. Principle and result are now identical; but since the principle is principle only by \textbf{positing}, and the result result only by \textbf{sublating} the movement, the movement itself is essential, and expresses the true distinction between principle and result.
% --- p. 103 ---
\setcounter{footnote}{0}%
\opage{103}With the principle the positing movement \textbf{begins}, with the result it \textbf{ends}. Since the principle is itself result, the movement must end in the same instant as it begins; but since the result is itself principle, it must also begin in the same instant as it ends. The pure positing of essence is thus an infinite converting. Ahead of every result there goes a positing that is negated in the result itself; but since the result is hereby converted into principle, the sublated positing again proceeds from its own negation and continues itself.

As the positing that must already be posited in order to be able to be posited, and for that reason posits its own sublation in order to posit itself, reflection is \textbf{presupposing}.
b) Presupposing reflection determines the essential identity and distinction of principle and result. In order that the principle can posit its result, it must itself already be. But in this abstract being-for-itself it is only a semblance; its true essence lies in the positing by which it posits the result. Before it has posited itself, it is thus not an essential principle, and before the principle has become essential, it can deposit no result; but conversely, it cannot posit itself either except by depositing its result. In the essence of principle and result there thus lies a double presupposition. α) In order to posit itself, the principle posits its own negation and sublation, namely its result, but it itself sublates this negation; for it posits itself only by positing its result as that which is not yet posited but is to be posited. According to this \textbf{primitive} presupposition the principle posits itself \textbf{in advance} by presupposing
% --- p. 104 ---
\setcounter{footnote}{0}%
\opage{104}its result. The result is presupposed as a semblance vanishing for the essence of the principle, and thereby presupposed as that which is to be posited; with this positing the essential negation of the principle sets in, and the result posits itself. But the negation of the principle is an essential moment in its self-position; the result thus posits itself as that which does not posit but is posited; it presupposes its presupposing principle; its true essence is its self-sublation, and its presupposing \textbf{consecutive}\footnote{Since God has in advance posited his world (primitive presupposing), the world now also presupposes its God (consecutive presupposing).}. γ) This presupposing, doubled in itself, is however by no means to be regarded as a composition of two heterogeneous fragments. The primitive and the consecutive presupposing are identical. But their identity is no tautology either; each of the two terms of reflection is itself the whole reflection. They are both equally independent, and in this their mutual independence each other's complete negation and distinction. The developed presupposing is thus one single positing reduplicated in itself, which uninterruptedly turns \textbf{against} itself, and thereby characterizes itself as an \textbf{opposing}. As opposing reflection doubles its own independence, it deposits two centers, and thus goes, through a double movement, back into itself. The principle takes itself back through the result, and writes, so to speak, its
result as a periphery around its own center. It is thus at once positing, posited,
% --- p. 105 ---
\setcounter{footnote}{0}%
\opage{105}presupposing and presupposed, and its movement, complete in itself, constitutes its immobility. Conversely the result reflects itself in its principle; its self-sublation is its self-position. But the self-position of the result is the essential negation of the principle: the result thus reflects its immobility in itself, and rests in the sublated principle as in itself alone. Herewith absolute identity is posited as absolute distinction. The principle is now no longer principle, for its transition into the result is completed; just as little can the reflection of the result have any specific significance, since the result is now rather itself principle. Each of the opposed terms is thus closed in itself, the movement has vanished, and only pure opposition remains\footnote{The essential distinction between pantheism and atheism, dualism and deism on the one side and theism on the other shows itself as a distinction between the one-sided and the all-round conception of the reflection-movement. Pure reflection describes a circle. Pantheism is now separated from atheism (pancosmism) in that the former makes God, the latter by contrast the world, the center of the circle. But one-sided reflection without essential opposition is an abstract unessentialness, and the single circle a big zero. Dualism (deism), on the other hand, keeps to opposition. For the dualist, determinate existence presents itself under the image of two circles touching each other. But since one here forgets that the presupposed primal opposition is posited in advance, thus not unoriginated but sprung from a deeper identity, the double center of dualism can also be said to represent pure incomprehensibility. Only theism grasps the true reflection-movement in its whole self-doubling. In recognizing the center of God and that of the world as sprung from the divine essence's own absolute centrality, theistic reflection sublates the circles of dualism: ⊙⊙ in the symbol of infinity ∞.}.
% --- p. 106 ---
\setcounter{footnote}{0}%
\opage{106}2. In the presupposed opposition each term is independent, \textbf{determinate} in itself, and must be taken as it is given. Herewith movement vanishes as semblance in the reflection of determination.

Determination-reflection is a) \textbf{distinguishing}, b) antinomic and c) \textbf{fixing}.

a) For distinguishing reflection
opposition presents itself in its pure immediacy, and is therefore held fast as a co-being of two identities touching each other: A = A, B = B.
As the reflection-determination is now posited as the identity acquiescing in itself, movement vanishes as semblance; it is hereby posited as something wholly irrelevant to identity itself, as a juxtaposing and comparing that takes place merely in the subjective consciousness of the observer, thus as an \textbf{external} reflection. The observer's eye rests now on A and now on B, and subjective reflection now impresses their distinction on itself: A is not B; B is a not-A; A is a not-A; A is not not-A. In the pure identity of A with itself there lies, then, no determination at all of any not-A whatever. Indeed the very proposition: A is A, expresses only a reflection-identity that falls outside A itself. For A is here posited twice; but in its pure identical being-for-itself A cannot separate itself from itself. The separation is thus owed to external reflection, which, in its effort to hold fast identity in its perfect determinateness and to \textbf{distinguish} between A and not-A (B), has % print: Die (for Øie) (misprint)
% --- p. 107 ---
\setcounter{footnote}{0}%
\opage{107}had to posit A over against it, and by the copula ``is'' to undo this opposition again. Pure determinateness, freed from all contradiction, can only be expressed as an interjection-identity: A!, and the substitution by which subjective reflection puts A in the place of A (A is? --- A!) is then hereby reduced to a perfect tautology\footnote{Tautological reflection, whose criterion is that it regards contradiction as a subjective error of thought that can be removed by a distinguishing between the identities that are, has found its most extensive development in Wolfianism. Wolf defines contradiction as a ``\textit{simultaneitas in affirmando et negando}'', without an inkling that he hereby defines the very principle of true reflection. In describing identity: ``\textit{Eadem dicuntur, quæ sibi invicem substitui possint}''. (\textit{Ontolog, p. 148}) he adds: ``\textit{Facta nimirum substitutione perinde esse debet, ac si nulla facta fuisset}'', without observing that every such substitution thereby becomes entirely meaningless.}.

But even this tautology, which strives to escape contradiction, has its real life-principle in contradiction itself. For subjective reflection runs through the following propositions in its movement: A is (perhaps) not-A; A is not not-A and finally A is A. That it puts A in the place of A has significance solely in that it would otherwise have to put not-A in the place of A. But in this way it puts precisely A in the place of (the feared) not-A. Only by contradiction itself can subjective reflection distinguish itself free of contradiction. Contradiction is thus no longer an error of thought but a necessity of thought,
% --- p. 108 ---
\setcounter{footnote}{0}%
\opage{108}and reflection is herewith recognized as \textbf{antinomic}. Through antinomic reflection it is brought to consciousness that duality, opposition and contradiction belong to reflection's own essence; but despite this the reflection-movement from which the determinations result is nevertheless posited as a semblance-essence, whose propositions must be sublated as soon as they are posited; true essence, by contrast, keeps itself in its pure indeterminable unity, and every attempt to determine it is annihilated in the ``antithetic'' of reflection. \emph{Thesis}: \textbf{True} essence is, when \textbf{it is more closely determined}, \textbf{absolute} identity.

Pure identity is \textbf{distinct} from distinction and \textbf{opposed} to opposition. Herewith it has sublated its abstract being-for-itself and itself posited itself as the distinction of distinction and the opposition of opposition; absolute identity is thus in itself opposition. \emph{Antithesis}: \textbf{True essence must} \textbf{be determined as absolute} \textbf{opposition}.

In every opposition the one term must necessarily presuppose the other; for in order to be opposed, they must both be \textbf{posited} against each other: opposition is the inclusion of the opposed terms. But since each of the terms in the opposition is opposed to the other, they must mutually exclude each other; opposition is thus the mutual exclusion of the terms. Inclusion and exclusion are themselves opposed to each other. Thus opposition is, then, on all sides in opposition with itself, and thereby, as absolute opposition,
% --- p. 109 ---
\setcounter{footnote}{0}%
\opage{109}opposed to itself. But the opposite of opposition is identity\footnote{Antinomic reflection is characteristic of Kantianism.}

The essence that remains in itself and is indeterminable in its pure determinateness is, however, not to be confused with the absolute indifference of being. Indifferent being is continually on the point of assuming any determination whatever, and is thus in an uninterrupted movement and unrest. The essence that is in itself, by contrast, keeps all reflection-determinations out, sets them off from itself as semblance and reflected glow, and thus acquiesces in its pure \textbf{transcendence}, lifted above all oppositions.

c) When this separation of identity and opposition, determination and movement, is carried through at every point, and thereby driven to the highest pitch, reflection shows itself as \textbf{fixing}. Fixing reflection holds fast, in the strictest way, the pure being of essence in its distinction from all negation. Even the reflection-determinations must in consequence be fixed as immovable. The reflection-determination fixed in itself\footnote{Fixing reflection has been carried through with a rare energy in the Herbartian philosophy.}
is now to be posited as a for-itself essentiality, whereby the primal identity of essence is shattered and separated into an innumerability of identity-fragments. Each of these fragments is itself an in-itself, a pure identity, a tautology, thus a contradiction. The omnipresence of tautology is also the omnipresence of contradiction; every specific determination that
% --- p. 110 ---
\setcounter{footnote}{0}%
\opage{110}subjective reflection attributes to the determinate it must therefore continually take back. The immovable moments of essence show themselves to reflection in an uninterrupted movement, although reflection continually assures us that they do not move; in their mutual determination they continually pass over into each other, and reflection must itself employ the negation of the one in order to determine the other, although it is nonetheless asserted that the one identity-fragment (``reals'')
does not concern the other. Reflection here entangles itself in an infinite self-deception; in fixing its determinations so one-sidedly, it fixes itself.

3) This illusion is sublated in \textbf{absolute or all-round} reflection, which, as a) \textbf{optical-intuitive}, b) \textbf{dialectical} and c) \textbf{dialectical-speculative}, lets determination and movement mirror each other. Transcendent essence, which hovers above the opposition of reflection, thereby becomes opposed to opposition itself, and is thus determined in itself as
the absolute identity of opposition. Essence now mirrors its immobility in its movement, it beholds itself in all its reflection-determinations, and its dark transcendence transforms itself, through optical-intuitive reflection, into clear \textbf{immanence}\footnote{Schelling.}. But the determinateness beheld in immanence is as it is, and it is not yet evident why it must move; movement is beheld as absolute, and it is not yet recognized why it must deposit the determinate. The identity of movement and determination is thus still tautologically immediate, and the contradiction evident.
% --- p. 111 ---
\setcounter{footnote}{0}%
\opage{111}b) The contradiction immanent in this identity of movement and determination first comes to true development in \textbf{dialectical} reflection.

Absolute identity is, as such, itself identity and opposition, and developed opposition likewise at once both opposition and identity. Each term in the opposition is thus for itself the whole identity and the whole opposition. It maintains its integrity by excluding the other term; identity excludes opposition just because it is itself opposition and has opposition in itself. Since each of the opposed terms now has its independence in itself by having its other, and has this other only by excluding it, it excludes its independence in the same instant as it would vindicate it; it steps out beyond itself and comes into movement just as it would remain in itself and fix its determinateness. According to this immanent contradiction, each of the terms has its inequality sublated in its own equality, and is thus to be posited as the \textbf{positive}; but conversely it also has an inequality reflected in equality, and must therefore just as well be posited as \textbf{the negative}\footnote{Dieß ist also derselbe Widerspruch, der das Positive ist, nemlich Gesetztseyn oder Negation, als Beziehung auf sich. Aber das Positive ist nur an sich dieser Widerspruch; das Negative dagegen der gesetzte Widerspruch; den in seiner Reflexion in sich, an und für sich Negatives oder als Negatives identisch mit sich zu seyn, hat es die Bestimmung daß es Nichtidentisches Ausschliessen der Identität sey. Es ist dies gegen die Identität identisch mit sich zu seyn, hiermit durch seine ausschließende Reflexion sich selbst von sich auszuschliessen. Das Negativeist also die ganze, als Entgegensetzung auf sich beruhende Entgegensetzung, der absolute, sich nicht auf Anderes beziehende Unterschied; er schließt als Entgegensetzung die Identität von sich aus; aber somit sich selbst, den als Beziehung auf sich bestimmt er sich als die Identität selbst, die er ausschließt. Hegel.}.
% --- p. 112 ---
\setcounter{footnote}{0}%
\opage{112}c) But when the positive in its development is itself the negative, then the identity of developed positivity and negativity, more closely considered, is developed tautology itself. The contradiction is then not yet resolved, and cannot possibly vanish as long as determination itself is not yet more closely determined. The positive must then be posited against itself as a This! against This! But this specific antithesis reflection, as such, cannot grasp; for to pure reflection This = This, as A - A. Fixed determinateness culminates in immediacy, i.e. in the negation of the reflection-movement. The demonstrable positivity-moment must be beheld as it is posited by essence itself. It is not, however, sensed as something that immediately is; its specific positivity is at the same time its specific negativity; it is, as the immovable, a moment within the dialectical movement itself. Integrated by this determination-moment, reflection is \textbf{dialectical-speculative}. As essence, as infinite negativity, deposits the positive, specific determinations, and thereby at the same time posits them as the negative, immovable moments of the reflection-movement, the thetic distinctions are preserved and mirrored in reflection itself. The dialectical stepping-forth of the thetic moments from the pleromatic center of essence thus constitutes the real presupposition, and reflects the own speculative ground-type of absolute essentiality.
% --- p. 113 ---
\setcounter{footnote}{0}%
\opage{113}
\underafsnit{b) The Duality of Essence.}
\paragraf{22}{The Primitive Essence, the Consecutive Essence, Subordination.}

In the absolute identity of all positive reflection-determinations essence has its infinite \emph{fullness}. As the negative unity of the fullness (transcendent hypostasis) it must deposit the specific moments of determination; but since these vanish as semblance in its own self-position (transcendental hypostasis), the fullness remains esoteric, and the opposition reflected in it is sublated in the \emph{primitive} essence's own for-itself identity (absolute hypostasis). The infinite positivity of the primitive essence is at the same time its infinite negativity; the esoteric fullness is therefore negated in the exoteric, and distinction enters in all-round specification. The independent unity of the exoteric fullness, reflected in itself, is hereby posited as the primitive essence's own central opposition, i.e. as \emph{consecutive} essence. But as the esoteric % print: esotoriske (misprint)
fullness, in the presupposing reduplicated in itself, subordinates the exoteric to itself, the primitive essence itself becomes subordinated to the consecutive, and with this reciprocal \emph{subordination} the central opposition perishes. % print: exotoriske (for exoteriske) (misprint)

Determinate existence has dissolved in its nothing; it vanished as semblance and shadow in pure essence. Essence in this its abstraction was itself a negative being;
% --- p. 114 ---
\setcounter{footnote}{0}%
\opage{114}but from its infinite negativity there developed at the same time the negation and sublation of semblance; the self-reflection of essence constituted itself as principle for all reflection-determinations, and determinate determinateness culminated in the immediacy that is. Herewith, then, determinate existence is reproduced as the being \textbf{posited} by essence.\footnote{In der Sphäre des Seyns war das Daseyn das Seyn, das die Negation an ihm hatte, und das Seyn der unmittelbare Boden und Element dieser Negation, die daher selbst die unmittelbare war. Dem Daseyn entspricht in der Sphäre des Wesens das Gesetztseyn. Es ist gleichfalls ein Daseyn, aber sein Boden ist das Seyn, als Wesen oder als reine Negativität; es ist eine Bestimmtheit oder Negation nicht als seyend, sondern unmittelbar als aufgehoben. Daß Daseyn ist nur Gesetztseyn dieß ist der Satz des Wesens, vom Daseyn. Hegel. Logik IItes Buch. p. 23.}
The whole wealth of specific determinations of determinate existence thus \textbf{is}, and must be taken as it \textbf{is} given; but it is not of itself; as something that immediately is, it is infinite fragility and self-dissolution; its becoming is an uninterrupted resulting, and it itself a result that presupposes its presupposing principle; according to this consecutive presupposing, all determinate-existence-moments are sublated in the negative unity of the primal essence. But as negative unity, essence is not an abstract medium or external receptacle in which the moments could be thought of as preserved as things that are. Their being in essence is precisely their non-being; in the self-position of the primal essence they are posited only as those that are to be posited.\footnote{Denn man kann nicht sagen, daß in Gott sey Feuer, Bitter oder Herbe, vielweniger Luft, Wasser oder Erde, allein man siehet, daß es daraus worden ist. Jacob Böhme.} Insofar as
% --- p. 115 ---
\setcounter{footnote}{0}%
\opage{115}essence is determined as that \textbf{which} posits itself by positing the infinite multiplicity of determinate existence, it is, with this pregnant identity, reflected in itself as in the absolute fullness
(πλήρωμα). Essence is here consecutively presupposed, without however really being posited; it is posited by posited being, thus by that which can posit nothing, since it is itself deprived of all independence. The reflection-movement which essence thus itself has to go through in order to be able to deposit the specific moments of the fullness, reflected in their being-for-itself, shall now be considered more closely: a) Since all reflection-determinations vanish in the primal unity of essence, this unity is reflected in itself as in the negation of the determinations; essence
affirms itself by negating all determinations; its indeterminability is its determinateness reflected in itself. But in this determinate indeterminability its being is infinitely affirmative, and essence is hereby presupposed as the \textbf{transcendent} hypostasis. As essence is hypostatized, it sublates the consecutive presupposing, and determines itself to be the all-positing. Hereby the reflection-determinations are then presupposed as those that are not yet posited. But they are to be posited; for the indeterminability of the transcendent essence is determined solely as the principle for all determinations.
b) Insofar as the transcendent hypostasis now posits itself, it must determine its determinations, and thereby deposit difference and multiplicity as that in which it is itself negated; but insofar as it is to posit itself as the all-determining, the reflexes of the determinations must be absorbed as vanishing moments in its own self-position; only thus can essence appropriate to itself all
% --- p. 116 ---
\setcounter{footnote}{0}%
\opage{116}reflection-determinations, assimilate them to itself, and make them its own. But hereby the transcendent essence at the same time transcends its pure abstract being-for-itself, and opens up its hidden fullness, in order to hold itself fast as the determinate identity of the reflexes. Through this reflection-movement essence is posited by itself, and has thus determined itself as the in-itself determinable, i.e. transcendental, \textbf{hypostasis}. Transcendentality, immediately taken, now contains nothing more than what must already have been given in transcendence itself; but, according to the reflection-movement, it does contain something more. For in pure transcendence all determinations had vanished, and essence was consecutively presupposed; in the transcendental hypostasis, by contrast, they are posited as reflexes in the self-position of essence, and essence has accomplished the movement by which it constitutes itself as the all-positing. What thus lay hidden in the involution of transcendence becomes clarified and posited in the evolution of transcendentality. According to the self-movement that sets in between the transcendent and the transcendental, essence now relates itself to itself as principle to result. The transcendental hypostasis presupposes the transcendent, as the result presupposes the principle. But the principle is only the principle of the result. The transcendent hypostasis with its hidden being must therefore posit its reflection-determinations as the self-negation and opposition in which it first comes to possess itself and its fullness. Both hypostases are thus, each in itself and each for itself, the whole essence with its whole fullness, and by this their being-for-itself reflected in each other as
% --- p. 117 ---
\setcounter{footnote}{0}%
\opage{117}mutual negations. Herewith the pregnant identity of the primal essence \textbf{is} reduplicated in itself as primal opposition and \textbf{duality}.
c) This duality in the primal essence is, however, not to be held fast as a distinction that merely is between essence and essence. Self-movement is no becoming, no successive transition. In order that essence can be a transcendent, it must already have posited itself as a transcendental: the principle is itself result, and the result principle. The negation and distinction that remove tautology lie in the differences and oppositions of the fullness; but since all moments are still only vanishing semblance in the self-development of essence, essence moves from itself to itself, and posits the reflexes of determination as the medium in which it mirrors itself alone. Essence developed in itself is thus at once both positing and posited, and holds fast in its \textbf{absolute hypostasis} the whole fullness as esoteric.

The absolute hypostasis presupposes the developed distinction between transcendence and transcendentality; it is thus neither the one nor the other of these, but keeps itself in its being-for-itself as the identity of duality and opposition. If it were only the developed tautology of transcendence and transcendentality, all determination would vanish again; essence would then once more be transformed into \textbf{a} transcendent abstractum, and the movement would have to begin anew; but if, conversely, the absolute hypostasis posited the transcendent and the transcendental outside itself and, so to speak, \textbf{beside} itself, then the self-affirmation of the primal essence would dissolve into sheer relativities; it would thereby
% --- p. 118 ---
\setcounter{footnote}{0}%
\opage{118}represent itself as a posited determinate existence, which would again have to be sublated and negated. The self-position of essence is then to be posited as absolute identity in absolute opposition. Each of the terms of the opposition is itself the whole identity; thereby identity is opposed to itself, but this opposition is the negation of identity in identity itself, thus its own self-determination, whereby the opposition is sublated in the for-itself identity. Thus the absolute unity of the primal essence is reflected in its developed triplicity (τρεις ὑποστασεις ἐν μια οὐσια), and essence has affirmed itself as \textbf{the primitive} essence.

As transcendent hypostasis the primitive essence is a principle that merely \textbf{is}, an essential semblance that reflects its own distinction: as transcendental hypostasis, by contrast, it is \textbf{posited} as principle, and is thereby its own result; and finally, as absolute hypostasis, it constitutes itself as the \textbf{self-positing}, in itself complete and essential principle.

But the \textbf{essential} principle is only the principle of the essential result; it is therefore reflected and completed in itself as that which is not yet completed. As the self-positing principle the primitive essence vindicates its own central identity with itself; this centrality lies in this, that each hypostasis can for itself be posited as the unity of the two others; as the transcendent is thus posited as absolute, it reflects in itself the transcendental as transcendent and the absolute as transcendental. The self-movement by which the primitive essence posits its own identity and distinction is thus, taken for itself, a semblance-movement,
% --- p. 119 ---
\setcounter{footnote}{0}%
\opage{119}and only as primitive presupposing absolutely essential. The reflection by which the essential principle is reflected in itself must therefore be doubled from all sides, and thereby show itself as a reflection by which it is reflected in other. The primal opposition between the transcendent and the transcendental hypostasis, from which the central identity results, must thus itself be reflected in a central opposition, whereby the primitive essence is negated in the non-primitive. The self-negation of the primal essence is thus reduplicated in itself: partly, namely, it is a negation by which essence posits itself as its own negativum, which, as shown above, produces the primal opposition in the primitive essence itself; partly it is a negation by which the primitive essence, as concrete for-itself identity, posits its non-essence, i.e. its essential other. Without such a central opposition the primal opposition is an empty semblance, and the central identity a transcendent abstractum; but, reflected in its central opposition, the primal opposition is itself the primitive presupposition, and the central identity of the primitive essence the self-positing principle, which in the central opposition deposits its essential result.

This central opposition of the primitive essence can now be sought neither in determinate existence, since this has itself vanished as semblance, nor in pure semblance, since this is only the negation of pure essence, which is itself semblance. The essential negation of the primitive essence must therefore be posited as the semblance reflecting in itself, i.e. as the unessential or consecutive essence.
% --- p. 120 ---
\setcounter{footnote}{0}%
\opage{120}Pure semblance is infinite distinction; it affirms
itself as semblance of semblance in uninterrupted self-separation.
The semblance reflected in itself thus mirrors back
an innumerability of semblances, of which each single one
is in turn a semblance reflected in itself, a semblance-hypostasis.
Insofar now as the semblance-movement everywhere tends
toward a hypostatizing, every semblance takes on a stamp of essentiality,
and vindicates for itself a relative independence;
but insofar as the essential semblance
is only a semblance-essence, its hypostatic
affirmation is absorbed as a moment in its essential negation:
the reflection-in-itself of the consecutive essence vanishes
everywhere as infinite reflection-in-other;
in thereby negating the central identity, and
transforming itself into a semblance-identity for the hypostatized
distinctions of the outer reflection-determinations, it posits
itself as the essence cast outward, with an exoteric
fullness.

The esoteric and the exoteric fullness negate and
reflect each other mutually. In the esoteric,
every distinction-moment can now vanish as semblance,
just because distinction itself is posited and hypostatized
in the exoteric; and conversely, the identity in the
exoteric can be posited as semblance-identity, since the identity
of the esoteric fullness is affirmed and centralized
in the primitive essence itself.

The central opposition between the primitive and
consecutive essence now also presents itself as a
true opposition between the \textbf{positive} and negative.
The primitive essence, which in its abstract
being-for-itself would sink down into a fleeting semblance
without any independence and reflection-in-itself, frees
itself from this abstraction by positing
the consecutive as its true negation, and
% --- p. 121 ---
\setcounter{footnote}{0}%
\opage{121}thus shows that it has infinite negativity in
itself; but it at the same time posits the consecutive essence
as an essential result, which must presuppose
its presupposing principle, and thereby goes uninterruptedly
back into itself, as the absolutely completed. All
negativity and non-independence is hereby posited as a
serving moment for its positivity and independence;
what steps forth in the exoteric fullness
is already given as negative being in the esoteric;
the reflection-movement by which the primitive
essence posits changeability, and passes over
into the exoteric fullness as into its own concrete
transcendentality, is thus the very same in which
it reflects its unchangeability, and vindicates its
eternal transcendence: the \textbf{hypostatized} semblances
of \textbf{the consecutive} essence \textbf{are} all \textbf{radiations
of the primitive} \textbf{light}.

The consecutive essence thereby shows itself as
the self-development of negativity. It proceeds in
all respects from its other. Reflected in the
primitive, it \textbf{is} before it becomes \textbf{posited}; it \textbf{is} presupposed,
as that which is to be posited. But reflected in
itself, it is posited before it comes to be;
its proper being proceeds from the self-position
by which it negates and sublates its presupposition.
But herein lies precisely its essential
contradiction: it is, in order to be posited by its other;
and it is posited by its other, in order to posit itself
and to be. This contradiction finds its resolution
in the fact that the peculiar reflection-movement
of the consecutive essence lets its whole positivity
be absorbed as a moment vanishing in the posited and developed negativity.
According to this peculiar
reflection-movement, the consecutive
% --- p. 122 ---
\setcounter{footnote}{0}%
\opage{122}essence \textbf{is} an independent, for-itself, self-identical
essence. But its identity is negative.
It posits itself as that which is posited by other;
its independence expresses its non-independence:
the position is negation. As identity
would hypostatize itself, it vanishes as semblance in an
innumerability of relative hypostases; and as each
separate essentiality would hypostatize and sustain its
being-for-itself, it posits its own self-contradiction;
as the result of a vanished independence, every
semblance-hypostasis vanishes in another self-contradictory independence. The distinction-moments in the
exoteric fullness are thus produced in order to be destroyed,
and destroyed in order to be reproduced: the essential
result is an identity dissolving itself.\footnote{According to its triplicity, the essence of God is infinitely affirmative; the world-essence, on the other hand, uninterruptedly seeks independence, without being able to attain it in itself. The independence it has in its peculiar self-movement is only a medium for the revelation of God's absolute independence. Could it, on the other hand, in actuality be thought of as torn loose from identity with God, it would thereby become diabolical; for the diabolical essence is precisely the contradiction ceaselessly hypostatizing itself, which can nevertheless never attain a true hypostasis.}

It is the primitive essence that uninterruptedly posits
and sublates the contradictions of the consecutive; and
it is, on the other hand, the negativity in the
consecutive that expresses the positivity in the primitive.
As they thus mutually presuppose each other,
the distinction of dignity is at the same time posited in and with
the reduplicated presupposition: their reciprocal
% --- p. 123 ---
\setcounter{footnote}{0}%
\opage{123}relation is a mutual \textbf{subordination}; for
subordination is the reflection of the non-independent in
the independent.

The consecutive essence cannot vindicate for itself
exclusive independence, for it would thereby
cease to be consecutive; nor can it give up
its independence and vanish, for the primitive itself
would thereby be annihilated. But the
contradiction that the independence of the consecutive essence
lies in its non-independence is removed in subordination.
That it has independence follows precisely
from its having in itself the primitive as its essential
principle; the independence it has in itself
is thus not its own; it is sacrificed in being subordinated.
But since, as a semblance reflected in itself,
it reflects, has and is in itself its own essential
distinction, it is not violated in its sacrifice
by an alien essence: its subordinating is
its own complete self-position, and its independence results uninterruptedly from its non-independence.

The primitive essence can then subordinate the
consecutive to itself only in that the latter subordinates
itself. The self-position of the consecutive is thus
the basis on which the independence of the primitive
rests, the real presupposition from which
its whole primitivity proceeds. But in this way
the primitive essence itself becomes consecutive, and the
consecutive primitive. The presupposing of subordination
transforms itself, then, in its development into a
mere interchanging of the independent and non-independent,
thus into a coordination of both. The
essential advance lies in this, that all the reflection-movements
and thetic moments that the
% --- p. 124 ---
\setcounter{footnote}{0}%
\opage{124}central opposition contains are now grasped and determined;
but since each of the terms in this opposition
is itself at once both the primitive
and the consecutive essence, and the one is thus in the ground
the very same as the other, the
whole opposition is itself a central contradiction, and
must consequently go to the ground.\footnote{When the central opposition which logic has hereby intimated between the essence of God and that of the world cannot maintain itself, but dialectic must again let them vanish into each other, the ground of this is that the categories in which God and the world are here apprehended are still too abstract. But the theistic type vindicates its significance nonetheless; for every central opposition that vanishes thereby forms a transition to fuller categories, in which it returns with greater clarity. The typical course of development is then this: the deeper the immanence, the higher the transcendence; the more concrete the universality, the firmer the centrality. Speculative logic therefore cannot stop before the theistic central opposition is posited with absolute freedom, so that the primitive essence can no longer vanish in the consecutive.}

\underafsnit{c) Absolute Essence.}
\paragraf{23}{Primal Ground, Content, Mediation.}

In absolute essence every determinately existent
independence is sublated; all is in the ground one; thus is
\emph{primal ground}. But distinction is just as
fundamental as identity. All is in the ground distinct; as the ground of all, the primal ground opens itself out in a \emph{content} of relative grounds. This double
% --- p. 125 ---
\setcounter{footnote}{0}%
\opage{125}ground-movement, whereby the primal ground partly sublates ground-elements, in order itself to ground its own
transcendence, partly gives itself up, in order to express its essential immanence in
the determinate relative grounds of determinate existence, is consummated in \emph{mediation}. In ``pure essence'' all being-determinations are absorbed
as reflexes; the one reflection-moment mirrors the
other; in ``the duality of essence'' the reflection in % print: Dualitet„ (closing quote set low) (misprint)
other becomes determined on all sides as a reflection-in-itself. Essence
sublates semblance, and is reflected in itself as essence in essence,
just as semblance conversely takes up essentiality
into its self-reflection. This central opposition
between the primitive and consecutive essence is reflected
again, through subordination, in one single ground-identity,
whereby the immanent, all-sided reflection is carried through,
and essence posits itself as \textbf{absolute essence}.
In absolute essence every essentiality-moment now steps forth
with absolute reflection-in-itself. The primitive
essence is absolute, for it sublates the consecutive as
a moment in its own self-reflection, as substrate for
its subsistence; each of its hypostases is consequently
also absolute, for it reflects the whole primitive essence
in itself. On the other hand, the consecutive
essence, which in its self-movement sublates the primitive,
and, when its reflection is complete, subordinates it
to itself, is likewise absolute; and since now each of the
consecutive hypostases again reflects all the others,
the absolute itself is posited in every moment of
absolute essence. Since no being, no semblance can
fall outside the absolute, the absolute in turn cannot
fall outside any of its moments whatsoever; every
moment of the absolute consequently becomes itself posited
as an absolutely independent, indeed as the absolute itself.
% --- p. 126 ---
\setcounter{footnote}{0}%
\opage{126}But the absolutely independent tolerates no otherness
outside itself. For if it left otherness behind,
this would in turn, by virtue of the absolute, itself become
an absolutely independent, and we would thus obtain two
absolutes, of which neither was absolute. As a moment
of the absolute becomes absolutely independent, it thereby
subordinates all the others, and transforms them into substrates.
Now as all the moments have equal claim to
absolute independence, so must they also all, without
exception, be transformed into substrates, and are consequently,
each for itself, absolutely non-independent. Here, then, is an
innumerability of independents, which are sublated and vanish
in the absolute independence.

When the independent is sublated, it goes to the ground.
Independent being is first the true determinate being;
to lose the independence of one's being is to lose one's proper
being; the transition from absolute independence to
absolute non-independence is therefore a transition as
immediate and abrupt as from being to non-being:
to \textbf{go to the ground is to be annihilated}. But the
absolutely independent is sublated only because it
itself has its non-independence in its independence, because the other
which constitutes its foundation and substrate would itself
be posited as absolutely independent; it thus goes to the ground
only insofar as it seeks its independent foundation;
in this respect it is so far from being annihilated
that by \textbf{going under} it much rather itself becomes
the solid \textbf{underlay} and foundation in which the
independent other must posit its true independence;
\textbf{there is thus no} \textbf{annihilation}. But if now
the whole multiplicity of independents is reduced to substrate
for the independence of the absolute, then this becomes an
independence in which nothing is independent, consequently an
absolute non-independence; if, conversely, the absolute
% --- p. 127 ---
\setcounter{footnote}{0}%
\opage{127}independence goes up into the multiplicity
and distinction of the independents, it precisely thereby goes under in sheer
relativities, and neither have we anything absolutely independent.
In the ground itself, then, lie the immediate identity and distinction of independence and
non-independence. In order that independence can be grounded, the independents
must go under; but in order to ground the independents,
independence itself goes to the ground. As
all is thus to be grounded, nothing becomes grounded.

This ground-contradiction in absolute essence can
first be removed by the true ground-movement. The absolute
grounds itself in order to ground its other, and
it grounds its other in grounding itself:
the absolute is absolute only in that it itself is the medium
between itself and other. The ground-movement is then,
as a consequence of this, a double one: partly a primitive, whereby
the absolute grounds itself as absolute; partly
a \textbf{consecutive}, whereby it lays itself down as ground for the
non-absolute. The one of these ground-movements grounds
in turn the other, and is thus the other's medium: grounding is \textbf{mediation}.

As the absolute independence in which all that is independent
has its ground, absolute essence is the
primitive ground, i.e.\ \textbf{primal ground}. Since pure absolute
independence sublates every independent being,
it is, in this its abstract being-for-itself, presupposed as
the negation of all independence, as that in which
as yet nothing is grounded, but from which all being, even
the absolute, is nevertheless to proceed and be grounded. In this
its transcendence the primal ground is thus determined as the
real nothing, \textbf{the negative identity of} all. The
non-absolute proceeds from its other, and precisely
therefore comes from the absolute as from its nothing; the absolute
itself is its own other, it originates from itself, and is
% --- p. 128 ---
\setcounter{footnote}{0}%
\opage{128}precisely therefore the unoriginated, i.e.\ that which has originated from nothing.
Nothing is thus that from which all originates, but only the
real nothing grounds a true, real being. In ``becoming''
pure being does indeed originate from pure nothing,
but here there is no grounding; for since being, as such,
is only an abstract immediacy, the nothing that is
is also an empty, abstract negation. In reflection,
essence is the nothing of semblance and semblance the nothing of essence,
but they alternate as fleeting reflexes in the clear light, and
lack the being of the ground. In mediation, on the other hand,
where every being is reflex, and every reflex a being,
where becoming is reflection, and reflection becoming,
real being goes forth from real nothing as from
its ground.

The \textbf{transcendent} primal ground, which concentrates all-being
in negative identity, is consequently no immediate
being-indifference which itself immediately and without any
ground turns over into difference; for its posited reflection
is its solid being-in-itself, whereby it holds fast
all in absolute unity; but neither is its being-in-itself
an abstract essence-reflection in which the one hypostasis
clarifies the other through and through; in the real nothing all
being collapses \textbf{immediately}, and the light of reflection
is extinguished in the night of the primal ground. As reflection goes
to the ground, it is sublated in becoming; as the real nothing,
the primal ground is therefore the immediate identity of the principle and the beginning of all-being. For insofar as it
is the negative presupposition of real being, the absolute
other from which the absolute itself is to proceed, it is
principle; insofar, on the other hand, as it has as yet grounded
no being, since it is itself nothing, it finds
itself in a becoming, and must therefore be posited as the
absolute beginning. With the absolute beginning
all begins; every distinction between principle and
% --- p. 129 ---
\setcounter{footnote}{0}%
\opage{129}result as between the primitive and consecutive, every
separation between the different terms of successive development
must also have its ground in the primal ground;
therefore all will come to be at one stroke. But as all is
at once to come to be and be grounded, and the grounding of the one
presupposes the going-under of the other, since it
presupposes its foundation, the transcendent primal ground,
as negative identity, i.e.\ as unity of conflicting tendencies,
is in an infinite tension, unrest and fermentation.\footnote{Jacob Böhme has depicted this dialectical self-contradiction of the primal ground in mystical colors in his writing \textit{De tribus principiis}: ``Man findet im Urkund die allerschreklichste und strengeste Geburt, alles Herbe, Bittere und Feuer: da kann man nicht sagen, daß es Gott sey, und ist doch der innerlich erste Quell, der in Gott dem Vater ist, nach welchem er sich einen zornigen, eifrigen Gott nennet, und derselbe Quell ist das erste \textit{principium}, und ist Gott der Vater in seinem Urkund, daraus diese Welt sich urkundet. — In diesem \textit{principio} stehet nichts als nur die allerschrecklichste Gebährung, die größte Aengstlichkeit, feindliche Wonne, gleich einem Schwefelgeist, und ist eben der Höllen Porten und Abgrund, darinnen Fürst Lucifer in Verlöschung seines Lichts geblieben''.}

This strife between nothing and being, which in the ground
is a strife between foundation and grounding, dissolves
the abstract transcendence of the primal ground. The ground
lays all under itself in order to be all; its non-independence
goes under in its absolute independence; it
makes itself the foundation for its own grounding.
Since the becoming of the ground thus consists in this, that,
according to the primitive ground-movement, it sublates the real nothing
in the real being, and posits itself as its own
consequence, it thereby becomes determined as the \textbf{positive}
% --- p. 130 ---
\setcounter{footnote}{0}%
\opage{130}\textbf{identity of all}, i.e.\ as the \textbf{transcendental
primal ground}.

Insofar as the transcendental primal ground presupposes
its real nothing, in presupposing a ground-movement from which it results, there is in the ground nothing, whose
immediate being is indeed sublated, nothing, which has indeed been
made foundation; but insofar, on the other hand, as it is the
transcendent primal ground that has itself become transcendental,
there is nothing either, whose real being is indeed posited in
and with the ground, nothing, which has indeed been grounded.
If the unity of the transcendent and transcendental
primal ground were immediate, there would be no ground-movement,
and consequently neither ground nor grounding.
The primal ground itself thus presupposes a ground-distinction, and
so rests in two mutually opposed foundations.
But if, conversely, the negative and positive identity of the primal ground
were not one and the very same ground-identity,
all ground-relation would cease, and nothing would be
grounded either. Since the true primal ground can first result
from the presupposed opposition, there is no
ground-relation that can connect the two primitive
foundations; they therefore relate indifferently to each other,
and are, each for itself, ungrounded, since they themselves
constitute the presupposition of the ground-movement. In the duality of the
principle, then, lies the absolute beginning.
But since the primitive ground itself results from the
dualistic principle, the two foundations contain only
the ground of the primal ground, and consequently have identity and the
mutual ground-relation in themselves; every foundation is thus
itself the medium between the ground and the grounding.
According to this development the identity of the negative
and positive identity is mediated, and the primitive ground
determined as the \textbf{absolute} primal ground.
% --- p. 131 ---
\setcounter{footnote}{0}%
\opage{131}The negative ground-identity, the real nothing, now
constitutes the whole absolute primal ground, but only on
the presupposition that the positive identity, the real
being, is at once posited both as its posited
other and also, according to the grounding, as its sublated
other, i.e.\ as its foundation; likewise the positive
ground-identity, the real being, will now also
constitute the absolute primal ground itself, which has overcome
and sublated the real nothing. The primitive ground is
thus in itself neither an abstract essentiality nor
an abstract determinate being, neither a real nothing nor a
merely real being, but the medium of both. The posited
being thus, according to the carried-through mediation,
comes no more from nothing than nothing from being; but
\textbf{everything must have its} \textbf{sufficient ground}.

According to the primitive ground-movement every distinction
and opposition has vanished as ground-element in
the identity of the primal ground; all is then in the ground one; the
primitive, absolute ground is the only ground of all.\footnote{Devotion, as it is described by \textit{Dr.} W. H. Rothe (The Doctrine of the Trinity and Reconciliation): ``He (man) is gone, the soul hovers in a sea of light, gentle feelings stream through the heart, clear figures of life pass by the soul; in visions, fleeting as morning clouds, in the depths of the heart his spirit moves, but not as his etc.''\ expresses the sinking down of subjective consciousness into the primal ground of life.} From this there now develops a double ground-opposition.
In the first place, the primal ground is opposed to itself; it is its
own other and as such itself both ground and consequence; next, it is opposed to its other, for according to its self-grounding
it has as yet become merely ground; but a
ground that grounds nothing \textbf{is} no ground. The
% --- p. 132 ---
\setcounter{footnote}{0}%
\opage{132}primitive ground-opposition is thus indeed a fundamental
presupposition for the consecutive one, but without it utterly
ungrounded. Only the ungrounded can be grounded;
for what is immediately identical with the ground itself
cannot follow from any ground-movement whatsoever, and
consequently cannot be grounded either; the other of the ground must
therefore be found present as something that immediately is, which does not
lie in the ground at all, but is utterly distinct from
it. But only that which proceeds from the ground itself
can be grounded, and only that which lies in the ground can
proceed from the ground: the ungrounded must nevertheless lie
in the absolute ground. The primitive ground thus has
in itself that which is distinct from it, and which it
itself is not: it has a content.

The content is, in the first place, the ground's own according to
the primitive ground-movement, identical with the primal ground
itself; its moments then continually coincide with the ground
and thereby with one another. The primal ground
broods over its \textbf{esoteric} content as over an eternal
mystery. But on the other hand the content is
the other of the ground, which emanates from the ground, that
in which the absolute ground goes out beyond itself and
posits its relativity; as the ground-moments of the content
thus separate themselves from the ground itself, they thereby part
from one another, and step forth, each
for itself, with a relative, grounded independence: the
mystery of the ground opens itself out in the \textbf{exoteric} content.
The esoteric and the exoteric content are in the ground
one, and precisely therefore ground-distinct. From the fundamental proposition
that all is one is now derived this one, \textbf{that all} in
the ground \textbf{is} distinct; from here the consecutive
ground-movement begins.\footnote{It is Professor Martensen's merit, in his ``Mester Eckart'', to have brought forward a cardinal point in mysticism of which, in philosophy as well as in theology, one has hitherto had only a very unclear representation, namely the relation of mystical consciousness to mystery and revelation. Taken logically, we can now locate the one-sidedness of speculative mysticism in this, that in its profound contemplation it uninterruptedly tries to reduce all the oppositions of determinate existence to the esoteric unity of the primal ground, but for that very reason is just as uninterruptedly compelled by the dialectical contradiction to return to the exoteric content.}
% --- p. 133 ---
\setcounter{footnote}{0}%
\opage{133}Insofar as the exoteric content in its distinction from
the absolute ground is nevertheless identical with it, it is
thereby determined as the primal ground's own alteration, i.e.\ as
an innumerability of \textbf{determinately existent grounds}; insofar as it is,
on the other hand, posited as the other of the ground, it represents
itself as a \textbf{grounded determinate existence}. The being-in-itself
of the absolute ground is the presupposition of its being
in other; its transcendence grounds its immanence;
it is thus itself the medium between the determinately existent
grounds and the grounded determinate existence, and mediation
lies in this, that the absolute becomes relative,
in order that the relative can become absolute. Since it is the primal ground
itself that is altered, the absolute has thereby
itself become relative; everything determinately existent has its measure
of the absolute in itself, and thus expresses in its
peculiar delimitation the whole absolute.\footnote{One finds oneself at this standpoint when one seeks the revelation of God's essence in a blade of grass.} But if the relative is to express the absolute,
it must itself be posited with absolute relativity. As an isolated,
for-itself something, the determinately existent has merely a
relative ground; only on the presupposition that it
is itself a moment in all the rest of determinate existence, in which
the altered primal ground exhausts its whole
% --- p. 134 ---
\setcounter{footnote}{0}%
\opage{134}multiplicity of relative grounds, can it be said to have the absolute
ground in itself. Thus: the determinately existent has its
ground in itself, and is immediately identical with it, for it
is because it is; but just as immediately it is now
also distinct from the ground; for it exists only on
the presupposition of another determinate existence, and consequently
has its ground \textbf{outside} itself. As the ground of
determinate existence, the altered primal ground is immanent in determinate existence
itself in that the one determinately existent grounds
the other. Everything determinately existent is thus at once
both ground and consequence. As consequence it arises
through the going-under of an other; as ground, on the other hand, it
itself goes under in order that an other can arise. But
also only insofar as determinate existence expresses the primal ground's
own alteration can the one determinately existent ground the
other. If everything determinately existent merely had the ground
immediately in itself, the multiplicity of determinate existence would
express the multiplicity of the determinately existent grounds of the altered primal ground;
the whole of determinate existence would then be
grounded at one stroke, and the ground-movement would cease;
the distinction of the ground from itself would then become so immediate
that the ground-identity, and with it the ground-relation, would entirely
vanish. Only in that the one determinately existent
uninterruptedly passes over into the other, as the ground into the
grounded, can the ground-movement be vindicated, and identity
uninterruptedly result from distinction. But if, conversely,
the one determinately existent merely had its ground in another
determinately existent, the one would have to presuppose the
other in an infinite series; distinction would then be
posited so immediately in determinate existence itself that the ground-movement
would lose itself in an uninterrupted change without grounding.
The presupposition that the one determinately existent
grounds the other then turns over just as immediately
into the presupposition that everything determinately existent is immediately
% --- p. 135 ---
\setcounter{footnote}{0}%
\opage{135}grounded by the primal ground.\footnote{Traducianism, which seeks the ground of the children's souls solely in the souls of the parents, is, e.g., just as one-sided as creationism, which seeks the ground of every soul immediately in God's creative will.} This abstract contradiction
is mediated in the consecutive ground-movement in that
\textbf{the ground-relations of determinate existence are specified}.

From the standpoint of immediate being the distinction of determinate existence
shows itself as an outward, superficial diversity;
here there is found only the indeterminate swarm of
``somes'' and ``others,'' which come and vanish in indifferent
becoming and change, and in their separation and
connection ceaselessly seek a highest goal. But
as being first finds its measure in essence, so it is
in turn the determination of essence to make itself the ground of
being and to ground determinate existence. Regarded from the standpoint of
absolute essence, the distinctions of determinate existence therefore present
themselves as essential and grounded;
the ground posits the determinate measure, and modality can
now no longer be dissolved into pure indifference. If
any something whatsoever could be laid as ground for any
other whatsoever, unity itself would be groundless,
and every fundamental distinction would vanish. It is then,
as a consequence of this, only certain particular determinately existents that
can be posited as ground-elements for certain others, and
become presuppositions for their grounding; with
essential measure, elective combination thus also returns,
and determinate existence is grouped, according to its ground-content,
in closed circles. As now the one and same
ground-relation which runs through all determinate existence must
specify itself into particular relations in order to become grounding,
the primal ground is mediated with every determinately existent
% --- p. 136 ---
\setcounter{footnote}{0}%
\opage{136}ground, and the contradiction is sublated in the self-determination of the ground. In the one determinately existent lies the determinate ground
of another determinately existent on the presupposition
that every ``something'' just as immediately refers
itself to a circle of ``others'' as to itself.
Insofar as the determinately existent has the ground in itself, it refers
itself to itself alone, and according to
this ground-relation cannot immediately contain
the ground of other; but insofar as it has its
ground outside itself, it has, on the other hand, an overreaching
content, a sum of determinations and
relations whereby it refers to certain other
determinately existents, as a something that is itself posited \textbf{because
of the circumstances}. The content which
everything determinately existent thus has in itself, and the content
it has outside itself, are then content-moments
of one and the same ground-relation, and in the mediation
of these moments lies the specific
ground. α) The ground is the identity, mediated with itself,
of the grounding and the grounded; the
determinately existent which is to ground, and that which is to
be grounded, have, as a consequence of this, the same content;
for otherness itself lies in the determinately existent
which contains the ground, since its relation to
other is only a moment in its relation to
itself. β) But the real ground lies, on the
other hand, in the fundamental distinction between the
grounding and grounded; the determinately existent which
is to ground must therefore conversely posit its relation
to itself as a moment for the relation
to other; the pure mediation of the ground then falls
outside the ground-elements through which it
is accomplished: every specific grounding
% --- p. 137 ---
\setcounter{footnote}{0}%
\opage{137}\textbf{presupposes its particular conditions}. γ) The one determinately existent thus contains the ground of the
other in that there lies a whole sum of
conditions between the ground and that which is to be grounded.
The conditions must be presupposed as immediately
given; they constitute only the discrete points
in the circle of determinate existence within whose periphery
the grounding is sought; as \textbf{one} such sum of
determinately existents they relate indifferently to the whole
ground-relation, and give themselves out as \textbf{unconditioned}. But
in their unconditioned being they are merely determinately existent,
and can consequently condition nothing, since their relation
to themselves among one another vanishes in abstract
immediacy: only in relation to the ground-relation
does their immediate being-for-itself turn over
into a being for other; for another ground-relation
the very same moments of determinate existence would
furnish other conditions; they are thus, as conditions of the ground,
themselves \textbf{conditioned}. If now
all conditions thus depend on the determinate ground-relation,
this must thereby step forth as unconditioned;
by a mere addition of conditions one procures
no grounding; here one must then begin
with a determinately existent that can be said to contain
the ground. But only insofar as the content of the conditions
is taken up as moments into the content of the determinately existent
ground can a real
distinction be posited between the grounding and that which is to be grounded; the peculiarity of the ground-relation is
thus nevertheless conditioned by the immediate determinate existence
of the conditions. Through the immanent unity of the pure mediation of the ground
and the immediacy of the conditions, the ground
becomes specified and determined. But the ground
% --- p. 138 ---
\setcounter{footnote}{0}%
\opage{138}so determined is in itself insufficient. For the
closed circle of the particular conditions is
only a term in the infinite series of all the other
circles of determinate existence, and the specific ground-relation,
which through its conditions is unconditionally brought out
as a for-itself determinateness, only a
moment in the universal ground-movement. In order
now to be freed from this indeterminate infinity,
the determinate ground must then at the same time be posited by an
unconditioned position of the primal ground itself. This
specific position, whereby the primal ground fixes the
relative ground over against the particular conditions,
in order to sublate both as moments in the mediation
of the sufficiently determined ground, now
itself bears the stamp of immediacy. According to this
immediacy the unconditioned position of the primal ground is
then also conditioned; for its specification has
the relative ground and its conditions as its
outer presupposition, whereby its relation to itself
vanishes in a relation to other. But the
relative grounds and their conditions, which constitute
the whole grounded determinate existence, are only an immediate
expression of the primal ground's own alteration;
the presupposition by which its specific position
is conditioned it has itself presupposed; it is thus conditioned
only by the conditions it has itself given itself,
and thus steps forth in the sufficiently determined
ground as \textbf{the absolutely unconditioned}. In the absolutely unconditioned all ground-determinations are now
posited, and thereby at the same time sublated. Determinate existence
with its relative grounds is only presupposition
and condition for the primal ground's own self-grounding;
therefore the primitive ground sublates the
presuppositions it has itself posited. In the % print: Forudsætuinger (misprint)
% --- p. 139 ---
\setcounter{footnote}{0}%
\opage{139}immediacy of determinate existence its content is altered, and this
alteration is a ferment whereby it consummates
its own real mediation, and lets all conditions
vanish as moments in the esoteric content.
As determinate being rises up out of the primal ground, so
must it also sink back into it; therefore
the determinately existent ceaselessly goes to the ground; the relative
grounds are in themselves only semblance-grounds. But
conversely the primal ground now also sublates itself, in
sublating its exoteric content; for its ground-identity
is itself condition and presupposition
for the development of the fundamental distinction.
In the outer conditions and the immediacy of the determinately existent grounds
lie the alteration and self-sublation of the primal ground; but as it must thus in turn
sublate its own self-sublation in order to posit itself,
the negation of the conditions shows itself as
an indispensable condition for the unconditioned
being-in-itself of the primal ground, for its esoteric content. Since
the exoteric content is posited by the esoteric, it
nevertheless always lies, in the ground, outside the ground;
it has an alien ground in itself, is a merely grounded,
a substrate for the true ground's mediation
with itself, and must therefore be sublated.
But as the outer multiplicity is sublated as
conditioned and posited, it precisely thereby itself becomes a
positing, and transforms itself into a real ground for
the esoteric content, which now becomes a posited and
grounded. As posited and grounded, the esoteric
content must in turn be sublated, and in its self-sublation
lies the grounding of the exoteric. The semblance-grounds
of determinate existence are all \textbf{grounded semblances};
they are \textbf{phenomena}.
% --- p. 140 ---
\setcounter{footnote}{0}%
\opage{140}Although now this mediation, whereby the transcendence
and being-for-itself of the primal ground uninterruptedly posits itself
as presupposition for its immanence and being
for other, can in the ground never cease, it is
nevertheless meaningless to wish further to hold fast
the separation of the opposed ground-determinations, since
they have all shown themselves to be one-sided abstractions.
The immediacy regained by the sublation of the ground then steps forth with unconditioned independence.
Since mediation is thus exhausted,
identity again turns immediately over into
distinction, and the latter into the former; the ground itself thereby steps into the background, and a new sphere opens.

\overgang

The identity of essence cannot be without distinction. For
pure essence, distinction is semblance; for the primitivity of essence and
ground, it is a consecutive essence, a
determinate existence of relative grounds. But, as a moment vanishing in identity,
distinction has not yet come
into its right; therefore identity must now conversely
be laid as ground for an independent development of phenomenal
oppositions, and reflection itself change
its character. For while the immanent, dialectical-speculative
reflection predominates in essentiality,
the outwardly distinguishing and fixing reflection
becomes, on the other hand, predominant in
B. The Sphere of the Phenomenon.

\afsnit{B. The Sphere of the Phenomenon.}
\underafsnit{a) The Ground-Phenomenon.}
\paragraf{24}{Matter, Form, Reflexivity.}

As the phenomenon, the immediate unity of
essence and semblance, is posited as result of the mediation of absolute
% --- p. 141 ---
\setcounter{footnote}{0}%
\opage{141}essence, the immediate falls outside the mediate. That which
immediately is in the phenomenon represents itself as \emph{matter}, the reflection falling outside matter, on the other hand, as form. This distinction is again just as immediately sublated: no matter without form --- as immediately
fixed: the same matter can have the most diverse
forms, and the same form contain the most diverse
matters. This outer reflection-relation is expressed in the
\emph{reflexivity} of matter and form.

The dialectical transition from being to essence is
an apparent advance, but an essential retreat. According to its pure immediacy, determinate existence represents
itself as if it were groundless only because it
needed no ground, as if it were itself the
unconditioned presupposition for all movement and development,
thus as the primitive. The essencelessness in being
itself is thus, regarded from this side, so far
from being a lack that it much rather expresses
the apriority whereby determinate existence becomes able
to posit essence. The movement from being to essence
is thus, according to this view, progressive: essence
\textbf{is a consequence of} \textbf{being}. But as the
self-development of being ends in essence, it becomes evident that the
unrest and movement which has driven being out beyond
itself stems precisely from its own lack and one-sidedness;
its transition into essence testifies to its origin
from essence. Since being thus shows itself as
the immediacy which essence itself has presupposed, in order by
the sublation of this presupposition to posit itself,
the movement from being to essence is \textbf{regressive}.
% --- p. 142 ---
\setcounter{footnote}{0}%
\opage{142}Essence, which comes forth as the last, \textbf{is} then the
proper first: \textbf{being is a consequence} \textbf{of essence}.

The contradiction that essence cannot come forth except
as a consequence of being, although being itself first comes
to be as a consequence of essence, finds its solution in
the mediation of the ground.

Insofar as the esoteric content of the ground presupposes
an exoteric one as that which is ceaselessly posited and is to be posited,
the ground itself follows from immediate being;
but as it is centralized into primal ground, it sublates
the immediacy of all-being in its own absolute independence,
thereby follows from its own consequence, and reduces
determinate existence to a consequence of the ground. The advance of being
to essence is then hereby posited as a retreat
to its proper source, and the ground has sublated
the illusion that essence should be a consequence
of immediate being. Since, however, in every true
illusion there nevertheless lies a semblance of truth,
this truth-semblance in the ground now also comes into
its right. For it has turned out that the esoteric
ground-content in itself would only be a semblance if
the exoteric with its relative grounds and conditions
were only semblance; if the one of these main sides vanishes,
the other vanishes with it. As it is thus
an illusion to wish to hold fast a for-itself
primal ground, and to let immediate determinate existence come
and vanish as superficial semblance, or, on the
other hand, to wish to dissolve the primal ground centralized in itself into
the relative grounds and their conditions, so it is
also a logical illusion if one\footnote{As Hegel. Werke 4. B. p. 110—118.} lets both these
main sides go under in an immediacy that is.
% --- p. 143 ---
\setcounter{footnote}{0}%
\opage{143}If the absolutely unconditioned is to be determined as that \textbf{which} sublates the primal ground and its conditions in an identity of essence and semblance
that is, then the advance of the ground-movement
to ``the fact itself'' again becomes a retreat
to the absolute primitivity of immediacy. This
indissoluble, adamantine immediacy then reduces the
preceding mediation to a merely outward reflection,
a subjective thought-movement. All thus ends in a % print: Tankeb.vægelse (misprint)
sterile unconditionedness from which one can come neither
forward nor back. The absolutely unconditioned cannot be
held fast as \textbf{posited} without at the same time presupposing the
presupposing ground; it cannot gather itself together and
sublate ground and conditions without at the same time having to
let go of the primal ground as that which can posit other conditions;
it cannot step forth as determinate without
reflecting the ground-movement by which it is determined.
``The fact'' is thus this: that the inexhaustible primal ground
always yields determinate grounds, that the stream of mediation
breaks its waves on the immediacy in which
it is itself crystallized. If the content of the ground is to be
more closely determined, the determination can no longer
be kept in the pure streaming of mediation, but must be posited
and sought in the medium in which mediation has gone
together with immediacy, where the identity of essence
has become identical with the distinction of semblance, where the
determinate ground has been absorbed into the grounded, thus
in \textbf{a being}, \textbf{which so follows from} \textbf{essence}
\textbf{that essence and the ground} \textbf{themselves follow along}.

The phenomenon is now this immediate, and thereby at the same time
determinate, identity of essence and semblance. As distinction
is sublated, reflection has vanished. Essence
is no longer essence, because its separation from
semblance has ceased; for the same reason semblance is no longer
semblance either; both reflexes hereby take on \textbf{one}
% --- p. 144 ---
\setcounter{footnote}{0}%
\opage{144}stamp of immediate being. But the immediacy of identity
is won through the mediation of distinction and opposition;
semblance therefore presses forward
anew, gives its ground-content a tinge of essentiality,
and reflection thus returns again, only
with changed character. Essence is determined by semblance,
and has posited its content in the reflection of semblance, and
semblance is determined immediately as that which gives
essence its determinate shape. Essence determined by semblance
now no longer keeps itself in a distant
background, but is itself the semblance in which it shines;
conversely, semblance is no longer the empty reflex either,
which vanishes and deceives, but represents the fundamental
essence: the reflection of the phenomenon thus reflects its
own sublation. As immediately as the determined
here coincides with the determination, so immediately
do they again separate from each other. This alternation of
identity and distinction, according to which that which is to be determined
just as immediately has the determination in itself
as outside itself, characterizes the phenomenon as that which is grounded
in itself, as something that is, which is at once in
the ground and yet outside the ground, as a ground-phenomenon, whose universal moments are \textbf{matter} and
\textbf{form}.

Matter is the essence that is, the fixed presupposition
of essentiality, which approaches palpable
immediacy. In opposition to this, form shows itself
as the pure activity of determinability, reflected in itself.
Matter and form thus seem to relate indifferently
toward each other.

As essential immediacy, matter is indifferent
toward reflection and form; therefore it represents
itself in the sphere of the phenomenon as a substrate that takes on
all the determinations of being.
% --- p. 145 ---
\setcounter{footnote}{0}%
\opage{145}That matter is that which is in the ground-phenomenon
is seen also from its reproducing at once both quality and quantity. Its ground-parts,
the atoms, express the original quality-distinction, and
its infinite divisibility gives us back quantitative
discreteness. But matter does not come to be, although it \textbf{is}
in a constant becoming; it has unchangeability in common
with essence itself, although it is uninterruptedly subject
to change. Since material identity is just as
little able to posit its own distinction as
material distinction is capable of reflecting its identity,
matter, as such, remains \textbf{inert} and
obscure. Matter is in itself indeterminable, and must
nevertheless be determined; its obscurity negates reflection,
and nevertheless presupposes the light of reflection: matter's own
indeterminability thus expresses its claim to
be determined by its other, i.e.\ by form.

Form, taken for itself, is pure determining, which
has as yet determined nothing, since that which it is to determine
falls outside the determination itself; it is the \textbf{clear},
\textbf{freely} hovering \textbf{reflection}, which in the being of matter
reflects its own self-sublation. Formal identity
mirrors its distinction; it thereby points out beyond
itself, comes into movement, and determines its own
sublation; for the same reason formal
distinction reflects identity as its negation and sublates
itself. Form is thereby posited as the active principle,
as infinite self-movement.

In their immediate opposition matter and
form are, each for itself, equally finite and non-independent.
If matter could not determine itself because its
identity and distinction fell indifferently apart from each other,
neither can pure abstract form consummate
its self-determination, since its distinction and identity
% --- p. 146 ---
\setcounter{footnote}{0}%
\opage{146}uninterruptedly sublate each other. The unity of material
and formal distinction is first the true determinate
distinction; for insofar as distinction has a material side,
it cannot be sublated as a semblance vanishing in identity,
and insofar as it at the same time
has a formal side, it points out beyond itself, and determines
itself to identity. As matter without form
is a shadow, so form outside matter is a
semblance; only through their mutual affirmation do they
both obtain independence. Matter's negation of all form
is its own essential form. By keeping itself as
the indeterminate over against the determining form, it gives
itself a determinateness. Its indifference toward form
is its own indifference toward itself; when
it then takes up form and lets itself be determined, it takes up
no more than what it already has; form expresses
only matter's own independence. On the
other hand form, when it keeps itself in its
independent distinction over against matter, becomes reflected in
its otherness. But now form-identity is precisely determined
as that which in all distinction and opposition wins
itself; matter is thereby posited as the other sublated
in the negativity of form, i.e.\ as form's own independent
subsistence.

This developed identity of matter and form shows
that the one does not ground the other; for since there
is no transition, there can be no ground-relation
either. As in the ground itself there is no matter,
so in matter there is no ground either, and
as form cannot follow from matter, since it
is itself matter's own immediate presupposition, so
matter cannot follow from form either, since abstract
form presupposes the being of matter. The subsistence of matter
and form in and through their phenomenal
% --- p. 147 ---
\setcounter{footnote}{0}%
\opage{147}identity is just as immediately a subsistence in phenomenal distinction. Matter, which in its independence has sublated
form, is thereby already formed, just as form,
which has posited matter as its sublated other, is
materialized. Both are then at one stroke in the ground-phenomenon
itself.

If the phenomenal identity is more closely determined, it
shows itself at once both as essential and
as unessential. It is essential to matter to have
a determining form; but which determinate form it
is to have is as yet quite unessential to it; likewise
it is essential to form to be materialized; but which
matter form is to determine is entirely indifferent,
since the determination of matter depends solely and alone
on form. Insofar now as matter is reflected in
form as in its negation, it thereby becomes delimited;
but insofar as it itself subsists in form, it
just as immediately sublates its non-being in being,
and is posited as \textbf{determinate} matter, whereby form
is then also fixed into immediately \textbf{determinate} form. This
determinateness, with which all mediation has ceased, lies in
the \textbf{surface}, which precisely shows itself as the pure identity
and distinction of matter and form. In superficial
determinateness matter coincides with its form; for
immediate form is the surface of matter; in the surface
matter ceases, and vanishes as semblance in the being
of form; as the determinateness of matter, the surface is
matter's negative, its pure otherness, i.e.\ form itself. But on the other hand superficial
form must now also vanish as reflex in the being of
matter. At the standpoint of being, all determinateness is
only a limit which lets the delimiting fall
outside the delimited; but the phenomenon, whose determinateness
has through the reflection of essentiality been raised to
% --- p. 148 ---
\setcounter{footnote}{0}%
\opage{148}the immediacy of being, on the other hand identifies the sides of the
limit, and thereby posits its surface for itself. The surface
is consequently matter's own determinate being; so that it
is not matter's semblance, but matter itself as that which is,
that shows itself in the surface in the most determinate way:\footnote{The senses of sight and of feeling give us the surface from two opposite sides. Sight properly sees only the form, not the matter; when we nonetheless seem to see the matter itself, this comes from the reflex of matter being sublated in the being of form. Feeling, on the other hand, refers itself immediately to matter, not to form; we only seem to feel the form, because its reflex is sublated in the determinate matter. If sight posits form as a premise from which one infers to matter, feeling conversely posits matter as the premise from which one infers to form. That the inference as such does not come to consciousness stems from reflection being sublated in the phenomenal identity.}
the superficial determinateness of matter and form is
then, as the posited identity of both, immanent and essential.

But since mediation has vanished, superficial
identity just as immediately turns over into distinction,
whereby the posited determinateness shows itself as unessential.
Matter, which is indifferent toward its determinate form,
at the same time falls outside the surface, and thereby shows itself
as an unformed mass which can receive innumerable other
superficial determinations. On the other hand
form, as the \textbf{moving}, \textbf{articulating}
principle, cannot acquiesce in the surface's determinateness as something that is either; for since formal identity reflects its
own distinction, it must, as a consequence of this, reduce its
superficiality fixed in matter to moment and presupposition
for other determinate forms; only by sublating
the posited determinatenesses as reflexes in an innumerability
of form-determinations can form sustain its pure
% --- p. 149 ---
\setcounter{footnote}{0}%
\opage{149}self-determining. Since every determinate surface can thus
be exchanged to infinity with any
other whatsoever, the superficial distinction between matter and
form is posited, and the determinateness of the surface reduced to
a pure unessentiality.

This distinction between matter and form shows itself
still more determinately when it is apprehended in relation to the surface's
own essential and unessential side. For when
matter uninterruptedly tends to remove the surface's
essential determinateness, and as a ``rudis indigestaque moles'' to sink into its chaos, it lies, on the other hand,
in the forming and shaping principle to posit
superficial determinateness as a moment preserved in the immanent formation,
or to convert the surface's
unessentiality into its essentiality. Although the mass
withdraws from the determinate surface into the indeterminable,
it nevertheless cannot escape form, since
the latter itself reaches out beyond every fixed determinateness.
Since form thus first becomes essential and immanent by sublating its superficial distinction,
the mass, in withdrawing from the surface, falls on
every point to the essential form, and is thereby
posited as its \textbf{stuff}. By giving itself as stuff,
matter is in every respect subjected to form. Insofar
as the surface is still indifferent and unessential,
the indeterminate stuff resembles the mass formless in itself;
but since it is form itself which, by sublating its
own superficiality, has penetrated the mass, and masters
the stuff, the essentiality of the surface is hereby
at the same time posited; for the formation of the stuff consists precisely
in this, that the form and the surface are determined. The stuff's
essential form is superficial, since it lies in
the stuff's own elasticity that it cannot immediately
have any fixed determinateness. The
% --- p. 150 ---
\setcounter{footnote}{0}%
\opage{150}superficiality of form is to be posited as the stuff's pure formability; but
already in that matter is determined as stuff, it has,
as such, the determining form in itself; the specific
expression of the determining form is immediate
form-determinateness, surface; consequently superficial
form is also essential.

The all-sided determination of stuff and form develops
more closely thus: α) Stuff and form are immediately
identical; the immediacy of identity is pure
indeterminateness, and this indeterminateness is the ground of the surface or
of determinateness. With grounding, mediation
returns; as identity posits itself as ground,
it sublates its own immediacy, form comes
into movement, and fixes the surface as the result and consequence
of the immanent formation. But when the consequence is
posited, the ground-movement ceases, and the immanent formation
is stamped out on the surface\footnote{In this presupposition all plastic art is grounded.}. The stuff itself is hereby determined
from a double side: partly it becomes, namely,
in itself infinitely determined by the immanent
form, and thereby raised to the \textbf{content of form}, partly
it becomes determined from the side of finitude, that is to say,
delimited by the superficial form. Since now the sides of finitude
and infinity, according to the sublated
reflection, no longer fall outside each other, since the
unessential and superficial form immediately expresses
the essential and immanent, the stuff has, in the ground,
no other content-determination than that which shows
itself on the surface. This unity of both main sides of the formation
is not a tautological ``Einerlei'',
for mediation has stepped between; but just as little
is it to be regarded as the reflex of a fundamental opposition,
for the presupposed mediation is sublated; in the immediate coincidence of both
% --- p. 151 ---
\setcounter{footnote}{0}%
\opage{151}sides lies, then, the \textbf{specific peculiarity} of the determinate
stuff. β) The
one stuff steps forth with its peculiarity over against
the other; as the peculiarity of the stuff is given with
the determinate form, so conversely the specific
determinateness of form is fixed in the stuff. Since the distinction between
the heterogeneous stuffs is posited superficially, the separation
still shows itself only as outward and indifferent; form
seems to obtrude itself on the stuff, to come to it from outside,
in order to make a divorce between stuff and stuff. Nonetheless
superficial form is already presupposed as
the immanent formation's own determinateness; form-separation
therefore begins just as much from within, so that
form itself in every peculiar stuff presupposes
a distinctive content. How every stuff is in itself
determined as the content of form, it can then show only
in relation to another peculiar stuff of which
the same holds. Insofar as the separation is outward,
the relation can only be determined superficially; the surfaces of the stuffs
touch each other. But in the medium of touching
the one stuff immediately coincides with the
other; every touching produces a superficial re-forming,
in sublating the outer distinction; since now the surface
itself is the immanent form's own determinateness,
every superficial re-forming at the same time produces an
inner transformation: the heterogeneous stuffs interpenetrate
one another\footnote{These two modes of view already began to assert themselves in the old Ionian philosophy of nature. The mechanical tendency proceeded from a consideration of the multiplicity that is, and explained the origin of things from the separation and connection of the primal qualities, unchangeable in themselves, i.e.\ from the all-sided touching of the constituents. But by this hypothesis the qualities' own origin itself became just as inexplicable as their superficial sundering and union. The dynamic philosophy, on the other hand, let the qualities' own determinateness vanish as moments in the infinite transformation, without being able to give any sufficient explanation of the essential and phenomenal distinction of things. The one-sidedness of both theories is sublated in the dialectical identity of all-sided touching and immanent transformation.}. Since all
% --- p. 152 ---
\setcounter{footnote}{0}%
\opage{152}interpenetration is only an all-sided touching, material
distinction is hereby vindicated; the one particle falls outside
the other, so that the stuffs are not confounded, and
do not flow together in chaotic indeterminability; but insofar as
the infinite touching is an interpenetration, and
thereby at the same time a transformation, form itself is posited
as the sundering, articulating principle, and materialism
sublated. From this mutual specification of matter and form
it now becomes evident that, as form
by means of matter becomes distinct from itself,
and thereby separated into a multiplicity of finite forms,
so conversely the obscure, indifferent unity of matter becomes
separated and sundered by means of form. The form-unity,
infinite in its pure formation, which would everywhere
raise every opposition into a fleeting reflected gleam of
the light of identity, is bound and fixed in the dark, inert
mass, finds in the stuff it penetrates ceaseless
barriers, and thus falls out from itself into
a series of material forms. In matter's outward,
indifferent difference, then, lies the ground of
finitude. But regarded from the other side, matter
shows itself at the same time as the indifference which effaces
all distinction. As matter and mass, the ground-stuff is
only one, and like itself at all points; form, on the other hand,
as the pure self-separation of essentiality, is the
% --- p. 153 ---
\setcounter{footnote}{0}%
\opage{153}reflection-movement which drives the mass-like chaos over
into distinction, sunders, articulates and shapes it. It is
the one, infinite form that makes matter into the
many particular matters, into \textbf{one} discrete content of
heterogeneous stuffs: the finite determinateness of matter has its
ground in form. Just as matter, then, from its
side is the medium between the determinate form's identity
with itself and distinction from other forms (a medium
which immediately vanishes in the determinateness of form), so
form in turn is the medium between matter's specific
unity with itself and distinction from other matters;
hereby the ground-phenomenon, indeterminate in itself, is sundered
and split into a \textbf{multiplicity of determinate} phenomena.

The determinate form's reflection in itself is its being
in matter, and the determinate matter's being-in-itself its
reflection in form. The reflection of both in each other shows
itself, then, just as immediately as their independent being-for-itself,
since the reflexes from both sides are removed as
vanishing moments; conversely their being-for-itself
thereby at the same time vanishes as the reflex of their being in
each other: they are, each for itself, \textbf{reflexive}. This
unity of relation and reflection is again, when it is
more closely determined, to be regarded as a) \textbf{pure},
b) \textbf{conditioned} and c) \textbf{unconditioned reflexivity}.
a) According to pure reflexivity it is enough for the determinate
matter to have a determinate form; as form is thus bound in matter because of the immediacy of determination, so it is just as
immediately freed to being-for-itself and independence by
virtue of the outwardness of reflection. With regard
to superficial formation, the one and same
% --- p. 154 ---
\setcounter{footnote}{0}%
\opage{154}form can be imparted to utterly heterogeneous matters\footnote{The same form of Jupiter can, e.g., be expressed in a statue of clay, plaster, bronze, marble etc.},
and as regards essential form, the same
matter is able to undergo the most diverse
metamorphoses, insofar, namely, as it is superficial
re-forming that, according to an infinitely all-sided
touching, brings about transformation, and insofar as
the peculiarity of the stuff is solely and alone
\textbf{to be} expressed in the determinateness of form. When the stuff
by a change of form becomes another stuff,
the otherness lies not in matter, which on its
side continues to be like itself, but in form.
The result is that in the same degree as phenomenal
\textbf{identity} has won independence,
\textbf{phenomenal separation and} diversity has also
been \textbf{made independent}\footnote{Ice, water, steam can show themselves as the same stuff in different forms.}. Since the phenomenal separation of matter and form
nevertheless presupposes identity, they themselves determine
their mutual specific determinateness; this specification
comes to appearance in conditioned reflexivity.
For insofar as form itself is both immanent
and superficial, the determinate phenomena are separated
from each other in that their essential continuity
is reflected in an unessential discreteness; phenomenal
distinction is only a semblance-distinction, since the
one phenomenon can by transformation become the
other. But on the other hand essential
continuity nevertheless comes to appearance only in discreteness; the independence of the phenomena depends on their
phenomenal distinction; in this independence
% --- p. 155 ---
\setcounter{footnote}{0}%
\opage{155}the semblance of discreteness is essentially immediate, and thereby
at the same time material; continuity, on the other hand, which cannot
posit itself as determinate phenomenon, remains unessential
and formal. According to the immediacy of discreteness,
phenomenal being-for-itself is posited as
reflection in other phenomena; every form-distinction
is superficial, indifferent and changeable. The discrete
centers remain standing as substrates
whose alternating forms are imparted from outside. But when
every phenomenon fetches its form from outside, and
outside the given phenomenon nothing can be given other
than phenomena, then phenomenon is uninterruptedly determined
by phenomenon; their exclusion is an assimilation,
so that the outwardness according to which the phenomenon
had to fetch its determination from other phenomena
is thus an otherness which lies in the determinate
phenomenon's own determination. Hereby the discrete
form is posited as immanent, so that every phenomenon
can now be said to determine its own form. But
in this way phenomenal being-for-itself is also
again transformed into a being for other, for the phenomenon
determines itself only by determining
other phenomena. The result is thus this:
every phenomenon is determined by other phenomena,
so that it is thereby posited as a stuff indifferent toward its superficial
form; but just as immediately the
phenomenon is determined in and through itself; it is a content
specified in itself, which sets formation its
determinate measure: the mutual non-dependence of form and stuff
and their reflection in each other, or their reflexivity,
is all-sidedly conditioned. The phenomenon is determined only insofar as it comes
to appearance (for what does not come to appearance cannot
be a phenomenon either); but the phenomenon comes
% --- p. 156 ---
\setcounter{footnote}{0}%
\opage{156}to appearance only in the superficial form, in other phenomena;
its ground is dissolved into sheer conditions; the independence of the phenomenon is thus negative;
it shows itself immediately as non-independence. But
this non-independence is in turn negative; every phenomenon
sublates the conditions in its own ground;
the ground itself vanishes in phenomenal independence,
and posits its phenomenon unconditionedly.

Distinction is now driven to its immediate extreme.
Form-separation is material, for it is outward
and immediate; but matter is at the same time all-sidedly
formed, since all distinction has become superficial.
The superficiality of the forms is nothing other than
\textbf{their} material being. In every determinate form
the determinateness of matter is now at the same time posited so immediately
that every change sublates the independence of the phenomenon.
Reflexivity is consequently unconditioned.
The phenomenon can be altered; therein shows itself its
conditioned independence; but it is altered only
when the given conditions are sublated and new ones have entered;
the changed conditions give themselves each
time as moments in the absolute determinateness of the unconditioned.
As the conditions are, so too is the
unconditionedness in which they are sublated, and so too is
the phenomenon determined. Since now there can be
no phenomenal change without changed conditions,
and the changed conditions are always concentrated
in the unconditioned, the transition from phenomenon
to phenomenon is a transition from the unconditioned
to the unconditioned. This transition is formal,
tautological and untrue; distinction vanishes as
a semblance, since the unconditioned is always identical with
itself. But just as immediately it represents
itself as outward and material; the transition happens by
% --- p. 157 ---
\setcounter{footnote}{0}%
\opage{157}a leap, since the one unconditioned \textbf{is} indifferent
toward the other; this material separation vanishes
again as semblance in the essential identity of the unconditioned, and so on to infinity. Since
matter and form, as a consequence of this, can no longer
be separated, they must necessarily lose their
particular significance, and be absorbed into a world of independent, freely emerging \textbf{phenomena}.

\underafsnit{b) Existence.}
\paragraf{25}{Thing, Properties, Appearance.}

As the determinate phenomenon steps out of the inwardness of the ground and comes to \emph{existence}, its reflection-in-itself becomes an equally immediate being-for-itself,
as its reflection-in-other becomes a being-for-other; it thereby comes to subsist as the \emph{thing}, which has its \emph{properties}. The properties
show themselves to be the thing's own existence: the thing
makes its appearance in the properties; conversely, the properties vanish
in the thing's subsistence: the thing's being in its properties is only \emph{apparent}. This existential antinomy between immanence and
transcendence is sublated in \emph{appearance}. The opposition between the primitive and consecutive
essence has now developed into an opposition between

the unity of the ground that is in itself and the manifoldness of the independent
phenomena. This opposition is
absolute; for as it is sublated, it is posited, and each
time it vanishes, it comes to appearance again. In
% --- p. 158 ---
\setcounter{footnote}{0}%
\opage{158}mediation the ground has its independence, but the
mediating ground-movement must always yield \textbf{a} result that has being, in which immediacy returns.
The transition from the sufficient ground to the independent
phenomenon is an interchange of mediation and
phenomenal immediacy which continues to
infinity. Phenomenal independence is grounded
in the sublated mediation; but the sublation of mediation
lies in mediation itself; mediation thus becomes
what overreaches, sublates its own sublation, and
reduces the phenomenon to a moment vanishing in the independence
of the ground. The identity of the independence of ground
and of phenomenon is to be posited as
\textbf{existence}.

The existential unity of ground and phenomenon
is in the first place immediate, and as such phenomenal.
When the ground has separated itself, as the absolute
primal ground being-in-itself, from the determinately existent and relative
ground, when in this self-separation it has
presupposed and sublated its determinate amount of outer conditions,
then it joins all its extremes together and
fixes itself as the sufficiently determined ground. The
sufficient ground is, in the sum total of its sublated
conditions, unconditioned; its ground-elements are sundered into
particular determinateness, in order to reflect one another in mutual subsistence and negative
being-for-itself; the sufficient
ground thereby itself comes to appearance in striking immediacy:
\textbf{it exists as the} \textbf{independent} phenomenon.

On the other hand, the immediacy with which
the matter and form of the ground-phenomenon fall outside
one another is only a reflex of the overreaching
mediation of the ground. Although matter seems to be indifferent
to all the determinations of form, this
% --- p. 159 ---
\setcounter{footnote}{0}%
\opage{159}indifference is nevertheless only an essential unessentiality. As
matter is not of itself, so it cannot
be in itself either; it must become the foundation for \textbf{the} determining
form, and thereby shows itself to be a mediation-moment
that belongs to the ground. The absolute
primal ground itself needs the manifoldness of the phenomena
in order to ground its own unity therein; it presupposes
the whole superficial play of heterogeneous stuffs and material
forms in order therein to clarify its own immanent,
affirmative independence: the phenomenon exists
\textbf{therefore as a reflected semblance} \textbf{of the existing
ground}.

Ground and phenomenon thus have, in consequence of this,
the same content; the distinction is only formal: \textbf{all}
moments in the absolute are to be determined as ground-moments,
for there is no phenomenal determinateness
that is not posited by the ground, just as little
as there is any ground-determination that must not also
show itself as phenomenal. If absolute
existence is considered from the principal side of mediation, then
the absolute itself exists as the esoteric content of the primal ground,
in which the phenomenal oppositions vanish
as immanent reflexes; if absolute existence
is determined, on the contrary, from the other principal side as sublated
mediation, as reflected being, then the absolute
itself exists in the exoteric content of the ground, in the outer manifoldness
of the independent phenomena, so that
the primitive essence-fullness is absorbed in the consecutive
essence.

But in this existential identity lies an existential
distinction, according to which each of the principal sides itself
is to be posited as absolute existence.
% --- p. 160 ---
\setcounter{footnote}{0}%
\opage{160}1) \textbf{The} absolute exists as phenomenon.

In the independent phenomenon the sufficient
ground exhausts its determinate content; its dark pregnant
inwardness becomes clear, everything comes to appearance in
superficial immediacy. Since all the conditions of the phenomenon are sublated in its unconditioned independence,
the existing phenomenon refers itself to its own
relative center, as to a midpoint absolute in itself,
and rests on itself alone. With the same unconditioned
independence one phenomenon steps forth after the
other, and since each of these has its sufficient
ground in itself, the absolute centrality of the primal ground vanishes
in the relative centers of the innumerable phenomena.
The being-for-itself of the phenomena makes up their immediate
unconditioned existence; every phenomenon comes to appearance
in its own light; its relation to other phenomena,
which was to show its infinite conditionedness, is an unessential
semblance, since the conditions are everywhere sublated
in unconditioned independence. The mutual
relationship of the phenomena, their being for one another, is then, considered from
this side, only an external reflection.
But phenomenal being-for-itself is just as immediately
a being for other; reflection in itself is itself
superficial; it is essentially a reflection in other,
since it is phenomenal. According to the immediate
reflection-in-itself of the phenomenon, the conditions are sublated in
its unconditioned being-for-itself; according to the immediate
reflection-in-other, the unconditioned being-for-itself dissolves
into sheer conditions. Only in this abrupt, unmediated
interchange does superficial existence come to appearance.
The form of the absolute's existence in the series
of phenomena is therefore the infinite progress. The absolute
as such cannot exist in any single phenomenon,
for the unconditioned independence of the phenomenon
% --- p. 161 ---
\setcounter{footnote}{0}%
\opage{161}immediately turns over into non-independence; just as little can an absolute
existence be expressed in the sum of the phenomena, for
every existing phenomenon is burdened with the same
essential contradiction. Phenomenal existence
is thus infinitely fragmentary; it must incessantly
begin over again, and the progress that was to contain
the new beginning is only a semblance-progress, in which
phenomenal non-independence everywhere comes to appearance.
As the determinately existent had to perish
because in its immediacy of being it lacked all
ground, so the existing phenomenon must now also
perish, because it would sublate the absolute ground and exhaust it in superficial
determinateness. The more clearly
the particular independence emerges, the more sharply
the transition shows itself in which the independent
finds its dissolution. Now as the absolute thus
seems to run off into sheer relativity and to scatter its
content into peripheral centers, phenomenal
existence at the same time betrays its infinite negativity;
the primal ground steps forth anew as the leveling
principle, and takes back its absolute centrality.

2) The absolute exists as ground.

Only in relation to the immediate independence of the phenomena
does the primal ground come to appearance as the all-dissolving
ground; but everything that it thereby dissolves it has
itself posited; it therefore only takes back its own,
in order to affirm itself. It has already been
shown that the esoteric and exoteric content of the ground reflect
each other mutually; the esoteric content is in itself
ground-identity, and thus has at one stroke all distinction-moments
in its inner infinity, all principles of determinate existence
in its one ground-principle; but in this its
immediate primitiveness the primal ground has as yet no
% --- p. 162 ---
\setcounter{footnote}{0}%
\opage{162}existence; it is as yet only presupposed as the real beginning
of all-being, as essential becoming, as originality
itself. In contrast to this, the exoteric content presents
an evolution of the thoroughgoing
distinction. Here unity is posited in the form of manifoldness,
here is a sporadic plurality of principles and relative
beginnings, an innumerability of determinately existent grounds
and finite conditions. It has further been shown
that each of the terms in the central opposition must itself
be presupposed as absolute, in order that the ground-content, doubled
in itself, can be developed on all sides. The absolute
primal ground is, in its original essentiality, infinite mediation
of abstract identity, negative, transcendent unity,
and abstract diversity, negative, transcendental
manifoldness. This mediation grounds itself in its
own primitive movement. --- Only as the primal ground takes back
its concrete unity through the exposed
manifoldness is its esoteric content developed;
but in order that the many-sided sundering can be realized,
the exoteric process of finite grounds and conditions must be
presupposed as the medium through which the return
incessantly takes place. Therefore the consecutive
ground-movement must also complete its circuit,
and its exoteric content be uninterruptedly restored through
the medium of the primal ground. Were the exoteric
progress of relative grounds and separate conditions only
a fleeting semblance, did it not contain an absolutum relative
in itself, then the primal ground would lack the real medium
for its immanent explication, then it would be without
the otherness of discrete being in which the pure light of essentiality and
originality could refract its rays;
but conversely, if the primitive ground-movement perished
in the exoteric and consecutive one, if the refracted
rays were not cast back again into the depth of the primal ground,
% --- p. 163 ---
\setcounter{footnote}{0}%
\opage{163}then the ground-mystery would remain in eternal darkness, then
the original unity would have to be broken off in the alteration
of the ground, and the absolute deny itself.

In phenomenal existence the exoteric
mediation of the ground has now culminated. The rich
variations of manifoldness have come to appearance; the positive immediacy
has fixed itself in the relative
centers of the existing phenomena; the exoteric ground-content is posited with
absolute superficiality, and thereby presupposed as what is
absolute in itself. But this immediate positivity is at the same time
infinite negativity. α) In order to exist, the one
independent phenomenon must presuppose the existence of the other;
if there is one single existing phenomenon, then there is
at the same time an infinite manifoldness. The
independent phenomenon thus reflects in its immediate
existence an innumerability of independent phenomena.
That every phenomenon subsists in itself will then mean
that it subsists in the subsistence of all the others. β) But
the reflection of the phenomena in one another is superficial; the
relative centers repel one another mutually; the one
phenomenon is immediately the negation of the other; insofar as
the one exists, the other cannot exist; if
the one phenomenon therefore subsists in the subsistence of the others, then it subsists
eo ipso in its own non-subsistence. The negation
from which existential immediacy proceeds
is itself an existing immediacy from which negation
results. γ) This contradiction does not find its resolution
in the consecutive process, which on the contrary must be regarded
as its infinite continuation. The essential thing
is that existential diversity comes to appearance;
but as it becomes phenomenal, it becomes at the same time
immediate and thereby unessential. It is thus only
its unessential side that can vanish in the dissolving
ground, and it is this dissolution by which the
% --- p. 164 ---
\setcounter{footnote}{0}%
\opage{164}primal ground sublates its alteration and returns into
itself. Only this return is the true,
essential advance. The pleromatic center has,
during its inner self-separation, struck into the relative
centers, because the ground first comes to true subsistence
in existence. But in phenomenal existence
the primal ground is lost and its centrality sublated. The
pleromatic centrality therefore subsists as the phenomenal
center subsists, i.e. it subsists in its non-subsistence,
it exists in the negation of its existence.
Only as it consummates this negation does it affirm
all existential moments, and thereby posit itself
as what exists in itself, as the overreaching
centrality, as the primitive presupposition of the subsisting phenomena,
as the existential \textbf{depth}.
3) The absolute exists in the mediation of the central-existence and the phenomenal progress.

The manifoldness veiled in the unity of the primal ground
is a negative identity which sublates itself in order to
posit its distinction. That all-identity has in itself infinite
distinction, and that conversely manifoldness and infinite
distinction are reflected in themselves by reflecting identity,
is shown by pure mediation. Nevertheless this
mediation is only the form of the absolute's infinite
self-determining, and as such without any determinateness;
it is only formal. But formal determining posits
and presupposes determinateness; otherness emerges out of
the unity of reflection, and form deposits its content;
manifoldness is therefore to be determined and posited as
manifoldness, and the distinctions are to be stamped out in their diversity;
conversely the ground-unity is also to be posited
and articulated in its own concrete self-reflection, so that
all distinction-moments return and vanish in
% --- p. 165 ---
\setcounter{footnote}{0}%
\opage{165}it as real semblance, as immanent reflexes: the surface
is to be reflected in the depth, and the depth in the surface.
As now the outer manifoldness, reflected in itself,
becomes \textbf{posited}, the immediate
all-unity is at the same time \textbf{sublated}; the ground falls to the ground in the existing phenomena; the reflection in itself of ground-identity
vanishes as reflection in manifoldness. Here
is therefore an essential transition, a universal beginning,
by which phenomenal manifoldness sublates
the indeterminate ground-unity and begins the exoteric
progress. But on the other hand the reflection
in itself of manifoldness through the sublated
ground-identity is only the presupposition for the self-reflection of the all-unity
in the sublated manifoldness. When the primal unity
is \textbf{posited} according to its concrete identity with itself,
then the outer infinity of the existing phenomena
has at one stroke \textbf{vanished} from the series, and the determinate
distinction-reflexes must be sought in the clear
depth of the ground-existence. The immanent central-existence therefore results
from the return of the phenomena into the ground; it proceeds
from phenomenal conditions, and thus itself has
a \textbf{relative beginning}. In \textbf{the mediation of the
universal} \textbf{and relative beginning is} \textbf{the
thoroughgoing and} \textbf{true real beginning first to be posited}.
a) Every existing phenomenon is burdened with the
contradiction that its existence is at once infinitely conditioned by the universal mediation of the absolute,
and yet in its immediate determinateness restricted to
a peculiar circle of finite conditions. That every
phenomenon exists in unconditioned being-for-itself thus
follows from the fact that the one phenomenal circle of conditions
is limited by the other; the specific
ground-relations break off the mediation of the universal ground
to infinity, and \textbf{the} relative, partial beginnings
% --- p. 166 ---
\setcounter{footnote}{0}%
\opage{166}reduce the universal beginning to an external reflection.
But precisely thereby the partial beginning becomes
a self-contradiction; in the phenomenal series:
\dots{} B, C\dots{}. C cannot begin except on the presupposition
of B; the circle of conditions in which C
exists is itself conditioned by the circle of conditions
in which B has its existence; B again presupposes
A, and so on to infinity. Could but one
single phenomenon begin its existence, then
all the rest would begin at the same time: every immediately
determinate beginning is therefore immediately conditioned
by the universal beginning. Now as the phenomenal
series is split apart into outer manifoldness,
into partial beginnings in which there is no universal beginning,
so the universal beginning itself is only
presupposed as the veiled manifoldness of the ground-unity,
in which there is no partial beginning, because
no existential distinction has yet been posited.
The transition from the universal to the partial beginning
can nevertheless not take place by a leap; as
the phenomenal series sinks as semblance into the essence of the primal ground
without being able to begin, it presupposes
its presupposing principle, i.e. it posits all its
partial beginnings back into the concrete inwardness of the universal beginning,
and presupposes the latter as predisposed
to posit the specific distinction. Insofar as
the determinate phenomenon begins its existence on
the presupposition of separate conditions, externality
and otherness go back into an infinite regress,
and the conditions of the conditions flee away into the indeterminable;
but insofar as this external reflection
in the series of partial beginnings also has its
presupposition in the predisposed universal beginning,
all outer phenomenal conditions are reflected
% --- p. 167 ---
\setcounter{footnote}{0}%
\opage{167}in the central-existence's own depth. The immanent
reflection of the conditions in the universal beginning and their
external reflection in the partial beginnings of the series
are therefore one and the same reflection. For an existence
to begin, it is then required that its peculiar conditions
be found present in the series of phenomena, and next
that the interruption and absence lying in
phenomenal otherness be integrated by the essential
presence of the immanent conditions of the universal beginning,
which are presupposed in existential centrality.
As the conditions and the relative grounds are,
so too is the central-existence in its primal ground
disposed to posit the new beginning; but
conversely the outer specific conditions and relative
grounds are the central-existence's own presuppositions,
by whose sublation it posits the unconditioned. The
finite grounds and phenomenal conditions spin themselves
out as a tautological continuation of the one relative
existence in the other; but only when the lightning of centrality
strikes down into the circle of conditions is
the real beginning posited, and a new existence emerges
out of the stream of mediation.
The central-existence now has its contradiction in this, that it
is at once to contain the universal beginning specified
in itself, the mediation of all conditions
in one absolute unconditionedness, and thus to subordinate
the phenomenal progress to itself, while it must nevertheless also
presuppose the immediate
separation of the phenomena, in order by their sublation to gain in existential
substance, and be disposed to continue the series of phenomena. This contradiction finds its
resolution in the immanent self-movement by which
the ground-existence uninterruptedly releases its otherness.
The self-development of the absolute into existence is a self-continuing
explication of the double
% --- p. 168 ---
\setcounter{footnote}{0}%
\opage{168}content of the primal ground. As the ground-unity discloses its mystery
and expends its content in the cycle of the phenomena,
then this content shows itself in the first place as its
\textbf{own} content; for it has now become clear how
it dissolves the existence of the phenomena in order to take
the specified content-moments back into its own
existential self-reflection; next, it is an \textbf{other}
content, for it has likewise been shown that the dissolution of the determinate
phenomena yields stuff and material
for other phenomena in the exoteric series.
The ``\textbf{becoming}'' of immediacy is posited in the central-existence
as a moment that is sublated in the ``\textbf{has been}''
of reflection, while on the other hand, as regards the phenomenal
progress, the essential ``has been''
shows itself as a vanishedness that must constantly
be replaced anew in the immediate becoming of the partial beginning;
when the existing phenomenon, in accordance with its
pure superficiality, uninterruptedly takes its beginning
from other, partly from other phenomena, and partly
from the depth of the primal existence, the central-existence takes,
through the presupposed sublation of otherness, i.e. through that of the partial beginnings,
its immanent real beginning
from itself; what is therefore, with respect
to the series, infinite regress and progress, that
is, with respect to the central-existence, infinite self-deepening
and self-development. By this immanent
self-development the central-existence vindicates its absolute
dignity, and subordinates the phenomenal series to itself, so that the dialectical pulse of centrality beats
in every existing phenomenon. The double foundation of the primal ground required for its development
an esoteric and exoteric mediation, a primitive
and consecutive ground-movement. By virtue of the
doubleness of the content, mediation and
% --- p. 169 ---
\setcounter{footnote}{0}%
\opage{169}formation have therefore become doubled, and this developed
form-distinction is posited as a moment in the one absolute
formation. The content shapes its own opposition
out of the unity of the ground; the opposition is therefore
not immediately dualistic. But the same ground-content
that in formation posits its many-sidedly developed
identity at the same time develops this content-identity
in a double formation. The primitive and
consecutive ground-movement are both absolute in their mutual
reflection. By this immanent
form-distinction the content becomes distinct from itself.
The external form-distinction leaves the content
unchanged in what is essential; the finite form presupposes
a given content as stuff and material; but
the immanent form, which penetrates, forms
and determines all content-moments, thereby gives
the content its absolute determinateness, i.e. it determines
it as its determinate content, so that the one
content can be separated from the other only by virtue of form.
If absolute existence therefore has a formation
reduplicated in itself, if its development has two
immanent primal forms, then it also has, in its concrete content-identity, two opposed contents.

Only in the movement of mediation can the central-existence
separate itself from the phenomenal existences; but
as the many-sided existential mediation has now been carried through,
immediacy returns. The central-existence
exists in the mediation of all the presuppositions of the absolute,
in the identity of all phenomenal and immanent
conditions; it exists immediately as the
absolutely unconditioned. But the same holds of
phenomenal existence; it too takes its real beginning
in all presuppositions of the absolute; its
separate circle of conditions is likewise mediated
% --- p. 170 ---
\setcounter{footnote}{0}%
\opage{170}through all the immanent conditions of the central-existence;
it too exists in its immediacy as the absolutely
unconditioned. Absolute form is then, in consequence
of this, posited in every content, and conversely every
content is peculiarly determined by the many-sidedly determining
form. Thus the double ground-content,
by virtue of the one determining form, taken immediately,
is one and the same content, and yet it is just as
immediately an altogether distinct content because
of the innumerable form-determinations posited in absolute
formation. The absolute is in all formations
always identical with itself, and therefore one single
unconditioned in itself; but if the exoteric content is unconditionally
sublated in the esoteric by virtue of the immediate,
existential unity, then the esoteric has just as unconditionally
vanished in the exoteric because of the immediate,
existential manifoldness and distinction. As
absolute unconditionedness makes its appearance in the central-existence,
it breaks through at the same stroke
in all relative existences, and sparkles in innumerable
centers. In this immediate identity of central
and phenomenal existence \textbf{that which} exists makes
its appearance, partly as the \textbf{thing} in \textbf{itself}, partly as the \textbf{things}
in \textbf{their} properties.

According to its immediate being-in-itself, the central-existence
is to be held fast as the abstract essence-existence,
as blind transcendence, the pure negativum of
all phenomena; in consequence of this it cannot then
be determined more precisely either; for all determinations vanish
as reflexes of an outer manifoldness that does not
concern its own essential reflection in itself at all;
according to this abstract being-for-itself, that which exists
is therefore posited as the transcendent, as the thing in itself
(``Das Ding an sich''). When it is said that the thing
% --- p. 171 ---
\setcounter{footnote}{0}%
\opage{171}in itself cannot be cognized, this means in other words
that there is nothing at all in it to cognize, since
it excludes all distinction and relation; the thing in itself
is thus to be regarded as the existing tautology,
by which the proposition: A \textbf{is} = A is transformed into:
A \textbf{exists} = A. That the thing in itself, despite this
tautology, nevertheless does not altogether vanish and become
no thing at all, it owes to its relation to the denied
phenomena. As unity and identity belong
to the thing in itself, so manifoldness
and distinction, on the contrary, fall over onto the phenomena: these are
then not merely distinct from the thing, but also from one another.
But on the other hand the
existing identity is always conditioned by the existing
distinction; the thing in itself is itself determined by the fact that it
excludes all phenomenal determinations. The distinguishing
reflection by which the phenomena are separated from
the thing in itself is therefore the very same by which they separate
themselves from one another. That which gives every phenomenon
existing identity, the unconditioned being-for-itself
that keeps itself indifferent toward every relation
to the other phenomena, that which exists in the phenomenon,
which remains when we abstract from all
relative determinations, is consequently also a thing in itself.
If identity is to exist unconditionally in and for itself, then
distinction must be excluded; but by this exclusion
distinction itself is fixed in immediate being-for-itself,
and thereby comes to exist as immediate identity.
The one thing in itself thus becomes at one
stroke a manifoldness of things in themselves.

The reflection of the thing in its being-for-itself reflects
at the same time its being in other. This other is now itself
a thing; that the thing is reflected in itself as in its other
will then mean that it is reflected in another thing.
% --- p. 172 ---
\setcounter{footnote}{0}%
\opage{172}Insofar as those that exist are, in their manifoldness, things,
they keep themselves, each one, in the extreme of distinction and opposition;
but as they are determined by determining
one another mutually, they at the same time identify their existential distinction: in this mutual reflection the things \textbf{have}
\textbf{properties}.

The property is in the first place the thing's \textbf{own} determinateness:
the thing's determination in relation to other
comes to appearance in the property as its own immediate
self-determination: as the property is, so too
is the thing constituted; next, the property is an expression
of the thing's determinateness in relation to the property of another
thing. What holds of the things themselves
then holds also of the properties: their reflection-in-other
is their own reflection-in-itself. According to this
double reflection, the things and the properties each
by themselves have a certain independence, a relative subsistence.

Everything that exists is always a something; but not every
something can for that reason be said to exist. What ``something''
is in immediate being, that the \textbf{thing} is as
the phenomenal essence; \textbf{something} is, as such, an
immediacy that has not yet run through reflection;
the \textbf{thing}, on the other hand, is an essence-existence whose
immediacy reflects the sublated reflection; something is, as the determinately existent, the pure negation of
\textbf{other}; the \textbf{thing}, on the other hand, in identity with its
constitution, is nevertheless distinct from the \textbf{property};
something \textbf{is} quality, but the thing has property; therein
lies the distinction.

A more essential likeness obtains between the thing
and the \textbf{phenomenon}. The existing phenomenon
is the immediate unity of thing and properties;
insofar as the thing is phenomenon, it can be dissolved into
matters, and run through a series of finite
% --- p. 173 ---
\setcounter{footnote}{0}%
\opage{173}changes: \textbf{the thing consists of matters}\footnote{Hegel's Logik, 2tes Buch. p. 155 (Werke 4de Band) and J. L. Heiberg's speculative Logik p. 51.}; but the phenomenon
nevertheless stands below the thing in existential development,
for phenomenal unconditionedness has not yet
reflected the essential distinction \textbf{of the thing and the properties}.
Because of its pure superficiality,
phenomenal existence cannot lose a
single one of its determinations without being at the same time
altered thereby; the phenomenon is in this respect a Noli
me tangere; between thing and properties, on the other hand,
there obtains the separation that the thing itself is posited as
the essential, while the properties in their mutual
diversity show themselves as the unessential; therefore
the thing can very well lose one of its properties without
therefore being dissolved at once\footnote{When a square table is changed into a round one, it becomes another phenomenon, but is nevertheless the same thing. A flower without scent is also a flower.}. That the thing has properties
is also expressed by saying that it \textbf{subsists} in its
properties, or, what is the same, that it subsists
in the subsistence of the properties. Insofar as there is now
a double subsistence here, the thing's existence is not
exhausted in any single property whatever; but since the double
subsistence is nevertheless grounded in one identity,
the thing as such cannot bear a change
of all its properties either, without itself having to be
altered thereby at the same time.

In the sharp sundering of thing and properties
the distinguishing and fixing reflection plays its very freest
game. In the distinction between the things and the properties
the things' own mutual distinction is held fast.
In the change of the properties there does indeed take place a
% --- p. 174 ---
\setcounter{footnote}{0}%
\opage{174}change in the things; but since these, by subsisting in
the properties, also have a separate subsistence in themselves,
they are not \textbf{for one another} mutually, and consequently
independent of the changes occurring in the properties.
In such an extreme of independence the things
can be fixed, just because they have the properties, in whose
determination it lies precisely to let the things be
for one another, and thereby reflect the stream of changes
in which the things themselves keep unchanged.
As the whole exchange of relative determinations is, for the things
themselves, an indifferent and external reflection, so
too are the things indifferent toward one another;
they can meet in the most diverse relations without
their firm subsistence suffering thereby. The changes reflected
in the properties therefore do not touch the things'
essence, but only their phenomenal side.

If, on the other hand, the identity of thing and properties is held fast,
then the essence of the things becomes reflected in its phenomenon,
and the unity of both posited as the true existence of the things.
The distinction that comes to appearance between essence
and phenomenon then coincides with the things'
own subsistence, and as the thing itself exists in immediate
unity of its essential and unessential side,
so too the properties have in themselves an essential and
unessential existence. By losing an unessential property
the thing itself suffers an unessential change, whereas
it is essentially altered when an essential property
vanishes. The reflection of the things and of the properties
is therefore fundamentally one and the very same, and what is
peculiar in reflection lies in this, that it both
\textbf{is} sublated and yet not sublated.

If the thing's reflection-in-itself is \textbf{one} with its reflection
in the properties, then the thing's concrete determinateness
immediately coincides with its properties, and these
% --- p. 175 ---
\setcounter{footnote}{0}%
\opage{175}vanish as reflexes in the thing's ``haecceity''; the phenomenon is then negated as a mere phenomenon, and thereby
sublated in the determination of its ``noumenon'': the thing
\textbf{makes} its appearance in its properties. But since the reflection of the properties
in the thing is now also \textbf{one} with
their own reflection in themselves, independence falls over
onto the properties; the thing's determinate phenomenon is
thus posited in its distinction from the thing itself, and the latter,
on the contrary, presupposed as a ``noumenon'' indeterminable in itself:
the thing's existence in the properties is \textbf{only} apparent.
The existential antinomy, that the thing is at once
both transcendent and immanent, lies in
appearance \textbf{itself}.

The things in their phenomenal determinateness are immediately
opposed to one another, and thus form a series
of existing phenomena. Could essential existence
be exhausted in phenomenal immediacy, then
the thing itself would be identical with its appearance, and
every illusion would vanish. But now the thing in its
phenomenon is at the same time distinct from the phenomenon; it is
then, in consequence of this, opposed to its own appearance, and
its immediate subsistence in the phenomenon illusory.
What is essential in appearance, namely the thing in itself,
is therefore just \textbf{that} which does not make its appearance; thereby
appearance itself sublates its own essentiality.
On the other hand this essentiality is at the same time
posited in its sublation; for however distinct appearance
may be from the thing itself, it nevertheless lets the thing's
abstract being-for-itself come to appearance; what, on the other hand,
it seems to conceal is the thing's own concrete determinateness.
Now it has, however, been shown that haecceity,
by which the thing is determined in itself, and distinct from every
other thing, is necessarily reflected in a cycle
of essential and unessential properties. Either
% --- p. 176 ---
\setcounter{footnote}{0}%
\opage{176}the thing in itself must then remain an empty abstraction that lacks
all existence, since it lacks appearance,
or else its otherworldly existence must itself become an
appearance in otherworldly phenomena. The illusion that
lies in the abstract opposition between thing and
phenomenon vanishes as it makes its appearance, or
in other words: the antinomy of appearance is only apparent.

In appearance existential reflection culminates.
That which exists is only something apparent; for it always subsists
in another subsistence; this other subsistence
is itself negative, and subsists at the same time in its sublation
as something apparent. But since every existence, just
because it is apparent, subsists in the negation of
its subsistence, and the one existent is, in relation to
the other, this negation, the one existence subsists
in and by the other, and thus obtains,
because of appearance, a subsistence reflected in itself.
α) The thing in itself suffers from the contradiction that there is to
be an abstract independence without semblance of reflection,
and precisely thereby becomes an abstract semblance without any independence.
The thing's unity in itself is only a semblance-unity,
for it exists in a manifoldness of things in themselves;
but the thing's essential manifoldness is also a
semblance-plurality, for it exists only in the one unreflected
thing in itself. The thing's absolute essentiality is then an essence
apparent in itself, and must therefore make its appearance
in \textbf{its other}, i.e. in properties and phenomena.
β) But by coming to appearance in the properties, the thing
itself falls into a new contradiction. Its existence oscillates
now between the two extremes, either to expend its
independence in the subsistence of the properties, or to
absorb the properties in its subsistence. If the first happens,
then the ground-existence is exhausted in the existing
% --- p. 177 ---
\setcounter{footnote}{0}%
\opage{177}phenomena, and since the phenomena negate every opposition between an essential and unessential existence, the things themselves vanish in the changing forms of matter; if the latter happens, then the indeterminable \textit{Noumena} of the things fall back into a merely apparent essentiality, and flow away as non-independent moments in the absolute central-existence. \textit{γ}) This contradiction dissolves in the phenomena's own apparent being; for this too subsists in the reflection of abstract independence and non-independence. As the things first come to appearance in themselves by making their appearance in the phenomena, so conversely the phenomena, by making their appearance in one another, become the things' own appearances. Appearance is the \textbf{posited} unity of the identity and distinction of existence. Distinction exists, for every existence has its particular subsistence; its determinate form and content separate it from other contents in other forms; its independence excludes every other independence; it is, as such, immediate and positive. In absolute negativity there therefore exists an innumerability of positive contents; in absolute independence innumerable independents subsist: existential plurality \textbf{is} the form in which \textbf{unity \textbf{makes its appearance}.} But in the positivity of plurality lies its negativity; the existing distinction is in itself apparent; therefore identity comes to existence. Since the opposition between the distinct contents is in itself apparent, it must \textbf{make its \textbf{appearance}, \textbf{as a} manifoldness} of distinction-moments, in the absolute content-identity; there exists therefore only one absolute content in all relative contents, only one absolute form in all relative forms, only one absolute positivity in all negativity, only one absolute independence in all independents. Now as both forms of all-existence are in appearance
% --- p. 178 ---
\setcounter{footnote}{0}%
\opage{178}posited with the same unconditioned validity, the existential opposition between the depth of centrality and the superficial phenomena has ceased to be a contradiction. As the immanent and essential appearance subsists in mediation, whose all-unity is the depth, so the unessential appearance exists in immediacy, whose phenomenal distinction is superficial. Appearance is itself identity of mediation and immediacy. The things in their phenomena are immediate existences; they come to appearance as determinately existent. The thing's determinate being also lies in its appearance; for what is to make its appearance must already have gone through its becoming. But this determinate being is at the same time essential, and as becoming is reflected in determinate being, origin comes to appearance. The thing's opposition to its phenomenon and immediate appearance is hereby reflected as an essential distinction between its \textbf{determinate existence} and its \textbf{origin.} According to the immediacy of appearance, the one thing seems to originate from the other, and yet origin itself does not come to appearance; for mediation, which here shows itself as the restless, superficial turnover of distinct contents and forms, uninterruptedly wrestles with immediacy, and is as such merely apparent. As the phenomenal existence of the things is to go together with itself into absolute subsistence, it repels itself from itself. But by making its appearance in its own otherness, i.e. in its merely apparent originality, it at the same time lets its \textbf{essential origin,} the absolute other by which it is itself posited, come to appearance, and subsists in its existential manifoldness as a result that presupposes its presupposing principle, i.e. \textbf{the \textbf{depth} of centrality.}
% --- p. 179 ---
\setcounter{footnote}{0}%
\opage{179}In the central-existence immediacy is conversely a
reflected semblance of mediation. The existential depth cannot
immediately come to appearance, for thereby it would
itself be transformed into a superficial phenomenon; just as little
can it be a thing in itself, since the thing's distinction from
the phenomenon is itself a reflex of the essentiality reflected
in itself: the central-existence therefore exists in the form of
essential mediation. According to this mediation its
immediate independence is its own infinite negativity;
hereby its content is presupposed as the absolute originality
of the ground; but since this originality
is apparent in itself, it gives the things their
origin, and posits the outer existential manifoldness
as the medium in which it makes its appearance; as
negative existence it therefore subsists in its own negation,
i.e. in the positive immediacy of distinction and manifoldness;
but since this outer immediate existence now
also subsists in its negativity, since it is itself
only an appearance of the apparent unity,
its independence is posited in its sublation. By
subsisting in outer existence, the central-existence subsists
in the sublation of existential manifoldness;
as it hereby negates its own negation, it therefore exists
in the appearance of all immediate existences
as the identity affirmative in itself and totalized. When
the central-existence then subsists in the positive
unity of mediation which in its content sublates all distinction-moments
as negative and apparent, then
the content of phenomenal-existence subsists in an immediate distinction
whose unity in itself is apparent and negative.
By positing this double subsistence in one many-sided
appearance, absolute existence passes over into
totality.
% --- p. 180 ---
\setcounter{footnote}{0}%
\opage{180}\underafsnit{c) \textbf{Totality}.}
\paragraf{26}{Inner and Outer, Force and Product, Preformation.}

In appearance all existential moments come
to appearance, both in their immediacy and in
their mutual transition and immanent reflection,
whereby all-existence is posited as \emph{totality}. According to
its essential side, totality is determined as
a realm of laws, according to its immediacy, on the other hand,
as a world of phenomena: an abstract opposition that is sublated in the essential relationship, i.e. in
the reflection of the reflected and the immediate totality. According to the totalized form-distinction, which
is at the same time a distinction in content, the essential
relationship presupposes the totality reflected in the inwardness of mystery and
the depth of selfhood on the one side,
and the immediacy-totality unfolded in an external series
(whole and parts) on the other,
and thereby determines itself into a relationship between the inner and the \emph{outer}. Inwardness, incommensurable with its outer forms,
expresses itself in inexhaustible productivity; the identity of inner and outer is reflected in force, and totality develops
into a relationship between force and \emph{product}. As the formation of a new shape presupposes the
producing force, so the specific productivity of force
presupposes produced shapes in the
% --- p. 181 ---
\setcounter{footnote}{0}%
\opage{181}infinite, so that every real formation reflects its
\textbf{preformation}. According to the reflection of appearance, every moment in the existence of the absolute is thus thoroughly

clarified that the ground-existence can be posited as the medium for
the infinite progress of phenomenal-existence, just as % print: Phænominalexistentsens (misprint)
phenomenal existence is conversely posited as
the medium through which the ground-existence gains
its self-deepening: the esoteric and exoteric content
shine through each other. As, therefore, nothing
is lost when only the ground-existence is held fast, since
it itself contains all the moments of phenomenal-existence,
and posits them out of itself in order to affirm its own independence
in their negativity, so
no essential term of absolute existence is lost either
if phenomenal-existence is posited in its
self-reflection. For as it is posited by other, it reflects
its absolute otherness from all sides, and is thus
so penetrated by the content of the ground-existence that the unity of the depth
makes its appearance in its superficial
manifoldness. By being posited in the form of appearance,
all-existence has developed into \textbf{totality}.

Totality is then now the most adequate form for
the identity of the absolute in all independent distinctions.
It lies in totality that every existence has its
subsistence, that the sundering and separation are carried
through; but it also follows, on the other hand, from
the essence of totality that no existence can be isolated
into abstract being-for-itself and fall outside the wholeness.
Both demands of totality are fulfilled in appearance,
which at once posits every independence
as infinitely \textbf{negative} and infinitely positive.
% --- p. 182 ---
\setcounter{footnote}{0}%
\opage{182}But insofar as totality can be developed from
any existence whatever, it is as yet, according to appearance,
entirely formal, and lacks the peculiar
determinateness. In pure appearance all existences are
sublated: they are thereby all equally finite and
infinite, equally mediate and immediate, equally
central and peripheral. The point of departure for the totality-development
is, in consequence of this, entirely indifferent.
The category of appearance presents itself as a
transparent totality-form; but totality itself
requires a more precise determination.

The pure extremes of totality are: on the one side
the immanent essence of appearance, an identity that
sublates every phenomenal distinction and affirms
totality as the all-embracing and all-determining unity of the things; on the other side the self-separation of appearance
into the concrete immediacy of the phenomena,
whereby the sundering is consummated, and totality
shows itself in the whole manifoldness of particular existences.
Since, further, each of these sides with
all its independence is a moment of appearance,
and consequently itself has only an apparent independence,
they are not externally compounded, nor is a third needed to form their union:
the own reflection of appearance makes up their
true subsistence. The reflected subsistence of both sides
is now to be determined more precisely as the \textbf{law} and haecceity.

a) In the totality of its essential and unessential
properties, i.e. in its immediately determinate appearance, the thing is posited with its haecceity; but
as \textbf{this} thing (hæc res) it is thereby at the same time
limited and distinct from all other determinate
things, each of which expresses its
% --- p. 183 ---
\setcounter{footnote}{0}%
\opage{183}haecceity. The whole totality stamps itself out in the outer, infinite manifoldness of haecceities. In
the immediate determinateness and essential
delimitation of the things lies their finitude; one
thing arises and passes away after another: haecceity
is perishable and mutable. But as
the immediate independence of the things is apparent,
so their changeability and mutability is
also apparent; however many things
arise and pass away, there will nevertheless always remain an innumerability of determinate things.
The amount of subsisting things is infinite; the amount of them that arise and
pass away is consequently also infinite: it is the
whole totality that incessantly arises, and the whole
totality that incessantly passes away; for in essence
quantity is sublated; absolute reflection
reckons with infinite magnitudes. Since the outer manifoldness and distinction of haecceity is thus transformed to
infinity into an indifferent and unessential
semblance, the changeability of the things is negated
and at the same time taken up into their positive unchangeability,
and the law posited as their \textbf{essential} foundation,
as their \textbf{immanent subsistence}.

By being posited as the foundation of the things
the law acquires a certain likeness to the essential thing,
to the ``thing in itself'', since both present themselves as
the true, universal essence of the determinate things. Nevertheless
the law cannot be hypostatized into a thing, since
it negates every independent immediacy.
When the ``\textbf{thing} in \textbf{itself}'' is determined as the
foundation that is, which relates itself indifferently to
the phenomena of existence, the ``\textbf{law}'', on the other hand,
has its subsistence in \textbf{the} subsisting things; the ``thing
in \textbf{itself}'' rejects appearance, and shows itself
% --- p. 184 ---
\setcounter{footnote}{0}%
\opage{184}thereby as the \textbf{transcendent} essence; the ``\textbf{law}'',
on the other hand, makes its appearance in the appearance of the things,
and holds only as the \textbf{immanent} essence.
What the ``\textbf{thing} in \textbf{itself}'' is in the form of \textbf{sublated}
reflection, the \textbf{law} is in the form of \textbf{posited} reflection;
when the thing in itself, by an otherworldly abstract immediacy, pushes the phenomenal determinations
away from itself, the law itself comes openly
forward as the essentiality in the exchange of the phenomena.\footnote{Although Kant demonstrated that we cannot cognize ``das Ding an sich'' (the transcendent essence), he was nevertheless convinced that we can investigate the laws of things; for the law (the transcendental, i.e. apparent and immanent essence) holds not of \textit{Noumena} but of \textit{Phænomena}.} When the changing manifoldness of the things is posited
as reflex in their immanent unity, it is hereby
at the same time presupposed that this essence-identity is absorbed
as reflex in the determinate immediacy.
Since now the law itself is the essence of the things, and since the
independent immediacy of the things proceeds precisely
from their own essentiality, haecceity springs
from the law's own positing reflection. This
contradiction, that haecceity at once proceeds
from the law and yet falls outside the law, is at
once determined by the law and yet indifferent to the law,
at one stroke is posited and sublated by % print: Lvven (for Loven) (misprint)
the law, is expressed in the proposition that \textbf{the thing is
subject to the law}. The subsistence of the thing and of the law
is hereby doubled: appearance is the thing's
truth; it therefore lets the essence of the law
come to appearance as the otherness in which it
itself subsists; nevertheless it does not subsist as law,
% --- p. 185 ---
\setcounter{footnote}{0}%
\opage{185}but as immediate independence. As regards
the law, it too has haecceity
as the other in which it comes to existence,
but nevertheless does not subsist as \textbf{this} thing,\footnote{The stone is heavy; it therefore has an essential subsistence in its gravity, for however it may otherwise be constituted, it is nevertheless subject to the law of gravity. Yet it does not exist as gravity, but as stone, i.e. as the other in which the law of gravity can come to appearance. If someone sees a human being die merely as this human being, then he senses only a fact; if, on the other hand, his thought sees humanity die in this human being's death, he cognizes the law of mortality.} but
as its negation. Both principal sides are thus,
according to the reflection of appearance, equally positive and
equally negative, but according to the immediate distinction
between law and haecceity abstractly opposed to each other.
If now the thing's subsistence in the law is posited as
the positive, then haecceity is a vanishing
semblance, a negativum through which the \textbf{unity} of totality
is affirmed; if, on the other hand, the thing's subsistence in
immediate independence is posited as the
positive, then the law is the negation by which haecceity
is determined, and the \textbf{allness} of totality is affirmed.
While the law, therefore, in the first respect is
presupposed as the foundation of the things, in the
latter respect it is posited as the things' \textbf{immanent
and essential limit}. The determinateness of the things
is their finitude reflected in the law, which is
thus posited in a double form. In the form
of essential identity the finitude of the things comes
to appearance as a common vanishing in the essence of the law,
as in their pure infinity and
% --- p. 186 ---
\setcounter{footnote}{0}%
\opage{186}universality: \textbf{they are all equal} \textbf{before the law}; in the form
of essential distinction finitude comes to appearance
as the limit that the thing cannot overstep without
its subsistence being altered; the law sunders and separates
the things: \textbf{thus far and no further}! is the
commandment that holds for each single existence. The separation
of the things and their transition into one another make them
sides in the immanent relation that is determined
as their law. But the things are at the same time more than
mere sides in the law of the relationship; they are haecceities whose immediate independence shows itself indifferent
toward the mutual relation. --- The medium
between law and haecceity is the essential property.
As \textbf{property}, namely, it is a moment
in haecceity, a determination in the for-itself
immediacy by which the thing itself falls
outside the law and the relationship; as \textbf{essential} property
it is itself indifferent toward haecceity, reflects
its unchangeability in the latter's changeability,
and shows the thing's determinateness by, and dependence
on, the law. Since now the same thing has a plurality
of properties, since every property shows that
the thing's reflection in itself is posited as its reflection
in other things, and the thing's independence as a
side in the immanent reflection, every
property can by itself be posited as essential, while the
rest are sublated as unessential, whereby the thing's
haecceity shows itself to be subject to a plurality of distinct
laws.\footnote{Man is rational, and as such subject to the laws of duty; man is bodily, and according to this property subject to the laws of gravity.} By reflecting its
essential unity in the immediate distinction of the things,
% --- p. 187 ---
\setcounter{footnote}{0}%
\opage{187}the law comes to appearance in a \textbf{manifoldness of
distinct laws}.
In the immediate diversity of the things the essential distinction of the laws
comes to appearance; as the essential relations of the things
express the manifoldness of the laws,
so the outer, unessential relations, on the other hand, present
the indifference of haecceity toward the law
itself. Both relations are, however, appearances,
and consequently moments in one and the same reflection.
If the determinate law is to come to appearance,
it presupposes a concurrence of given conditions;
it reflects its essentiality in indifferent
relations. From this follows the \textbf{thetic} and
\textbf{hypothetical} character of the law.\footnote{In experimental physics the point is precisely to hit upon the right conditions, so that the law can come to appearance in its pure universality.} α) The law
is abstractly thetic; as, being the things' own
essence, it lacks all inner reflection in itself,
it at the same time lacks all mediation, and gives no
ground for its validity. Nevertheless the law
is an essentiality; it is the immanent
reflection of unchangeability in the changeability of the things, and
in consequence of this cannot be broken off. If the law is once
given as presupposition, then at the same time
all its consequences are given. The development of the laws
from the law thus has an \textbf{analytic}
character. β) But as the apparent essence
the law has no validity outside its appearance,
and outside the existence of the things no fulfillment whatever.
As the immediate determinateness of the things
is essentially determined and limited by
the law, so conversely the determinateness of the law is in turn
% --- p. 188 ---
\setcounter{footnote}{0}%
\opage{188}determined by the immediacy of the things, and its fulfillment
infinitely conditioned. The opposed sides in
the law-determined relationship must be factually given; they
are not posited by the law itself; but they are presupposed in order
that the law can be posited. Therefore the manifold
laws come to appearance sporadically in the existential manifoldness of the things; the one law cannot immediately
be deduced from the other; but the leading back of the laws
to the law must take place by the \textbf{synthetic} way.\footnote{Mathematical cognition is at once analytic and synthetic.} γ)

Nevertheless the thetic and hypothetical determinateness of the law
is in what is essential one and the same reflection.
That the law comes to appearance in thetic form stems from
its lack of self-deepening; its reflection
in itself, by which it levels every immediate
distinction, is immediately its reflection in the existing
things, by which it itself passes over into plurality
and distinction. With the one thetic law, then, there is
given a manifoldness of thetic law-determinations.
All theses here immediately coincide,
for they are at one stroke. The one law therefore cannot
fetch its meaning from the other; the one
has no validity to lend the other, but each
has validity in itself. The separation between
the ground-law and the derived laws therefore does not lie
in the essence of the laws at all; for every determinate law
is the ground-law itself in particular determinateness. Insofar as
the one law has equal validity with the other,
they show themselves mutually indifferent; the one does not seem
to presuppose the other. But since the manifoldness
of the laws is an appearance of the
same thesis, their mutual
% --- p. 189 ---
\setcounter{footnote}{0}%
\opage{189}indifference itself only apparent; they presuppose one another
mutually: if one law could fall away, then
they could all vanish; if the one law persists,
then they all subsist; by this mutual presupposing
the thetic determinateness turns over into the
hypothetical, and the latter into the former. Just because the thesis
is thus at once both given in the hypothesis
and also posited as its extreme, the mutual coherence of the laws is of itself so evident, and yet a riddle.\footnote{In the Pythagorean theorem the thesis is immediately given with the hypothesis; nevertheless a Pythagoras was needed to deduce the former from the latter.}

The laws are the bond of the things; what is enigmatic in
the mutual coherence of the laws is at the same time the secret
in the essential connection of the things. That the thing,
which in its immediate independence only seems
to rest on itself, must nevertheless stand and fall
with other existences far removed from its outer contact and immediate
surroundings, expresses the riddle of connection in the most surprising form.\footnote{Therefore a Newton's genius was required to see the law of the heavenly bodies in the fall of an apple.}
The riddle lies in the fixed separation between the essence of the laws
and the existence of the things. The medium
through which the distinction of the laws reflects its identity
is as yet presupposed as an abstract otherness which
is indeed subject to the law, but nevertheless falls outside the law.
It is the immediate haecceity of the things that
steps in between law and law, and thereby conceals the connection.
For the same reason the law too now forms,
from its side, an abstract opposition to the things;
for its unity and distinction keep themselves exclusively in
reflection, and develop into a manifoldness of
% --- p. 190 ---
\setcounter{footnote}{0}%
\opage{190}\textbf{essential} relations. The existence of the things, on the other hand, is
reduced to immediacy; it is scattered and gathered
in \textbf{unessential} and indifferent relations. But
however essential this separation may seem from one side,
just as essential is its sublation
in the many-sided development of totality, whereby the riddle
is solved.\footnote{It goes without saying that we are not thinking here of the empirical but of the logical solution.}.
Totality, formal in its pure appearance,
has through the opposition of the laws and the things gradually
gained a concrete content. Each of its
principal sides makes its appearance in a peculiar form, and
since the form on both sides \textbf{is} essential and immanent,
its content too has an essential determinateness.
As the realm of laws is the calm prototype of the phenomenal world,
beheld in the mirror of reflection,
so the phenomenal world is conversely the restless copy of the laws,
regarded in the light of immediacy.
It remains to show how these
two sides relate to one another.
If they are, each by itself, only sides, where then is the medium
in which they have their independence? If
they are each by itself totalities, if the one totality has
in its development separated itself into two equally independent
totalities, how will it then be possible
to lead so absolute a dualism back to the original
ground-unity? a) Totality is in the first place determined as phenomenal-existence, and the law is its own essential
side. As the foundation of the things, the essence of the law

coincides with existence's own immanent
subsistence. The content-determinations that vanish
% --- p. 191 ---
\setcounter{footnote}{0}%
\opage{191}as sides in the law-determined relationship \textbf{are} reflexes
of the things' own immediate determinateness, essence-moments
in haecceity itself. Since now the things
are thus not merely sublated as sides in the law,
but also posited as independent phenomena in
haecceity, the phenomenal world has not
only the content of the law in itself, but besides this
a manifoldness of unessential determinations indifferent to the law,
hence the whole content doubled
in itself. Thus phenomenal-existence becomes
what overreaches; it subsists in the totality of the laws,
as in its otherness; this subsistence is
nevertheless its own reflection-in-itself; it has otherness itself
in itself, and subsists in essential reflection.
But as phenomenal-existence, through its own
otherness, returns into itself as reflected
world, it becomes from the opposite side eo
ipso alien to itself; for the immediacy
with which it was identical from the first has now become
its otherness; in order to subsist in this otherness,
it sublates its reflected subsistence again,
and comes to appearance in the outer manifoldness of the things.
b) Totality must next be determined as reflected
totality, so that phenomenal-existence becomes a
side of the law's own overreaching content. Since
all things subsist by virtue of the law, which posits their
immanent limit, the subsistence
by which they fall outside one another is reflected as
a moment in the subsistence according to which they fall
together with one another, and are identified with the essence of the law.
From the standpoint of the phenomenal world
the former is the positive, the latter the negative
subsistence; from the standpoint of the law, on the other hand, the relation is
the reverse. Now as this distinction between
% --- p. 192 ---
\setcounter{footnote}{0}%
\opage{192}the positive and negative vanishes again in the law's
own dialectic, the content of the law itself becomes posited as
the whole totality. The immediate existence of the things
is their positive determinate being in the form of finitude.
But since the law is hereby posited as the immanent
limit, the thing's positive determinateness is at the same time
an expression of the law's own negative determination;
by coming to appearance in the thing as in an
immediate other, it becomes to itself \textbf{an} other, i.e.
it separates itself into a plurality of determinate laws.
The law hereby takes up otherness as an integrating
moment in its own essence. The one thing is in
its subsistence dependent on the other, for the law is
their foundation; it is at the same time distinct from the
other, for the law is their limit. Now as the
one thing subsists in itself by subsisting in the
other's subsistence, it subsists in the giving up of
its independence and transition into another existence;
but since every other existence to which
the transition takes place again subsists in its otherness,
the immediate independence of the things returns
uninterruptedly through its negation; it therefore has
its immediacy in reflection and its reflection
in immediacy, so that this immanent
transition and return is the self-movement of the law.
Besides its own essential content, the realm of laws
then also comprises the changing manifoldness
of phenomenal-existence, and only gains the rest of essentiality
in the unrest of movement.
c) Insofar as the \textbf{immediate totality}, by virtue
of its own reflection, fixes itself in the form of immediacy,
a for-itself realm of
laws is an unessential abstraction: \textbf{only the} phenomenal \textbf{world exists}. Insofar as the

\medskip
\noindent[The text breaks off here, at the end of p.~192 and of the Third Fascicle; Fascicle 4 is not in the scanned volume.]

\clearpage
% PDF 202: the title of Fascicle 3, bound at the end of KB's copy
\begin{center}
\vspace*{4em}
{\Huge\textbf{The Speculative Logic}}\\[1em]
{\Large in Its Main Outlines}\\[1em]
by\\[1em]
{\Large\textbf{Lic. Rasmus Nielsen,}}\\[0.5em]
Professor Extraordinary of Philosophy at the University of Copenhagen.\\[2em]
{\large\textbf{Third Fascicle.}}\\[1em]
\textit{1842.}
\end{center}
\end{document}
